\originalsection{第9次作业}
\sourceinfo{Jeremy Orloff, MIT 18.04, Spring 2018, Problem Set 9, 题面 PDF 第1--2页, 官方解答 PDF 第1--8页; MIT OpenCourseWare, CC BY-NC-SA 4.0. 本节为中文翻译、推导补全与原生 TikZ 重绘}

\begin{exercise}
\textbf{原题 PS9.1。}\label{ass:ps9-1}
(a) 从均匀流出发, 在点 $2\ii$ 加一个涡. 求复势, 由此计算流函数, 并画若干流线. 可手绘, 也可用 Matlab 或 Mathematica 等软件.
(b) 从 (a) 的同一流场出发, 用 Milne--Thomson 定理使流体绕圆柱 $|z|=1$ 流动. 求复势、流函数并画若干流线.
(c) 解释为何 (a)(b) 的流场在远离原点处都近似均匀流.
\sourceinfo{题面第1页; 官方解答第1--3页}
\end{exercise}
\begin{solution}
据官方解答翻译并补全, 并修正涡位置的排印错误.
(a) 均匀流的复势为 $Uz$, 速度为 $(\operatorname{Re}U,-\operatorname{Im}U)$. 位于 $2\ii$ 的涡可取复势 $\ii Q\log(z-2\ii)$, $Q\in\R$ 表示强度. 为与官方图比较, 以下取 $U=1$, 图中再取 $Q=1$.
\[
 \Phi_a(z)=z+\ii Q\log(z-2\ii),\qquad
 \psi_a(x,y)=y+Q\log|z-2\ii|
 =y+\frac Q2\log\bigl(x^2+(y-2)^2\bigr).
\]
对数需在避开涡的适当割线域选支; 势的实部可以多值, 但上述流函数和速度是单值的. 求导得到
\[
 \Phi_a'(z)=1+\frac{\ii Q}{z-2\ii},\qquad
 (u,v)=\left(1+\frac{Q(y-2)}{x^2+(y-2)^2},
 -\frac{Qx}{x^2+(y-2)^2}\right).
\]
流线是 $\psi_a=c$. 当 $Q>0$, 在 $1+2\ii$ 处速度向下, 因而涡为顺时针. 停滞点满足 $\Phi_a'=0$, 得 $z=(2-Q)\ii$; $Q=1$ 时为 $\ii$.

编者校订: 官方第1页复势框内写 $z+\ii Q\log z$, 图注也说涡在 $0$, 但原题、图中涡位置及其后流函数均为 $2\ii$. 此处保留原题的 $2\ii$, 统一改为 $\log(z-2\ii)$; 不把涡移至原点.

(b) 圆定理对初始复势 $F$ 给出 $F(z)+\overline{F(1/\overline z)}$. 若保留一般复数 $U$, 镜像均匀流项为 $\overline U/z$. 取 $U=1$ 后得到
\[
 \Phi_b(z)=z+\frac1z+\ii Q\log(z-2\ii)
 -\ii Q\log\left(\frac1z+2\ii\right),\qquad |z|>1,
\]
外部涡 $2\ii$ 仍须从流动域中删去. 其流函数为
\[
 \psi_b(x,y)=y-\frac{y}{x^2+y^2}
 +Q\log|z-2\ii|-Q\log\left|\frac1z+2\ii\right|.
\]
在 $|z|=1$ 上有 $1/z=\overline z$, 从而 $|1/z+2\ii|=|z-2\ii|$, 并且 $y-y/|z|^2=0$. 所以圆柱边界上 $\psi_b=0$, 是流线, 流体不穿越壁面. 用于画方向箭头的导数是
\[
 \Phi_b'(z)=1-\frac1{z^2}+\frac{\ii Q}{z-2\ii}
 +\frac{\ii Q}{z(1+2\ii z)}.
\]
下图取 $U=Q=1$; 曲线为上述流函数等值线的数值采样, 以 TikZ 折线绘制. 黑箭头表示归一化速度方向; 圆环标记涡位置, 实心点表示停滞点. 右图仅画圆柱外的流动.
% Native TikZ contours of the exact stream functions, U=Q=1.
% Derived from PS9.1 official equations, with log(z-2i) correction.
% x grid: [-5,5], 601 samples; y grid: [-3,5.5], 511 samples.
% psi_a=y+log(abs(z-2i)).
% psi_b=y-y/abs(z)^2+log(abs(z-2i))-log(abs(1/z+2i)).
% Contour levels: -3,-2.5,-2,-1.5,-1,-.5,0,.3,.6,.9,1,
%                 1.3,1.6,1.9,2.2,2.5,3,3.5,4,4.5,5.
% The vortex disk abs(z-2i)<.065 is masked; for b abs(z)<=1.005
% is masked, and the exact unit-circle boundary is added separately.
% Interpolated contour vertices retained at stride 3 plus endpoints;
% coordinates rounded to 5 decimals. The paths approximate level sets.
% Velocity arrows use exact complex derivatives, normalized length .48.
% All graphical primitives below are native, editable TikZ paths.
\begin{center}
\begin{tikzpicture}[x=.49cm,y=.49cm,>=Stealth]
\draw[->](-5.2,0)--(5.5,0) node[right]{$x$};\draw[->](0,-3.2)--(0,5.9) node[above]{$y$};
\begin{scope}\clip (-5,-3) rectangle (5,5.5);
\draw[main,line width=.35pt] plot coordinates {(-5.00000,-2.95112) (-4.96667,-2.94738) (-4.91667,-2.94176) (-4.86667,-2.93612) (-4.83333,-2.93235) (-4.78333,-2.92669) (-4.73333,-2.92102) (-4.69497,-2.91667) (-4.65000,-2.91155) (-4.60000,-2.90585) (-4.55000,-2.90014) (-4.51667,-2.89633) (-4.46667,-2.89061) (-4.41667,-2.88487) (-4.38333,-2.88104) (-4.33333,-2.87530) (-4.28333,-2.86954) (-4.25000,-2.86570) (-4.20000,-2.85993) (-4.15000,-2.85416) (-4.11398,-2.85000) (-4.06667,-2.84453) (-4.01667,-2.83875) (-3.96991,-2.83333) (-3.93333,-2.82910) (-3.88333,-2.82331) (-3.83333,-2.81751) (-3.80000,-2.81365) (-3.75000,-2.80786) (-3.70000,-2.80207) (-3.66667,-2.79821) (-3.61667,-2.79242) (-3.56667,-2.78664) (-3.53333,-2.78278) (-3.48333,-2.77701) (-3.43333,-2.77124) (-3.39364,-2.76667) (-3.35000,-2.76164) (-3.30000,-2.75590) (-3.25000,-2.75017) (-3.21667,-2.74635) (-3.16667,-2.74064) (-3.11667,-2.73495) (-3.08333,-2.73116) (-3.03333,-2.72549) (-2.98333,-2.71984) (-2.95000,-2.71609) (-2.90000,-2.71048) (-2.85000,-2.70489) (-2.80608,-2.70000) (-2.76667,-2.69563) (-2.71667,-2.69011) (-2.66667,-2.68463) (-2.63333,-2.68099) (-2.58333,-2.67555) (-2.53333,-2.67016) (-2.50000,-2.66658) (-2.45000,-2.66125) (-2.40000,-2.65596) (-2.35000,-2.65072) (-2.31667,-2.64725) (-2.26667,-2.64208) (-2.21667,-2.63697) (-2.18082,-2.63333) (-2.13333,-2.62856) (-2.08333,-2.62359) (-2.03333,-2.61869) (-2.00000,-2.61545) (-1.95000,-2.61065) (-1.90000,-2.60591) (-1.85000,-2.60125) (-1.81667,-2.59818) (-1.76667,-2.59364) (-1.71667,-2.58918) (-1.66667,-2.58480) (-1.63333,-2.58193) (-1.58333,-2.57770) (-1.53333,-2.57355) (-1.48333,-2.56949) (-1.44772,-2.56667) (-1.40000,-2.56295) (-1.35000,-2.55916) (-1.30000,-2.55547) (-1.25000,-2.55189) (-1.21667,-2.54957) (-1.16667,-2.54617) (-1.11667,-2.54290) (-1.06667,-2.53974) (-1.01667,-2.53670) (-0.96667,-2.53379) (-0.93333,-2.53191) (-0.88333,-2.52921) (-0.83333,-2.52664) (-0.78333,-2.52421) (-0.73333,-2.52191) (-0.68333,-2.51976) (-0.63333,-2.51774) (-0.60000,-2.51648) (-0.55000,-2.51471) (-0.50000,-2.51308) (-0.45000,-2.51160) (-0.40000,-2.51027) (-0.35000,-2.50910) (-0.30000,-2.50808) (-0.25000,-2.50721) (-0.20000,-2.50650) (-0.15000,-2.50595) (-0.10000,-2.50556) (-0.05000,-2.50532) (0.00000,-2.50524) (0.05000,-2.50532) (0.10000,-2.50556) (0.15000,-2.50595) (0.20000,-2.50650) (0.25000,-2.50721) (0.30000,-2.50808) (0.35000,-2.50910) (0.40000,-2.51027) (0.45000,-2.51160) (0.50000,-2.51308) (0.55000,-2.51471) (0.60000,-2.51648) (0.63333,-2.51774) (0.68333,-2.51976) (0.73333,-2.52191) (0.78333,-2.52421) (0.83333,-2.52664) (0.88333,-2.52921) (0.93333,-2.53191) (0.96667,-2.53379) (1.01667,-2.53670) (1.06667,-2.53974) (1.11667,-2.54290) (1.16667,-2.54617) (1.21667,-2.54957) (1.25000,-2.55189) (1.30000,-2.55547) (1.35000,-2.55916) (1.40000,-2.56295) (1.44772,-2.56667) (1.48333,-2.56949) (1.53333,-2.57355) (1.58333,-2.57770) (1.63333,-2.58193) (1.66667,-2.58480) (1.71667,-2.58918) (1.76667,-2.59364) (1.81667,-2.59818) (1.85000,-2.60125) (1.90000,-2.60591) (1.95000,-2.61065) (2.00000,-2.61545) (2.03333,-2.61869) (2.08333,-2.62359) (2.13333,-2.62856) (2.18082,-2.63333) (2.21667,-2.63697) (2.26667,-2.64208) (2.31667,-2.64725) (2.35000,-2.65072) (2.40000,-2.65596) (2.45000,-2.66125) (2.50000,-2.66658) (2.53333,-2.67016) (2.58333,-2.67555) (2.63333,-2.68099) (2.66667,-2.68463) (2.71667,-2.69011) (2.76667,-2.69563) (2.80608,-2.70000) (2.85000,-2.70489) (2.90000,-2.71048) (2.95000,-2.71609) (2.98333,-2.71984) (3.03333,-2.72549) (3.08333,-2.73116) (3.11667,-2.73495) (3.16667,-2.74064) (3.21667,-2.74635) (3.25000,-2.75017) (3.30000,-2.75590) (3.35000,-2.76164) (3.39364,-2.76667) (3.43333,-2.77124) (3.48333,-2.77701) (3.53333,-2.78278) (3.56667,-2.78664) (3.61667,-2.79242) (3.66667,-2.79821) (3.70000,-2.80207) (3.75000,-2.80786) (3.80000,-2.81365) (3.83333,-2.81751) (3.88333,-2.82331) (3.93333,-2.82910) (3.96991,-2.83333) (4.01667,-2.83875) (4.06667,-2.84453) (4.11398,-2.85000) (4.15000,-2.85416) (4.20000,-2.85993) (4.25000,-2.86570) (4.28333,-2.86954) (4.33333,-2.87530) (4.38333,-2.88104) (4.41667,-2.88487) (4.46667,-2.89061) (4.51667,-2.89633) (4.55000,-2.90014) (4.60000,-2.90585) (4.65000,-2.91155) (4.69497,-2.91667) (4.73333,-2.92102) (4.78333,-2.92669) (4.83333,-2.93235) (4.86667,-2.93612) (4.91667,-2.94176) (4.96667,-2.94738) (5.00000,-2.95112)};
\draw[main,line width=.35pt] plot coordinates {(-5.00000,-2.39575) (-4.95000,-2.38947) (-4.90123,-2.38333) (-4.86667,-2.37897) (-4.81667,-2.37263) (-4.76975,-2.36667) (-4.73333,-2.36202) (-4.68333,-2.35563) (-4.63948,-2.35000) (-4.60000,-2.34492) (-4.55000,-2.33846) (-4.51038,-2.33333) (-4.46667,-2.32766) (-4.41667,-2.32115) (-4.38237,-2.31667) (-4.33333,-2.31025) (-4.28333,-2.30368) (-4.25000,-2.29929) (-4.20000,-2.29269) (-4.15000,-2.28607) (-4.11667,-2.28165) (-4.06667,-2.27500) (-4.01667,-2.26833) (-3.98333,-2.26388) (-3.93333,-2.25718) (-3.88333,-2.25047) (-3.85000,-2.24599) (-3.80000,-2.23925) (-3.75615,-2.23333) (-3.71667,-2.22799) (-3.66667,-2.22122) (-3.63311,-2.21667) (-3.58333,-2.20990) (-3.53333,-2.20310) (-3.50000,-2.19856) (-3.45000,-2.19174) (-3.40000,-2.18491) (-3.36667,-2.18036) (-3.31667,-2.17352) (-3.26667,-2.16668) (-3.23333,-2.16212) (-3.18333,-2.15527) (-3.14486,-2.15000) (-3.10000,-2.14386) (-3.05000,-2.13701) (-3.01667,-2.13245) (-2.96667,-2.12561) (-2.91667,-2.11878) (-2.88333,-2.11422) (-2.83333,-2.10740) (-2.78333,-2.10060) (-2.75000,-2.09607) (-2.70000,-2.08928) (-2.65605,-2.08333) (-2.61667,-2.07801) (-2.56667,-2.07128) (-2.53230,-2.06667) (-2.48333,-2.06011) (-2.43333,-2.05345) (-2.40000,-2.04902) (-2.35000,-2.04241) (-2.30000,-2.03584) (-2.26667,-2.03147) (-2.21667,-2.02497) (-2.16667,-2.01850) (-2.13333,-2.01422) (-2.08333,-2.00784) (-2.03333,-2.00152) (-2.00000,-1.99733) (-1.95000,-1.99111) (-1.90000,-1.98494) (-1.86667,-1.98088) (-1.81667,-1.97483) (-1.76667,-1.96887) (-1.73333,-1.96493) (-1.68333,-1.95910) (-1.63333,-1.95336) (-1.60000,-1.94959) (-1.55000,-1.94401) (-1.50000,-1.93853) (-1.45162,-1.93333) (-1.41667,-1.92964) (-1.36667,-1.92446) (-1.31667,-1.91940) (-1.28333,-1.91610) (-1.23333,-1.91125) (-1.18333,-1.90654) (-1.13333,-1.90198) (-1.10000,-1.89902) (-1.05000,-1.89470) (-1.00000,-1.89053) (-0.95000,-1.88653) (-0.90840,-1.88333) (-0.86667,-1.88024) (-0.81667,-1.87669) (-0.76667,-1.87333) (-0.71667,-1.87015) (-0.66667,-1.86716) (-0.63333,-1.86528) (-0.58333,-1.86262) (-0.53333,-1.86016) (-0.48333,-1.85791) (-0.43333,-1.85587) (-0.38333,-1.85404) (-0.33333,-1.85243) (-0.28333,-1.85104) (-0.23910,-1.85000) (-0.20000,-1.84922) (-0.15000,-1.84842) (-0.10000,-1.84785) (-0.05000,-1.84751) (0.00000,-1.84740) (0.05000,-1.84751) (0.10000,-1.84785) (0.15000,-1.84842) (0.20000,-1.84922) (0.23910,-1.85000) (0.28333,-1.85104) (0.33333,-1.85243) (0.38333,-1.85404) (0.43333,-1.85587) (0.48333,-1.85791) (0.53333,-1.86016) (0.58333,-1.86262) (0.63333,-1.86528) (0.66667,-1.86716) (0.71667,-1.87015) (0.76667,-1.87333) (0.81667,-1.87669) (0.86667,-1.88024) (0.90840,-1.88333) (0.95000,-1.88653) (1.00000,-1.89053) (1.05000,-1.89470) (1.10000,-1.89902) (1.13333,-1.90198) (1.18333,-1.90654) (1.23333,-1.91125) (1.28333,-1.91610) (1.31667,-1.91940) (1.36667,-1.92446) (1.41667,-1.92964) (1.45162,-1.93333) (1.50000,-1.93853) (1.55000,-1.94401) (1.60000,-1.94959) (1.63333,-1.95336) (1.68333,-1.95910) (1.73333,-1.96493) (1.76667,-1.96887) (1.81667,-1.97483) (1.86667,-1.98088) (1.90000,-1.98494) (1.95000,-1.99111) (2.00000,-1.99733) (2.03333,-2.00152) (2.08333,-2.00784) (2.13333,-2.01422) (2.16667,-2.01850) (2.21667,-2.02497) (2.26667,-2.03147) (2.30000,-2.03584) (2.35000,-2.04241) (2.40000,-2.04902) (2.43333,-2.05345) (2.48333,-2.06011) (2.53230,-2.06667) (2.56667,-2.07128) (2.61667,-2.07801) (2.65605,-2.08333) (2.70000,-2.08928) (2.75000,-2.09607) (2.78333,-2.10060) (2.83333,-2.10740) (2.88333,-2.11422) (2.91667,-2.11878) (2.96667,-2.12561) (3.01667,-2.13245) (3.05000,-2.13701) (3.10000,-2.14386) (3.14486,-2.15000) (3.18333,-2.15527) (3.23333,-2.16212) (3.26667,-2.16668) (3.31667,-2.17352) (3.36667,-2.18036) (3.40000,-2.18491) (3.45000,-2.19174) (3.50000,-2.19856) (3.53333,-2.20310) (3.58333,-2.20990) (3.63311,-2.21667) (3.66667,-2.22122) (3.71667,-2.22799) (3.75615,-2.23333) (3.80000,-2.23925) (3.85000,-2.24599) (3.88333,-2.25047) (3.93333,-2.25718) (3.98333,-2.26388) (4.01667,-2.26833) (4.06667,-2.27500) (4.11667,-2.28165) (4.15000,-2.28607) (4.20000,-2.29269) (4.25000,-2.29929) (4.28333,-2.30368) (4.33333,-2.31025) (4.38237,-2.31667) (4.41667,-2.32115) (4.46667,-2.32766) (4.51038,-2.33333) (4.55000,-2.33846) (4.60000,-2.34492) (4.63948,-2.35000) (4.68333,-2.35563) (4.73333,-2.36202) (4.76975,-2.36667) (4.81667,-2.37263) (4.86667,-2.37897) (4.90123,-2.38333) (4.95000,-2.38947) (5.00000,-2.39575)};
\draw[main,line width=.35pt] plot coordinates {(-0.01667,1.91159) (-0.05000,1.92706) (-0.07045,1.95000) (-0.08221,2.00000) (-0.06667,2.04154) (-0.03779,2.06667) (0.00000,2.07649) (0.03779,2.06667) (0.06667,2.04154) (0.08221,2.00000) (0.07045,1.95000) (0.05000,1.92706) (0.01667,1.91159) (-0.01667,1.91159)};
\draw[main,line width=.35pt] plot coordinates {(-5.00000,-1.84138) (-4.95000,-1.83440) (-4.91667,-1.82973) (-4.86667,-1.82269) (-4.82409,-1.81667) (-4.78333,-1.81088) (-4.73333,-1.80374) (-4.70000,-1.79896) (-4.65000,-1.79176) (-4.60000,-1.78452) (-4.56667,-1.77968) (-4.51667,-1.77238) (-4.47774,-1.76667) (-4.43333,-1.76013) (-4.38333,-1.75273) (-4.35000,-1.74777) (-4.30000,-1.74031) (-4.25347,-1.73333) (-4.21667,-1.72779) (-4.16667,-1.72023) (-4.13333,-1.71517) (-4.08333,-1.70754) (-4.03412,-1.70000) (-4.00000,-1.69475) (-3.95000,-1.68703) (-3.91667,-1.68186) (-3.86667,-1.67407) (-3.81932,-1.66667) (-3.78333,-1.66102) (-3.73333,-1.65314) (-3.70000,-1.64786) (-3.65000,-1.63992) (-3.60869,-1.63333) (-3.56667,-1.62661) (-3.51667,-1.61857) (-3.48333,-1.61320) (-3.43333,-1.60511) (-3.40000,-1.59970) (-3.35000,-1.59156) (-3.30000,-1.58338) (-3.26667,-1.57792) (-3.21667,-1.56969) (-3.18333,-1.56419) (-3.13333,-1.55592) (-3.09765,-1.55000) (-3.05000,-1.54207) (-3.00000,-1.53373) (-2.96667,-1.52815) (-2.91667,-1.51977) (-2.88333,-1.51417) (-2.83333,-1.50576) (-2.79920,-1.50000) (-2.75000,-1.49169) (-2.70064,-1.48333) (-2.66667,-1.47758) (-2.61667,-1.46909) (-2.58333,-1.46343) (-2.53333,-1.45493) (-2.50000,-1.44926) (-2.45000,-1.44076) (-2.40635,-1.43333) (-2.36667,-1.42658) (-2.31667,-1.41808) (-2.28333,-1.41242) (-2.23333,-1.40393) (-2.20000,-1.39828) (-2.15000,-1.38983) (-2.11153,-1.38333) (-2.06667,-1.37578) (-2.01667,-1.36739) (-1.98333,-1.36181) (-1.93333,-1.35348) (-1.90000,-1.34795) (-1.85000,-1.33969) (-1.81129,-1.33333) (-1.76667,-1.32604) (-1.71667,-1.31794) (-1.68333,-1.31258) (-1.63333,-1.30460) (-1.60000,-1.29932) (-1.55000,-1.29148) (-1.50000,-1.28374) (-1.46667,-1.27863) (-1.41667,-1.27107) (-1.38333,-1.26609) (-1.33333,-1.25873) (-1.28333,-1.25150) (-1.25000,-1.24677) (-1.20000,-1.23979) (-1.15266,-1.23333) (-1.11667,-1.22852) (-1.06667,-1.22200) (-1.02460,-1.21667) (-0.98333,-1.21156) (-0.93333,-1.20557) (-0.88496,-1.20000) (-0.85000,-1.19610) (-0.80000,-1.19073) (-0.75000,-1.18562) (-0.71667,-1.18236) (-0.66667,-1.17769) (-0.61667,-1.17330) (-0.56667,-1.16922) (-0.53333,-1.16666) (-0.48333,-1.16309) (-0.43333,-1.15983) (-0.38333,-1.15691) (-0.33333,-1.15432) (-0.28333,-1.15209) (-0.23333,-1.15020) (-0.20000,-1.14914) (-0.15000,-1.14785) (-0.10000,-1.14693) (-0.05000,-1.14637) (0.00000,-1.14619) (0.05000,-1.14637) (0.10000,-1.14693) (0.15000,-1.14785) (0.20000,-1.14914) (0.23333,-1.15020) (0.28333,-1.15209) (0.33333,-1.15432) (0.38333,-1.15691) (0.43333,-1.15983) (0.48333,-1.16309) (0.53333,-1.16666) (0.56667,-1.16922) (0.61667,-1.17330) (0.66667,-1.17769) (0.71667,-1.18236) (0.75000,-1.18562) (0.80000,-1.19073) (0.85000,-1.19610) (0.88496,-1.20000) (0.93333,-1.20557) (0.98333,-1.21156) (1.02460,-1.21667) (1.06667,-1.22200) (1.11667,-1.22852) (1.15266,-1.23333) (1.20000,-1.23979) (1.25000,-1.24677) (1.28333,-1.25150) (1.33333,-1.25873) (1.38333,-1.26609) (1.41667,-1.27107) (1.46667,-1.27863) (1.50000,-1.28374) (1.55000,-1.29148) (1.60000,-1.29932) (1.63333,-1.30460) (1.68333,-1.31258) (1.71667,-1.31794) (1.76667,-1.32604) (1.81129,-1.33333) (1.85000,-1.33969) (1.90000,-1.34795) (1.93333,-1.35348) (1.98333,-1.36181) (2.01667,-1.36739) (2.06667,-1.37578) (2.11153,-1.38333) (2.15000,-1.38983) (2.20000,-1.39828) (2.23333,-1.40393) (2.28333,-1.41242) (2.31667,-1.41808) (2.36667,-1.42658) (2.40635,-1.43333) (2.45000,-1.44076) (2.50000,-1.44926) (2.53333,-1.45493) (2.58333,-1.46343) (2.61667,-1.46909) (2.66667,-1.47758) (2.70064,-1.48333) (2.75000,-1.49169) (2.79920,-1.50000) (2.83333,-1.50576) (2.88333,-1.51417) (2.91667,-1.51977) (2.96667,-1.52815) (3.00000,-1.53373) (3.05000,-1.54207) (3.09765,-1.55000) (3.13333,-1.55592) (3.18333,-1.56419) (3.21667,-1.56969) (3.26667,-1.57792) (3.30000,-1.58338) (3.35000,-1.59156) (3.40000,-1.59970) (3.43333,-1.60511) (3.48333,-1.61320) (3.51667,-1.61857) (3.56667,-1.62661) (3.60869,-1.63333) (3.65000,-1.63992) (3.70000,-1.64786) (3.73333,-1.65314) (3.78333,-1.66102) (3.81932,-1.66667) (3.86667,-1.67407) (3.91667,-1.68186) (3.95000,-1.68703) (4.00000,-1.69475) (4.03412,-1.70000) (4.08333,-1.70754) (4.13333,-1.71517) (4.16667,-1.72023) (4.21667,-1.72779) (4.25347,-1.73333) (4.30000,-1.74031) (4.35000,-1.74777) (4.38333,-1.75273) (4.43333,-1.76013) (4.47774,-1.76667) (4.51667,-1.77238) (4.56667,-1.77968) (4.60000,-1.78452) (4.65000,-1.79176) (4.70000,-1.79896) (4.73333,-1.80374) (4.78333,-1.81088) (4.82409,-1.81667) (4.86667,-1.82269) (4.91667,-1.82973) (4.95000,-1.83440) (5.00000,-1.84138)};
\draw[main,line width=.35pt] plot coordinates {(-0.03333,1.84544) (-0.06667,1.85900) (-0.09752,1.88333) (-0.11667,1.90854) (-0.13321,1.95000) (-0.13675,1.98333) (-0.13212,2.01667) (-0.11667,2.05320) (-0.09255,2.08333) (-0.06667,2.10243) (-0.02906,2.11667) (0.01667,2.11911) (0.05000,2.11057) (0.08333,2.09123) (0.10778,2.06667) (0.12681,2.03333) (0.13544,2.00000) (0.13333,1.95078) (0.12130,1.91667) (0.10000,1.88593) (0.07813,1.86667) (0.04676,1.85000) (0.00000,1.84115) (-0.03333,1.84544)};
\draw[main,line width=.35pt] plot coordinates {(-5.00000,-1.50974) (-4.95000,-1.50232) (-4.91667,-1.49735) (-4.86667,-1.48986) (-4.82336,-1.48333) (-4.78333,-1.47727) (-4.73333,-1.46965) (-4.70000,-1.46455) (-4.65000,-1.45685) (-4.60575,-1.45000) (-4.56667,-1.44392) (-4.51667,-1.43609) (-4.48333,-1.43085) (-4.43333,-1.42294) (-4.39388,-1.41667) (-4.35000,-1.40965) (-4.30000,-1.40161) (-4.26667,-1.39623) (-4.21667,-1.38810) (-4.18333,-1.38266) (-4.13333,-1.37445) (-4.08622,-1.36667) (-4.05000,-1.36065) (-4.00000,-1.35231) (-3.96667,-1.34671) (-3.91667,-1.33828) (-3.88333,-1.33263) (-3.83333,-1.32410) (-3.78996,-1.31667) (-3.75000,-1.30978) (-3.70000,-1.30112) (-3.66667,-1.29531) (-3.61667,-1.28656) (-3.58333,-1.28069) (-3.53333,-1.27185) (-3.50000,-1.26592) (-3.45000,-1.25699) (-3.41108,-1.25000) (-3.36667,-1.24198) (-3.31900,-1.23333) (-3.28333,-1.22683) (-3.23333,-1.21766) (-3.20000,-1.21153) (-3.15000,-1.20227) (-3.11667,-1.19608) (-3.06667,-1.18674) (-3.03333,-1.18048) (-2.98333,-1.17106) (-2.95000,-1.16475) (-2.90000,-1.15524) (-2.86667,-1.14887) (-2.81667,-1.13928) (-2.78333,-1.13286) (-2.73333,-1.12319) (-2.69971,-1.11667) (-2.65000,-1.10698) (-2.61431,-1.10000) (-2.56667,-1.09065) (-2.52954,-1.08333) (-2.48333,-1.07420) (-2.44533,-1.06667) (-2.40000,-1.05765) (-2.36161,-1.05000) (-2.31667,-1.04101) (-2.27833,-1.03333) (-2.23333,-1.02430) (-2.19540,-1.01667) (-2.15000,-1.00751) (-2.11275,-1.00000) (-2.06667,-0.99069) (-2.03028,-0.98333) (-1.98333,-0.97383) (-1.94790,-0.96667) (-1.90000,-0.95698) (-1.86547,-0.95000) (-1.81667,-0.94014) (-1.78288,-0.93333) (-1.73333,-0.92336) (-1.69996,-0.91667) (-1.65000,-0.90667) (-1.61653,-0.90000) (-1.56667,-0.89010) (-1.53238,-0.88333) (-1.48333,-0.87370) (-1.44722,-0.86667) (-1.40000,-0.85752) (-1.36075,-0.85000) (-1.31667,-0.84162) (-1.27254,-0.83333) (-1.23333,-0.82605) (-1.18333,-0.81689) (-1.15000,-0.81088) (-1.10000,-0.80200) (-1.06667,-0.79618) (-1.01667,-0.78763) (-0.98333,-0.78205) (-0.93333,-0.77386) (-0.88796,-0.76667) (-0.85000,-0.76080) (-0.80000,-0.75333) (-0.76667,-0.74853) (-0.71667,-0.74158) (-0.66667,-0.73498) (-0.63333,-0.73078) (-0.58333,-0.72480) (-0.53333,-0.71921) (-0.50000,-0.71573) (-0.45000,-0.71085) (-0.40000,-0.70643) (-0.35000,-0.70249) (-0.31469,-0.70000) (-0.26667,-0.69700) (-0.21667,-0.69439) (-0.16667,-0.69231) (-0.11667,-0.69076) (-0.06667,-0.68975) (-0.01667,-0.68929) (0.03333,-0.68938) (0.08333,-0.69003) (0.13333,-0.69121) (0.18333,-0.69294) (0.23333,-0.69520) (0.28333,-0.69799) (0.31667,-0.70013) (0.36667,-0.70375) (0.41667,-0.70785) (0.46667,-0.71243) (0.50908,-0.71667) (0.55000,-0.72103) (0.60000,-0.72675) (0.65000,-0.73286) (0.68333,-0.73714) (0.73333,-0.74386) (0.77693,-0.75000) (0.81667,-0.75579) (0.86667,-0.76335) (0.90000,-0.76855) (0.95000,-0.77656) (0.99103,-0.78333) (1.03333,-0.79046) (1.08333,-0.79908) (1.11667,-0.80494) (1.16667,-0.81388) (1.20000,-0.81993) (1.25000,-0.82913) (1.28333,-0.83535) (1.33333,-0.84477) (1.36667,-0.85113) (1.41667,-0.86074) (1.45000,-0.86720) (1.50000,-0.87697) (1.53333,-0.88352) (1.58333,-0.89340) (1.61667,-0.90003) (1.66667,-0.91000) (1.70000,-0.91667) (1.75000,-0.92671) (1.78333,-0.93342) (1.83333,-0.94351) (1.86667,-0.95024) (1.91667,-0.96035) (1.95000,-0.96709) (2.00000,-0.97721) (2.03333,-0.98395) (2.08333,-0.99406) (2.11667,-1.00079) (2.16667,-1.01088) (2.20000,-1.01759) (2.25000,-1.02765) (2.28333,-1.03434) (2.33333,-1.04435) (2.36667,-1.05101) (2.41667,-1.06097) (2.45000,-1.06759) (2.50000,-1.07750) (2.53333,-1.08408) (2.58333,-1.09392) (2.61667,-1.10046) (2.66667,-1.11023) (2.70000,-1.11672) (2.75000,-1.12642) (2.78578,-1.13333) (2.83333,-1.14248) (2.87257,-1.15000) (2.91667,-1.15841) (2.96013,-1.16667) (3.00000,-1.17420) (3.04851,-1.18333) (3.08333,-1.18986) (3.13333,-1.19918) (3.16667,-1.20536) (3.21667,-1.21460) (3.25000,-1.22072) (3.30000,-1.22987) (3.33333,-1.23594) (3.38333,-1.24500) (3.41667,-1.25101) (3.46667,-1.25997) (3.50418,-1.26667) (3.55000,-1.27480) (3.59833,-1.28333) (3.63333,-1.28948) (3.68333,-1.29822) (3.71667,-1.30401) (3.76667,-1.31266) (3.80000,-1.31839) (3.85000,-1.32695) (3.88750,-1.33333) (3.93333,-1.34110) (3.98333,-1.34951) (4.01667,-1.35509) (4.06667,-1.36342) (4.10000,-1.36895) (4.15000,-1.37719) (4.18746,-1.38333) (4.23333,-1.39082) (4.28333,-1.39892) (4.31667,-1.40430) (4.36667,-1.41232) (4.40000,-1.41764) (4.45000,-1.42558) (4.49911,-1.43333) (4.53333,-1.43871) (4.58333,-1.44651) (4.61667,-1.45169) (4.66667,-1.45942) (4.71382,-1.46667) (4.75000,-1.47220) (4.80000,-1.47980) (4.83333,-1.48484) (4.88333,-1.49236) (4.93333,-1.49984) (4.96667,-1.50480) (5.00000,-1.50974)};
\draw[main,line width=.35pt] plot coordinates {(-0.06667,1.78306) (-0.09887,1.80000) (-0.12052,1.81667) (-0.15000,1.84977) (-0.16667,1.87852) (-0.18030,1.91667) (-0.18550,1.95000) (-0.18333,1.99529) (-0.17370,2.03333) (-0.15751,2.06667) (-0.13333,2.09817) (-0.11320,2.11667) (-0.08333,2.13601) (-0.04676,2.15000) (0.00000,2.15644) (0.04676,2.15000) (0.08333,2.13601) (0.11320,2.11667) (0.13333,2.09817) (0.15751,2.06667) (0.17370,2.03333) (0.18333,1.99529) (0.18550,1.95000) (0.18030,1.91667) (0.16667,1.87852) (0.15000,1.84977) (0.12052,1.81667) (0.09887,1.80000) (0.06667,1.78306) (0.01667,1.77074) (-0.03333,1.77308) (-0.06667,1.78306)};
\draw[main,line width=.35pt] plot coordinates {(-5.00000,-1.17921) (-4.95000,-1.17135) (-4.91667,-1.16608) (-4.86667,-1.15813) (-4.81667,-1.15012) (-4.78333,-1.14475) (-4.73333,-1.13665) (-4.70000,-1.13122) (-4.65000,-1.12302) (-4.61152,-1.11667) (-4.56667,-1.10922) (-4.51667,-1.10086) (-4.48333,-1.09525) (-4.43333,-1.08679) (-4.40000,-1.08111) (-4.35000,-1.07254) (-4.31592,-1.06667) (-4.26667,-1.05812) (-4.22021,-1.05000) (-4.18333,-1.04352) (-4.13333,-1.03466) (-4.10000,-1.02872) (-4.05000,-1.01976) (-4.01667,-1.01374) (-3.96667,-1.00466) (-3.93333,-0.99857) (-3.88333,-0.98937) (-3.85000,-0.98320) (-3.80000,-0.97388) (-3.76155,-0.96667) (-3.71667,-0.95819) (-3.67362,-0.95000) (-3.63333,-0.94229) (-3.58689,-0.93333) (-3.55000,-0.92617) (-3.50134,-0.91667) (-3.46667,-0.90985) (-3.41695,-0.90000) (-3.38333,-0.89330) (-3.33368,-0.88333) (-3.30000,-0.87653) (-3.25152,-0.86667) (-3.21667,-0.85953) (-3.17043,-0.85000) (-3.13333,-0.84230) (-3.09040,-0.83333) (-3.05000,-0.82484) (-3.01138,-0.81667) (-2.96667,-0.80714) (-2.93333,-0.79999) (-2.88333,-0.78920) (-2.85000,-0.78196) (-2.80000,-0.77102) (-2.76667,-0.76368) (-2.71667,-0.75259) (-2.68333,-0.74515) (-2.63333,-0.73392) (-2.60000,-0.72638) (-2.55731,-0.71667) (-2.51667,-0.70736) (-2.48333,-0.69968) (-2.43333,-0.68810) (-2.40000,-0.68032) (-2.35000,-0.66859) (-2.31667,-0.66071) (-2.27156,-0.65000) (-2.23333,-0.64086) (-2.20000,-0.63286) (-2.15000,-0.62078) (-2.11667,-0.61268) (-2.06667,-0.60047) (-2.03333,-0.59228) (-1.99704,-0.58333) (-1.95000,-0.57167) (-1.91667,-0.56337) (-1.86667,-0.55086) (-1.83333,-0.54248) (-1.79705,-0.53333) (-1.75000,-0.52142) (-1.71667,-0.51295) (-1.66667,-0.50021) (-1.63333,-0.49169) (-1.60000,-0.48316) (-1.55000,-0.47033) (-1.51667,-0.46176) (-1.47093,-0.45000) (-1.43333,-0.44032) (-1.40000,-0.43174) (-1.35000,-0.41889) (-1.31667,-0.41033) (-1.27633,-0.40000) (-1.23333,-0.38901) (-1.20000,-0.38053) (-1.15000,-0.36788) (-1.11667,-0.35950) (-1.07859,-0.35000) (-1.03333,-0.33879) (-1.00000,-0.33063) (-0.95000,-0.31853) (-0.91667,-0.31058) (-0.87155,-0.30000) (-0.83333,-0.29118) (-0.79855,-0.28333) (-0.75000,-0.27262) (-0.71667,-0.26547) (-0.66667,-0.25507) (-0.63333,-0.24839) (-0.58333,-0.23876) (-0.55000,-0.23264) (-0.50000,-0.22392) (-0.45525,-0.21667) (-0.41667,-0.21080) (-0.36667,-0.20386) (-0.33333,-0.19965) (-0.28333,-0.19398) (-0.23333,-0.18916) (-0.18333,-0.18522) (-0.15000,-0.18310) (-0.10000,-0.18069) (-0.05000,-0.17923) (0.00000,-0.17874) (0.05000,-0.17923) (0.10000,-0.18069) (0.15000,-0.18310) (0.18333,-0.18522) (0.23333,-0.18916) (0.28333,-0.19398) (0.33333,-0.19965) (0.36667,-0.20386) (0.41667,-0.21080) (0.45525,-0.21667) (0.50000,-0.22392) (0.55000,-0.23264) (0.58333,-0.23876) (0.63333,-0.24839) (0.66667,-0.25507) (0.71667,-0.26547) (0.75000,-0.27262) (0.79855,-0.28333) (0.83333,-0.29118) (0.87155,-0.30000) (0.91667,-0.31058) (0.95000,-0.31853) (1.00000,-0.33063) (1.03333,-0.33879) (1.07859,-0.35000) (1.11667,-0.35950) (1.15000,-0.36788) (1.20000,-0.38053) (1.23333,-0.38901) (1.27633,-0.40000) (1.31667,-0.41033) (1.35000,-0.41889) (1.40000,-0.43174) (1.43333,-0.44032) (1.47093,-0.45000) (1.51667,-0.46176) (1.55000,-0.47033) (1.60000,-0.48316) (1.63333,-0.49169) (1.66667,-0.50021) (1.71667,-0.51295) (1.75000,-0.52142) (1.79705,-0.53333) (1.83333,-0.54248) (1.86667,-0.55086) (1.91667,-0.56337) (1.95000,-0.57167) (1.99704,-0.58333) (2.03333,-0.59228) (2.06667,-0.60047) (2.11667,-0.61268) (2.15000,-0.62078) (2.20000,-0.63286) (2.23333,-0.64086) (2.27156,-0.65000) (2.31667,-0.66071) (2.35000,-0.66859) (2.40000,-0.68032) (2.43333,-0.68810) (2.48333,-0.69968) (2.51667,-0.70736) (2.55731,-0.71667) (2.60000,-0.72638) (2.63333,-0.73392) (2.68333,-0.74515) (2.71667,-0.75259) (2.76667,-0.76368) (2.80000,-0.77102) (2.85000,-0.78196) (2.88333,-0.78920) (2.93333,-0.79999) (2.96667,-0.80714) (3.01138,-0.81667) (3.05000,-0.82484) (3.09040,-0.83333) (3.13333,-0.84230) (3.17043,-0.85000) (3.21667,-0.85953) (3.25152,-0.86667) (3.30000,-0.87653) (3.33368,-0.88333) (3.38333,-0.89330) (3.41695,-0.90000) (3.46667,-0.90985) (3.50134,-0.91667) (3.55000,-0.92617) (3.58689,-0.93333) (3.63333,-0.94229) (3.67362,-0.95000) (3.71667,-0.95819) (3.76155,-0.96667) (3.80000,-0.97388) (3.85000,-0.98320) (3.88333,-0.98937) (3.93333,-0.99857) (3.96667,-1.00466) (4.01667,-1.01374) (4.05000,-1.01976) (4.10000,-1.02872) (4.13333,-1.03466) (4.18333,-1.04352) (4.22021,-1.05000) (4.26667,-1.05812) (4.31592,-1.06667) (4.35000,-1.07254) (4.40000,-1.08111) (4.43333,-1.08679) (4.48333,-1.09525) (4.51667,-1.10086) (4.56667,-1.10922) (4.61152,-1.11667) (4.65000,-1.12302) (4.70000,-1.13122) (4.73333,-1.13665) (4.78333,-1.14475) (4.81667,-1.15012) (4.86667,-1.15813) (4.91667,-1.16608) (4.95000,-1.17135) (5.00000,-1.17921)};
\draw[main,line width=.35pt] plot coordinates {(-0.08333,1.66574) (-0.11667,1.68163) (-0.14409,1.70000) (-0.16667,1.71914) (-0.19427,1.75000) (-0.21644,1.78333) (-0.23267,1.81667) (-0.24416,1.85000) (-0.25138,1.88333) (-0.25513,1.93333) (-0.25019,1.98333) (-0.24209,2.01667) (-0.22928,2.05000) (-0.21113,2.08333) (-0.18590,2.11667) (-0.16667,2.13613) (-0.13333,2.16193) (-0.10000,2.18011) (-0.06667,2.19235) (-0.02640,2.20000) (0.01667,2.20104) (0.05000,2.19645) (0.09215,2.18333) (0.12560,2.16667) (0.15016,2.15000) (0.18333,2.11943) (0.20000,2.09926) (0.22095,2.06667) (0.23625,2.03333) (0.25000,1.98419) (0.25447,1.95000) (0.25359,1.90000) (0.24826,1.86667) (0.23333,1.81835) (0.21667,1.78373) (0.20000,1.75759) (0.18036,1.73333) (0.15000,1.70444) (0.11939,1.68333) (0.08546,1.66667) (0.05000,1.65551) (0.00000,1.65009) (-0.05000,1.65551) (-0.08333,1.66574)};
\draw[main,line width=.35pt] plot coordinates {(-5.00000,-0.85012) (-4.96667,-0.84460) (-4.91667,-0.83626) (-4.88333,-0.83067) (-4.83333,-0.82223) (-4.80000,-0.81656) (-4.75000,-0.80800) (-4.70363,-0.80000) (-4.66667,-0.79358) (-4.61667,-0.78483) (-4.58333,-0.77896) (-4.53333,-0.77010) (-4.50000,-0.76414) (-4.45000,-0.75515) (-4.41667,-0.74911) (-4.36667,-0.73999) (-4.33044,-0.73333) (-4.28333,-0.72462) (-4.24073,-0.71667) (-4.20000,-0.70901) (-4.15241,-0.70000) (-4.11667,-0.69318) (-4.06667,-0.68357) (-4.03333,-0.67711) (-3.98333,-0.66735) (-3.95000,-0.66079) (-3.90000,-0.65088) (-3.86667,-0.64423) (-3.81667,-0.63416) (-3.78333,-0.62740) (-3.73333,-0.61717) (-3.70000,-0.61030) (-3.65042,-0.60000) (-3.61667,-0.59293) (-3.57120,-0.58333) (-3.53333,-0.57528) (-3.49319,-0.56667) (-3.45000,-0.55733) (-3.41638,-0.55000) (-3.36667,-0.53908) (-3.33333,-0.53169) (-3.28333,-0.52051) (-3.25000,-0.51300) (-3.20000,-0.50162) (-3.16667,-0.49398) (-3.12067,-0.48333) (-3.08333,-0.47462) (-3.04953,-0.46667) (-3.00000,-0.45491) (-2.96667,-0.44693) (-2.91667,-0.43484) (-2.88333,-0.42671) (-2.84250,-0.41667) (-2.80000,-0.40611) (-2.76667,-0.39777) (-2.71667,-0.38513) (-2.68333,-0.37662) (-2.64468,-0.36667) (-2.60000,-0.35506) (-2.56667,-0.34631) (-2.51767,-0.33333) (-2.48333,-0.32415) (-2.45000,-0.31516) (-2.40000,-0.30153) (-2.36667,-0.29236) (-2.33333,-0.28311) (-2.28333,-0.26909) (-2.25000,-0.25965) (-2.21623,-0.25000) (-2.16667,-0.23569) (-2.13333,-0.22596) (-2.10000,-0.21615) (-2.05000,-0.20128) (-2.01667,-0.19125) (-1.98333,-0.18114) (-1.93611,-0.16667) (-1.90000,-0.15547) (-1.86667,-0.14505) (-1.82957,-0.13333) (-1.78333,-0.11857) (-1.75000,-0.10781) (-1.71667,-0.09695) (-1.67527,-0.08333) (-1.63333,-0.06938) (-1.60000,-0.05817) (-1.56667,-0.04686) (-1.52717,-0.03333) (-1.48333,-0.01814) (-1.45000,-0.00647) (-1.41667,0.00531) (-1.38333,0.01720) (-1.33854,0.03333) (-1.30000,0.04738) (-1.26667,0.05964) (-1.23333,0.07201) (-1.20000,0.08448) (-1.15888,0.10000) (-1.11667,0.11611) (-1.08333,0.12895) (-1.05000,0.14188) (-1.01667,0.15490) (-0.98333,0.16801) (-0.94464,0.18333) (-0.90287,0.20000) (-0.86667,0.21454) (-0.83333,0.22799) (-0.80000,0.24148) (-0.76667,0.25501) (-0.73333,0.26855) (-0.69691,0.28333) (-0.65580,0.30000) (-0.61667,0.31581) (-0.58333,0.32919) (-0.55000,0.34243) (-0.51667,0.35551) (-0.48333,0.36837) (-0.44357,0.38333) (-0.40000,0.39924) (-0.36667,0.41092) (-0.33333,0.42206) (-0.29748,0.43333) (-0.25000,0.44707) (-0.21667,0.45569) (-0.16689,0.46667) (-0.13333,0.47275) (-0.08333,0.47942) (-0.03333,0.48307) (0.00000,0.48378) (0.03333,0.48307) (0.08333,0.47942) (0.13333,0.47275) (0.16689,0.46667) (0.21667,0.45569) (0.25000,0.44707) (0.29748,0.43333) (0.33333,0.42206) (0.36667,0.41092) (0.40000,0.39924) (0.44357,0.38333) (0.48333,0.36837) (0.51667,0.35551) (0.55000,0.34243) (0.58333,0.32919) (0.61667,0.31581) (0.65580,0.30000) (0.69691,0.28333) (0.73333,0.26855) (0.76667,0.25501) (0.80000,0.24148) (0.83333,0.22799) (0.86667,0.21454) (0.90287,0.20000) (0.94464,0.18333) (0.98333,0.16801) (1.01667,0.15490) (1.05000,0.14188) (1.08333,0.12895) (1.11667,0.11611) (1.15888,0.10000) (1.20000,0.08448) (1.23333,0.07201) (1.26667,0.05964) (1.30000,0.04738) (1.33854,0.03333) (1.38333,0.01720) (1.41667,0.00531) (1.45000,-0.00647) (1.48333,-0.01814) (1.52717,-0.03333) (1.56667,-0.04686) (1.60000,-0.05817) (1.63333,-0.06938) (1.67527,-0.08333) (1.71667,-0.09695) (1.75000,-0.10781) (1.78333,-0.11857) (1.82957,-0.13333) (1.86667,-0.14505) (1.90000,-0.15547) (1.93611,-0.16667) (1.98333,-0.18114) (2.01667,-0.19125) (2.05000,-0.20128) (2.10000,-0.21615) (2.13333,-0.22596) (2.16667,-0.23569) (2.21623,-0.25000) (2.25000,-0.25965) (2.28333,-0.26909) (2.33333,-0.28311) (2.36667,-0.29236) (2.40000,-0.30153) (2.45000,-0.31516) (2.48333,-0.32415) (2.51767,-0.33333) (2.56667,-0.34631) (2.60000,-0.35506) (2.64468,-0.36667) (2.68333,-0.37662) (2.71667,-0.38513) (2.76667,-0.39777) (2.80000,-0.40611) (2.84250,-0.41667) (2.88333,-0.42671) (2.91667,-0.43484) (2.96667,-0.44693) (3.00000,-0.45491) (3.04953,-0.46667) (3.08333,-0.47462) (3.12067,-0.48333) (3.16667,-0.49398) (3.20000,-0.50162) (3.25000,-0.51300) (3.28333,-0.52051) (3.33333,-0.53169) (3.36667,-0.53908) (3.41638,-0.55000) (3.45000,-0.55733) (3.49319,-0.56667) (3.53333,-0.57528) (3.57120,-0.58333) (3.61667,-0.59293) (3.65042,-0.60000) (3.70000,-0.61030) (3.73333,-0.61717) (3.78333,-0.62740) (3.81667,-0.63416) (3.86667,-0.64423) (3.90000,-0.65088) (3.95000,-0.66079) (3.98333,-0.66735) (4.03333,-0.67711) (4.06667,-0.68357) (4.11667,-0.69318) (4.15241,-0.70000) (4.20000,-0.70901) (4.24073,-0.71667) (4.28333,-0.72462) (4.33044,-0.73333) (4.36667,-0.73999) (4.41667,-0.74911) (4.45000,-0.75515) (4.50000,-0.76414) (4.53333,-0.77010) (4.58333,-0.77896) (4.61667,-0.78483) (4.66667,-0.79358) (4.70363,-0.80000) (4.75000,-0.80800) (4.80000,-0.81656) (4.83333,-0.82223) (4.88333,-0.83067) (4.91667,-0.83626) (4.96667,-0.84460) (5.00000,-0.85012)};
\draw[main,line width=.35pt] plot coordinates {(0.00000,1.38318) (-0.03333,1.38547) (-0.08333,1.39793) (-0.11667,1.41233) (-0.15000,1.43208) (-0.17411,1.45000) (-0.20000,1.47246) (-0.22616,1.50000) (-0.25000,1.52924) (-0.26667,1.55300) (-0.28519,1.58333) (-0.30259,1.61667) (-0.31722,1.65000) (-0.33333,1.69606) (-0.34319,1.73333) (-0.35000,1.76903) (-0.35539,1.81667) (-0.35628,1.86667) (-0.35211,1.91667) (-0.34641,1.95000) (-0.33333,1.99850) (-0.32030,2.03333) (-0.30426,2.06667) (-0.28413,2.10000) (-0.26667,2.12385) (-0.24416,2.15000) (-0.21667,2.17598) (-0.18513,2.20000) (-0.15771,2.21667) (-0.12249,2.23333) (-0.08333,2.24651) (-0.05000,2.25353) (0.00000,2.25745) (0.05000,2.25353) (0.08333,2.24651) (0.12249,2.23333) (0.15771,2.21667) (0.18513,2.20000) (0.21667,2.17598) (0.24416,2.15000) (0.26667,2.12385) (0.28413,2.10000) (0.30426,2.06667) (0.32030,2.03333) (0.33333,1.99850) (0.34641,1.95000) (0.35211,1.91667) (0.35628,1.86667) (0.35539,1.81667) (0.35000,1.76903) (0.34319,1.73333) (0.33333,1.69606) (0.31722,1.65000) (0.30259,1.61667) (0.28519,1.58333) (0.26667,1.55300) (0.25000,1.52924) (0.22616,1.50000) (0.20000,1.47246) (0.17411,1.45000) (0.15000,1.43208) (0.11667,1.41233) (0.08333,1.39793) (0.03333,1.38547) (0.00000,1.38318)};
\draw[main,line width=.35pt] plot coordinates {(-5.00000,-0.74081) (-4.95568,-0.73333) (-4.91667,-0.72671) (-4.86667,-0.71817) (-4.83333,-0.71243) (-4.78333,-0.70377) (-4.75000,-0.69795) (-4.70000,-0.68916) (-4.66667,-0.68326) (-4.61667,-0.67435) (-4.57390,-0.66667) (-4.53333,-0.65933) (-4.48333,-0.65021) (-4.45000,-0.64408) (-4.40000,-0.63483) (-4.36667,-0.62861) (-4.31667,-0.61922) (-4.28333,-0.61291) (-4.23333,-0.60338) (-4.20000,-0.59697) (-4.15000,-0.58729) (-4.11667,-0.58078) (-4.06667,-0.57094) (-4.03333,-0.56434) (-3.98333,-0.55434) (-3.95000,-0.54763) (-3.90000,-0.53748) (-3.86667,-0.53065) (-3.81667,-0.52033) (-3.78333,-0.51339) (-3.73333,-0.50290) (-3.70000,-0.49584) (-3.65000,-0.48517) (-3.61667,-0.47800) (-3.56667,-0.46714) (-3.53333,-0.45984) (-3.48879,-0.45000) (-3.45000,-0.44136) (-3.41427,-0.43333) (-3.36667,-0.42255) (-3.33333,-0.41493) (-3.28333,-0.40339) (-3.25000,-0.39563) (-3.20000,-0.38388) (-3.16667,-0.37597) (-3.12778,-0.36667) (-3.08333,-0.35593) (-3.05000,-0.34781) (-3.00000,-0.33551) (-2.96667,-0.32722) (-2.92458,-0.31667) (-2.88333,-0.30622) (-2.85000,-0.29770) (-2.80000,-0.28479) (-2.76667,-0.27609) (-2.73087,-0.26667) (-2.68333,-0.25402) (-2.65000,-0.24506) (-2.60679,-0.23333) (-2.56667,-0.22233) (-2.53333,-0.21309) (-2.48660,-0.20000) (-2.45000,-0.18963) (-2.41667,-0.18010) (-2.37019,-0.16667) (-2.33333,-0.15589) (-2.30000,-0.14604) (-2.25745,-0.13333) (-2.21667,-0.12101) (-2.18333,-0.11082) (-2.14826,-0.10000) (-2.10000,-0.08492) (-2.06667,-0.07438) (-2.03333,-0.06372) (-1.99090,-0.05000) (-1.95000,-0.03660) (-1.91667,-0.02555) (-1.88333,-0.01438) (-1.84094,0.00000) (-1.80000,0.01408) (-1.76667,0.02569) (-1.73333,0.03743) (-1.69805,0.05000) (-1.65195,0.06667) (-1.61667,0.07960) (-1.58333,0.09198) (-1.55000,0.10450) (-1.51667,0.11718) (-1.47480,0.13333) (-1.43333,0.14959) (-1.40000,0.16284) (-1.36667,0.17627) (-1.33333,0.18988) (-1.30000,0.20367) (-1.26667,0.21765) (-1.22982,0.23333) (-1.19129,0.25000) (-1.15339,0.26667) (-1.11667,0.28308) (-1.08333,0.29822) (-1.05000,0.31359) (-1.01667,0.32920) (-0.98333,0.34506) (-0.95000,0.36117) (-0.91667,0.37755) (-0.88333,0.39420) (-0.85000,0.41115) (-0.81667,0.42838) (-0.78333,0.44593) (-0.75000,0.46380) (-0.71667,0.48200) (-0.68430,0.50000) (-0.65487,0.51667) (-0.62595,0.53333) (-0.59756,0.55000) (-0.56667,0.56849) (-0.53333,0.58884) (-0.50000,0.60964) (-0.46667,0.63089) (-0.43728,0.65000) (-0.41220,0.66667) (-0.38333,0.68632) (-0.35000,0.70949) (-0.31667,0.73320) (-0.29353,0.75000) (-0.26667,0.77015) (-0.23333,0.79566) (-0.20634,0.81667) (-0.18333,0.83536) (-0.15000,0.86298) (-0.12566,0.88333) (-0.10000,0.90627) (-0.06939,0.93333) (-0.05000,0.95183) (-0.01679,0.98333) (-0.01630,1.01667) (-0.03333,1.03395) (-0.06339,1.06667) (-0.08333,1.08810) (-0.10745,1.11667) (-0.13333,1.14735) (-0.15000,1.16833) (-0.17395,1.20000) (-0.19821,1.23333) (-0.21667,1.25989) (-0.23333,1.28540) (-0.25268,1.31667) (-0.27211,1.35000) (-0.29024,1.38333) (-0.30705,1.41667) (-0.32255,1.45000) (-0.33673,1.48333) (-0.35000,1.51779) (-0.36631,1.56667) (-0.37575,1.60000) (-0.38383,1.63333) (-0.39331,1.68333) (-0.39948,1.73333) (-0.40173,1.76667) (-0.40210,1.81667) (-0.40000,1.85235) (-0.39415,1.90000) (-0.38350,1.95000) (-0.37378,1.98333) (-0.36150,2.01667) (-0.34641,2.05000) (-0.32813,2.08333) (-0.30595,2.11667) (-0.28333,2.14485) (-0.26309,2.16667) (-0.23333,2.19351) (-0.20199,2.21667) (-0.17437,2.23333) (-0.13982,2.25000) (-0.10000,2.26422) (-0.06667,2.27225) (-0.01667,2.27816) (0.03333,2.27699) (0.08333,2.26863) (0.11667,2.25896) (0.15000,2.24562) (0.18333,2.22831) (0.21667,2.20643) (0.24535,2.18333) (0.26667,2.16297) (0.29313,2.13333) (0.31667,2.10121) (0.33333,2.07418) (0.35000,2.04230) (0.36667,2.00311) (0.37896,1.96667) (0.38761,1.93333) (0.39665,1.88333) (0.40023,1.85000) (0.40239,1.80000) (0.40079,1.75000) (0.39780,1.71667) (0.39051,1.66667) (0.38333,1.63111) (0.37120,1.58333) (0.36105,1.55000) (0.34958,1.51667) (0.33333,1.47488) (0.31667,1.43681) (0.30000,1.40217) (0.28333,1.37011) (0.26667,1.34015) (0.25000,1.31201) (0.23192,1.28333) (0.20970,1.25000) (0.18627,1.21667) (0.16667,1.18982) (0.14855,1.16667) (0.12150,1.13333) (0.10000,1.10720) (0.07840,1.08333) (0.05000,1.05152) (0.03236,1.03333) (0.00000,1.00000) (0.03333,0.96759) (0.05147,0.95000) (0.08333,0.92119) (0.10648,0.90000) (0.13333,0.87718) (0.16535,0.85000) (0.18566,0.83333) (0.21667,0.80873) (0.24920,0.78333) (0.27113,0.76667) (0.30000,0.74539) (0.33333,0.72129) (0.36343,0.70000) (0.38759,0.68333) (0.41667,0.66374) (0.45000,0.64173) (0.48333,0.62020) (0.51529,0.60000) (0.54220,0.58333) (0.56964,0.56667) (0.60000,0.54858) (0.63333,0.52909) (0.66667,0.50998) (0.70000,0.49125) (0.73333,0.47287) (0.76667,0.45483) (0.80000,0.43712) (0.83333,0.41972) (0.86667,0.40264) (0.90000,0.38584) (0.93333,0.36932) (0.96667,0.35308) (1.00000,0.33710) (1.03333,0.32137) (1.06667,0.30588) (1.10000,0.29063) (1.13333,0.27561) (1.16667,0.26080) (1.20000,0.24622) (1.23333,0.23183) (1.26900,0.21667) (1.30882,0.20000) (1.34931,0.18333) (1.38333,0.16954) (1.41667,0.15619) (1.45000,0.14303) (1.48333,0.13002) (1.51801,0.11667) (1.56193,0.10000) (1.60000,0.08577) (1.63333,0.07347) (1.66667,0.06132) (1.70000,0.04930) (1.74492,0.03333) (1.78333,0.01987) (1.81667,0.00833) (1.85000,-0.00309) (1.89012,-0.01667) (1.93333,-0.03109) (1.96667,-0.04208) (2.00000,-0.05296) (2.04251,-0.06667) (2.08333,-0.07966) (2.11667,-0.09015) (2.15000,-0.10054) (2.20000,-0.11593) (2.23333,-0.12606) (2.26667,-0.13610) (2.31337,-0.15000) (2.35000,-0.16077) (2.38333,-0.17048) (2.42793,-0.18333) (2.46667,-0.19437) (2.50000,-0.20377) (2.54621,-0.21667) (2.58333,-0.22691) (2.61667,-0.23603) (2.66667,-0.24955) (2.70000,-0.25847) (2.73333,-0.26732) (2.78333,-0.28045) (2.81667,-0.28911) (2.85898,-0.30000) (2.90000,-0.31045) (2.93333,-0.31887) (2.98333,-0.33137) (3.01667,-0.33962) (3.05896,-0.35000) (3.10000,-0.35997) (3.13333,-0.36800) (3.18333,-0.37993) (3.21667,-0.38781) (3.26667,-0.39952) (3.30000,-0.40725) (3.34093,-0.41667) (3.38333,-0.42634) (3.41667,-0.43387) (3.46667,-0.44508) (3.50000,-0.45248) (3.55000,-0.46349) (3.58333,-0.47077) (3.63333,-0.48159) (3.66667,-0.48874) (3.71667,-0.49938) (3.75000,-0.50641) (3.79904,-0.51667) (3.83333,-0.52378) (3.87975,-0.53333) (3.91667,-0.54087) (3.96175,-0.55000) (4.00000,-0.55769) (4.04507,-0.56667) (4.08333,-0.57423) (4.12973,-0.58333) (4.16667,-0.59052) (4.21575,-0.60000) (4.25000,-0.60656) (4.30000,-0.61607) (4.33333,-0.62236) (4.38333,-0.63173) (4.41667,-0.63792) (4.46667,-0.64715) (4.50000,-0.65326) (4.55000,-0.66235) (4.58333,-0.66837) (4.63333,-0.67733) (4.66706,-0.68333) (4.71667,-0.69210) (4.76174,-0.70000) (4.80000,-0.70666) (4.85000,-0.71530) (4.88333,-0.72102) (4.93333,-0.72955) (4.96667,-0.73519) (5.00000,-0.74081)};
\draw[main,line width=.35pt] plot coordinates {(-5.00000,-0.41422) (-4.95000,-0.40538) (-4.91667,-0.39945) (-4.86667,-0.39048) (-4.82716,-0.38333) (-4.78333,-0.37535) (-4.73609,-0.36667) (-4.70000,-0.35999) (-4.65000,-0.35065) (-4.61667,-0.34438) (-4.56667,-0.33490) (-4.53333,-0.32853) (-4.48333,-0.31890) (-4.45000,-0.31243) (-4.40000,-0.30263) (-4.36667,-0.29605) (-4.31667,-0.28610) (-4.28333,-0.27941) (-4.23333,-0.26928) (-4.20000,-0.26248) (-4.15000,-0.25218) (-4.11667,-0.24525) (-4.06667,-0.23477) (-4.03333,-0.22772) (-3.98333,-0.21705) (-3.95000,-0.20987) (-3.90456,-0.20000) (-3.86667,-0.19170) (-3.82886,-0.18333) (-3.78333,-0.17317) (-3.75000,-0.16567) (-3.70000,-0.15430) (-3.66667,-0.14664) (-3.61667,-0.13505) (-3.58333,-0.12724) (-3.53861,-0.11667) (-3.50000,-0.10745) (-3.46667,-0.09941) (-3.41667,-0.08724) (-3.38333,-0.07903) (-3.33361,-0.06667) (-3.30000,-0.05822) (-3.26667,-0.04976) (-3.21667,-0.03694) (-3.18333,-0.02829) (-3.13898,-0.01667) (-3.10000,-0.00633) (-3.06667,0.00259) (-3.01667,0.01614) (-2.98333,0.02529) (-2.95000,0.03452) (-2.90000,0.04854) (-2.86667,0.05800) (-2.83333,0.06756) (-2.78333,0.08208) (-2.75000,0.09189) (-2.71667,0.10181) (-2.66737,0.11667) (-2.63333,0.12707) (-2.60000,0.13737) (-2.55963,0.15000) (-2.51667,0.16364) (-2.48333,0.17436) (-2.45000,0.18521) (-2.40521,0.20000) (-2.36667,0.21292) (-2.33333,0.22425) (-2.30000,0.23572) (-2.25913,0.25000) (-2.21667,0.26508) (-2.18333,0.27710) (-2.15000,0.28929) (-2.11667,0.30165) (-2.07683,0.31667) (-2.03341,0.33333) (-2.00000,0.34638) (-1.96667,0.35960) (-1.93333,0.37304) (-1.90000,0.38670) (-1.86667,0.40058) (-1.82876,0.41667) (-1.79027,0.43333) (-1.75256,0.45000) (-1.71667,0.46619) (-1.68333,0.48154) (-1.65000,0.49719) (-1.61667,0.51317) (-1.58333,0.52949) (-1.55000,0.54617) (-1.51667,0.56324) (-1.48333,0.58072) (-1.45000,0.59863) (-1.41727,0.61667) (-1.38776,0.63333) (-1.35897,0.65000) (-1.33087,0.66667) (-1.30000,0.68548) (-1.26667,0.70641) (-1.23333,0.72806) (-1.20069,0.75000) (-1.17667,0.76667) (-1.15000,0.78574) (-1.11667,0.81048) (-1.08713,0.83333) (-1.06635,0.85000) (-1.03333,0.87758) (-1.00774,0.90000) (-0.98333,0.92234) (-0.95459,0.95000) (-0.93333,0.97154) (-0.90670,1.00000) (-0.88333,1.02656) (-0.86384,1.05000) (-0.83792,1.08333) (-0.81667,1.11285) (-0.80000,1.13762) (-0.78177,1.16667) (-0.76250,1.20000) (-0.74488,1.23333) (-0.72881,1.26667) (-0.71414,1.30000) (-0.70000,1.33531) (-0.68333,1.38172) (-0.67203,1.41667) (-0.66206,1.45000) (-0.65000,1.49360) (-0.63979,1.53333) (-0.63157,1.56667) (-0.61959,1.61667) (-0.61165,1.65000) (-0.60000,1.69769) (-0.59100,1.73333) (-0.58207,1.76667) (-0.56764,1.81667) (-0.55716,1.85000) (-0.54575,1.88333) (-0.53327,1.91667) (-0.51667,1.95674) (-0.50000,1.99291) (-0.48333,2.02551) (-0.46667,2.05509) (-0.44923,2.08333) (-0.42629,2.11667) (-0.40023,2.15000) (-0.38333,2.16949) (-0.35401,2.20000) (-0.33333,2.21907) (-0.30000,2.24631) (-0.27121,2.26667) (-0.24420,2.28333) (-0.21288,2.30000) (-0.17488,2.31667) (-0.13333,2.33088) (-0.10000,2.33934) (-0.05000,2.34734) (0.00000,2.34995) (0.05000,2.34734) (0.10000,2.33934) (0.13333,2.33088) (0.17488,2.31667) (0.21288,2.30000) (0.24420,2.28333) (0.27121,2.26667) (0.30000,2.24631) (0.33333,2.21907) (0.35401,2.20000) (0.38333,2.16949) (0.40023,2.15000) (0.42629,2.11667) (0.44923,2.08333) (0.46667,2.05509) (0.48333,2.02551) (0.50000,1.99291) (0.51667,1.95674) (0.53327,1.91667) (0.54575,1.88333) (0.55716,1.85000) (0.56764,1.81667) (0.58207,1.76667) (0.59100,1.73333) (0.60000,1.69769) (0.61165,1.65000) (0.61959,1.61667) (0.63157,1.56667) (0.63979,1.53333) (0.65000,1.49360) (0.66206,1.45000) (0.67203,1.41667) (0.68333,1.38172) (0.70000,1.33531) (0.71414,1.30000) (0.72881,1.26667) (0.74488,1.23333) (0.76250,1.20000) (0.78177,1.16667) (0.80000,1.13762) (0.81667,1.11285) (0.83792,1.08333) (0.86384,1.05000) (0.88333,1.02656) (0.90670,1.00000) (0.93333,0.97154) (0.95459,0.95000) (0.98333,0.92234) (1.00774,0.90000) (1.03333,0.87758) (1.06635,0.85000) (1.08713,0.83333) (1.11667,0.81048) (1.15000,0.78574) (1.17667,0.76667) (1.20069,0.75000) (1.23333,0.72806) (1.26667,0.70641) (1.30000,0.68548) (1.33087,0.66667) (1.35897,0.65000) (1.38776,0.63333) (1.41727,0.61667) (1.45000,0.59863) (1.48333,0.58072) (1.51667,0.56324) (1.55000,0.54617) (1.58333,0.52949) (1.61667,0.51317) (1.65000,0.49719) (1.68333,0.48154) (1.71667,0.46619) (1.75256,0.45000) (1.79027,0.43333) (1.82876,0.41667) (1.86667,0.40058) (1.90000,0.38670) (1.93333,0.37304) (1.96667,0.35960) (2.00000,0.34638) (2.03341,0.33333) (2.07683,0.31667) (2.11667,0.30165) (2.15000,0.28929) (2.18333,0.27710) (2.21667,0.26508) (2.25913,0.25000) (2.30000,0.23572) (2.33333,0.22425) (2.36667,0.21292) (2.40521,0.20000) (2.45000,0.18521) (2.48333,0.17436) (2.51667,0.16364) (2.55963,0.15000) (2.60000,0.13737) (2.63333,0.12707) (2.66737,0.11667) (2.71667,0.10181) (2.75000,0.09189) (2.78333,0.08208) (2.83333,0.06756) (2.86667,0.05800) (2.90000,0.04854) (2.95000,0.03452) (2.98333,0.02529) (3.01667,0.01614) (3.06667,0.00259) (3.10000,-0.00633) (3.13898,-0.01667) (3.18333,-0.02829) (3.21667,-0.03694) (3.26667,-0.04976) (3.30000,-0.05822) (3.33361,-0.06667) (3.38333,-0.07903) (3.41667,-0.08724) (3.46667,-0.09941) (3.50000,-0.10745) (3.53861,-0.11667) (3.58333,-0.12724) (3.61667,-0.13505) (3.66667,-0.14664) (3.70000,-0.15430) (3.75000,-0.16567) (3.78333,-0.17317) (3.82886,-0.18333) (3.86667,-0.19170) (3.90456,-0.20000) (3.95000,-0.20987) (3.98333,-0.21705) (4.03333,-0.22772) (4.06667,-0.23477) (4.11667,-0.24525) (4.15000,-0.25218) (4.20000,-0.26248) (4.23333,-0.26928) (4.28333,-0.27941) (4.31667,-0.28610) (4.36667,-0.29605) (4.40000,-0.30263) (4.45000,-0.31243) (4.48333,-0.31890) (4.53333,-0.32853) (4.56667,-0.33490) (4.61667,-0.34438) (4.65000,-0.35065) (4.70000,-0.35999) (4.73609,-0.36667) (4.78333,-0.37535) (4.82716,-0.38333) (4.86667,-0.39048) (4.91667,-0.39945) (4.95000,-0.40538) (5.00000,-0.41422)};
\draw[main,line width=.35pt] plot coordinates {(-5.00000,-0.08995) (-4.96402,-0.08333) (-4.91667,-0.07456) (-4.87448,-0.06667) (-4.83333,-0.05891) (-4.78646,-0.05000) (-4.75000,-0.04301) (-4.70000,-0.03335) (-4.66667,-0.02685) (-4.61667,-0.01702) (-4.58333,-0.01041) (-4.53333,-0.00041) (-4.50000,0.00631) (-4.45000,0.01649) (-4.41667,0.02333) (-4.36836,0.03333) (-4.33333,0.04065) (-4.28900,0.05000) (-4.25000,0.05829) (-4.21099,0.06667) (-4.16667,0.07627) (-4.13333,0.08355) (-4.08333,0.09458) (-4.05000,0.10201) (-4.00000,0.11326) (-3.96667,0.12084) (-3.91667,0.13231) (-3.88333,0.14004) (-3.84079,0.15000) (-3.80000,0.15964) (-3.76667,0.16760) (-3.71667,0.17966) (-3.68333,0.18779) (-3.63382,0.20000) (-3.60000,0.20843) (-3.56667,0.21681) (-3.51667,0.22953) (-3.48333,0.23811) (-3.43762,0.25000) (-3.40000,0.25990) (-3.36667,0.26877) (-3.31667,0.28222) (-3.28333,0.29130) (-3.25000,0.30048) (-3.20000,0.31441) (-3.16667,0.32382) (-3.13331,0.33333) (-3.08333,0.34778) (-3.05000,0.35754) (-3.01667,0.36741) (-2.96667,0.38242) (-2.93333,0.39256) (-2.90000,0.40283) (-2.85568,0.41667) (-2.81667,0.42902) (-2.78333,0.43972) (-2.75000,0.45055) (-2.70113,0.46667) (-2.66667,0.47821) (-2.63333,0.48953) (-2.60000,0.50100) (-2.55517,0.51667) (-2.51667,0.53035) (-2.48333,0.54238) (-2.45000,0.55458) (-2.41667,0.56695) (-2.37331,0.58333) (-2.33333,0.59872) (-2.30000,0.61176) (-2.26667,0.62502) (-2.23333,0.63850) (-2.20000,0.65220) (-2.16540,0.66667) (-2.12629,0.68333) (-2.08798,0.70000) (-2.05046,0.71667) (-2.01667,0.73198) (-1.98333,0.74738) (-1.95000,0.76309) (-1.91667,0.77913) (-1.88333,0.79550) (-1.85000,0.81223) (-1.81667,0.82934) (-1.78333,0.84684) (-1.75000,0.86477) (-1.71667,0.88313) (-1.68678,0.90000) (-1.65791,0.91667) (-1.62970,0.93333) (-1.60000,0.95130) (-1.56667,0.97202) (-1.53333,0.99337) (-1.50000,1.01539) (-1.47363,1.03333) (-1.44971,1.05000) (-1.41667,1.07372) (-1.38333,1.09853) (-1.35975,1.11667) (-1.33333,1.13755) (-1.30000,1.16490) (-1.27830,1.18333) (-1.25000,1.20815) (-1.22240,1.23333) (-1.20000,1.25447) (-1.17059,1.28333) (-1.15000,1.30430) (-1.12258,1.33333) (-1.10000,1.35818) (-1.07803,1.38333) (-1.05007,1.41667) (-1.03333,1.43739) (-1.01054,1.46667) (-0.98561,1.50000) (-0.96667,1.52626) (-0.95000,1.55013) (-0.92754,1.58333) (-0.90571,1.61667) (-0.88451,1.65000) (-0.86667,1.67868) (-0.85000,1.70598) (-0.83333,1.73370) (-0.81371,1.76667) (-0.79396,1.80000) (-0.77419,1.83333) (-0.75427,1.86667) (-0.73405,1.90000) (-0.71667,1.92808) (-0.70000,1.95451) (-0.68140,1.98333) (-0.65908,2.01667) (-0.63569,2.05000) (-0.61667,2.07584) (-0.59825,2.10000) (-0.57134,2.13333) (-0.55000,2.15807) (-0.52711,2.18333) (-0.50000,2.21107) (-0.47683,2.23333) (-0.45000,2.25715) (-0.41795,2.28333) (-0.39574,2.30000) (-0.36667,2.32007) (-0.33333,2.34082) (-0.30000,2.35924) (-0.26667,2.37546) (-0.23333,2.38955) (-0.20000,2.40159) (-0.15000,2.41599) (-0.11667,2.42327) (-0.06667,2.43067) (-0.03307,2.43333) (0.01667,2.43400) (0.05000,2.43222) (0.10000,2.42620) (0.14709,2.41667) (0.18333,2.40690) (0.21667,2.39582) (0.25000,2.38274) (0.28524,2.36667) (0.31722,2.35000) (0.35000,2.33074) (0.38333,2.30878) (0.41667,2.28432) (0.43876,2.26667) (0.46667,2.24253) (0.49439,2.21667) (0.51667,2.19420) (0.54244,2.16667) (0.56667,2.13878) (0.58500,2.11667) (0.61112,2.08333) (0.63333,2.05319) (0.65000,2.02969) (0.67036,2.00000) (0.69229,1.96667) (0.71345,1.93333) (0.73333,1.90113) (0.75000,1.87363) (0.76667,1.84585) (0.78406,1.81667) (0.80383,1.78333) (0.82363,1.75000) (0.84360,1.71667) (0.86385,1.68333) (0.88333,1.65184) (0.90000,1.62548) (0.91667,1.59977) (0.93872,1.56667) (0.96171,1.53333) (0.98333,1.50308) (1.00000,1.48055) (1.02342,1.45000) (1.05000,1.41675) (1.06667,1.39666) (1.09252,1.36667) (1.11667,1.33973) (1.13819,1.31667) (1.16667,1.28726) (1.18743,1.26667) (1.21667,1.23866) (1.24056,1.21667) (1.26667,1.19341) (1.29790,1.16667) (1.31800,1.15000) (1.35000,1.12429) (1.38141,1.10000) (1.40362,1.08333) (1.43333,1.06165) (1.46667,1.03813) (1.49811,1.01667) (1.52320,1.00000) (1.55000,0.98262) (1.58333,0.96158) (1.61667,0.94115) (1.65000,0.92129) (1.68333,0.90196) (1.71631,0.88333) (1.74652,0.86667) (1.77742,0.85000) (1.80901,0.83333) (1.84130,0.81667) (1.87432,0.80000) (1.90805,0.78333) (1.94253,0.76667) (1.97774,0.75000) (2.01372,0.73333) (2.05000,0.71687) (2.08333,0.70204) (2.11667,0.68748) (2.15000,0.67318) (2.18333,0.65913) (2.21667,0.64532) (2.25000,0.63173) (2.28762,0.61667) (2.33004,0.60000) (2.36667,0.58587) (2.40000,0.57321) (2.43333,0.56074) (2.46667,0.54846) (2.50836,0.53333) (2.55000,0.51849) (2.58333,0.50679) (2.61667,0.49524) (2.65154,0.48333) (2.70000,0.46704) (2.73333,0.45601) (2.76667,0.44511) (2.80318,0.43333) (2.85000,0.41845) (2.88333,0.40801) (2.91667,0.39768) (2.96365,0.38333) (3.00000,0.37238) (3.03333,0.36246) (3.07572,0.35000) (3.11667,0.33812) (3.15000,0.32856) (3.19199,0.31667) (3.23333,0.30510) (3.26667,0.29588) (3.31258,0.28333) (3.35000,0.27323) (3.38333,0.26432) (3.43333,0.25112) (3.46667,0.24242) (3.50185,0.23333) (3.55000,0.22103) (3.58333,0.21261) (3.63333,0.20012) (3.66667,0.19188) (3.70159,0.18333) (3.75000,0.17161) (3.78333,0.16361) (3.83333,0.15175) (3.86667,0.14393) (3.91226,0.13333) (3.95000,0.12465) (3.98499,0.11667) (4.03333,0.10575) (4.06667,0.09829) (4.11667,0.08722) (4.15000,0.07990) (4.20000,0.06904) (4.23333,0.06186) (4.28333,0.05120) (4.31667,0.04415) (4.36667,0.03369) (4.40000,0.02677) (4.44912,0.01667) (4.48333,0.00969) (4.53129,0.00000) (4.56667,-0.00709) (4.61488,-0.01667) (4.65000,-0.02358) (4.69993,-0.03333) (4.73333,-0.03980) (4.78333,-0.04940) (4.81667,-0.05575) (4.86667,-0.06520) (4.90000,-0.07145) (4.95000,-0.08074) (4.98333,-0.08689) (5.00000,-0.08995)};
\draw[main,line width=.35pt] plot coordinates {(-5.00000,0.23163) (-4.96667,0.23798) (-4.91667,0.24757) (-4.88333,0.25402) (-4.83333,0.26378) (-4.80000,0.27034) (-4.75000,0.28027) (-4.71667,0.28694) (-4.66667,0.29705) (-4.63333,0.30384) (-4.58333,0.31413) (-4.55000,0.32104) (-4.50000,0.33152) (-4.46667,0.33856) (-4.41667,0.34923) (-4.38333,0.35641) (-4.33618,0.36667) (-4.30000,0.37461) (-4.26061,0.38333) (-4.21667,0.39316) (-4.18333,0.40068) (-4.13333,0.41208) (-4.10000,0.41976) (-4.05000,0.43140) (-4.01667,0.43924) (-3.97136,0.45000) (-3.93333,0.45913) (-3.90000,0.46720) (-3.85000,0.47945) (-3.81667,0.48771) (-3.76758,0.50000) (-3.73333,0.50867) (-3.70000,0.51720) (-3.65000,0.53013) (-3.61667,0.53885) (-3.57453,0.55000) (-3.53333,0.56103) (-3.50000,0.57005) (-3.45155,0.58333) (-3.41667,0.59301) (-3.38333,0.60236) (-3.33333,0.61656) (-3.30000,0.62616) (-3.26667,0.63586) (-3.21870,0.65000) (-3.18333,0.66057) (-3.15000,0.67065) (-3.10858,0.68333) (-3.06667,0.69635) (-3.03333,0.70684) (-3.00000,0.71746) (-2.95090,0.73333) (-2.91667,0.74457) (-2.88333,0.75565) (-2.85000,0.76688) (-2.80190,0.78333) (-2.76667,0.79558) (-2.73333,0.80733) (-2.70000,0.81925) (-2.66118,0.83333) (-2.61667,0.84977) (-2.58333,0.86228) (-2.55000,0.87498) (-2.51667,0.88787) (-2.48333,0.90096) (-2.44402,0.91667) (-2.40307,0.93333) (-2.36667,0.94843) (-2.33333,0.96249) (-2.30000,0.97680) (-2.26667,0.99135) (-2.23333,1.00616) (-2.20000,1.02124) (-2.16667,1.03660) (-2.13333,1.05225) (-2.10000,1.06820) (-2.06667,1.08447) (-2.03333,1.10106) (-2.00000,1.11799) (-1.96667,1.13529) (-1.93333,1.15295) (-1.90000,1.17101) (-1.86667,1.18948) (-1.83333,1.20838) (-1.80000,1.22773) (-1.76667,1.24755) (-1.73529,1.26667) (-1.70853,1.28333) (-1.68230,1.30000) (-1.65000,1.32100) (-1.61667,1.34325) (-1.58333,1.36614) (-1.55891,1.38333) (-1.53333,1.40169) (-1.50000,1.42625) (-1.46872,1.45000) (-1.44726,1.46667) (-1.41667,1.49097) (-1.38528,1.51667) (-1.36537,1.53333) (-1.33333,1.56080) (-1.30779,1.58333) (-1.28333,1.60534) (-1.25307,1.63333) (-1.23333,1.65196) (-1.20088,1.68333) (-1.18333,1.70065) (-1.15088,1.73333) (-1.13333,1.75132) (-1.10268,1.78333) (-1.08333,1.80380) (-1.05589,1.83333) (-1.03333,1.85780) (-1.01012,1.88333) (-0.98333,1.91287) (-0.96493,1.93333) (-0.93493,1.96667) (-0.91667,1.98688) (-0.88976,2.01667) (-0.86667,2.04189) (-0.84389,2.06667) (-0.81667,2.09567) (-0.79673,2.11667) (-0.76667,2.14750) (-0.74760,2.16667) (-0.71667,2.19677) (-0.69564,2.21667) (-0.66667,2.24303) (-0.63972,2.26667) (-0.61667,2.28599) (-0.58333,2.31268) (-0.55612,2.33333) (-0.53298,2.35000) (-0.50000,2.37244) (-0.46667,2.39354) (-0.43333,2.41306) (-0.40000,2.43100) (-0.36667,2.44740) (-0.33333,2.46228) (-0.30000,2.47566) (-0.26667,2.48755) (-0.22616,2.50000) (-0.18333,2.51097) (-0.15000,2.51784) (-0.10000,2.52558) (-0.05000,2.53018) (0.00000,2.53171) (0.05000,2.53018) (0.10000,2.52558) (0.15000,2.51784) (0.18333,2.51097) (0.22616,2.50000) (0.26667,2.48755) (0.30000,2.47566) (0.33333,2.46228) (0.36667,2.44740) (0.40000,2.43100) (0.43333,2.41306) (0.46667,2.39354) (0.50000,2.37244) (0.53298,2.35000) (0.55612,2.33333) (0.58333,2.31268) (0.61667,2.28599) (0.63972,2.26667) (0.66667,2.24303) (0.69564,2.21667) (0.71667,2.19677) (0.74760,2.16667) (0.76667,2.14750) (0.79673,2.11667) (0.81667,2.09567) (0.84389,2.06667) (0.86667,2.04189) (0.88976,2.01667) (0.91667,1.98688) (0.93493,1.96667) (0.96493,1.93333) (0.98333,1.91287) (1.01012,1.88333) (1.03333,1.85780) (1.05589,1.83333) (1.08333,1.80380) (1.10268,1.78333) (1.13333,1.75132) (1.15088,1.73333) (1.18333,1.70065) (1.20088,1.68333) (1.23333,1.65196) (1.25307,1.63333) (1.28333,1.60534) (1.30779,1.58333) (1.33333,1.56080) (1.36537,1.53333) (1.38528,1.51667) (1.41667,1.49097) (1.44726,1.46667) (1.46872,1.45000) (1.50000,1.42625) (1.53333,1.40169) (1.55891,1.38333) (1.58333,1.36614) (1.61667,1.34325) (1.65000,1.32100) (1.68230,1.30000) (1.70853,1.28333) (1.73529,1.26667) (1.76667,1.24755) (1.80000,1.22773) (1.83333,1.20838) (1.86667,1.18948) (1.90000,1.17101) (1.93333,1.15295) (1.96667,1.13529) (2.00000,1.11799) (2.03333,1.10106) (2.06667,1.08447) (2.10000,1.06820) (2.13333,1.05225) (2.16667,1.03660) (2.20000,1.02124) (2.23333,1.00616) (2.26667,0.99135) (2.30000,0.97680) (2.33333,0.96249) (2.36667,0.94843) (2.40307,0.93333) (2.44402,0.91667) (2.48333,0.90096) (2.51667,0.88787) (2.55000,0.87498) (2.58333,0.86228) (2.61667,0.84977) (2.66118,0.83333) (2.70000,0.81925) (2.73333,0.80733) (2.76667,0.79558) (2.80190,0.78333) (2.85000,0.76688) (2.88333,0.75565) (2.91667,0.74457) (2.95090,0.73333) (3.00000,0.71746) (3.03333,0.70684) (3.06667,0.69635) (3.10858,0.68333) (3.15000,0.67065) (3.18333,0.66057) (3.21870,0.65000) (3.26667,0.63586) (3.30000,0.62616) (3.33333,0.61656) (3.38333,0.60236) (3.41667,0.59301) (3.45155,0.58333) (3.50000,0.57005) (3.53333,0.56103) (3.57453,0.55000) (3.61667,0.53885) (3.65000,0.53013) (3.70000,0.51720) (3.73333,0.50867) (3.76758,0.50000) (3.81667,0.48771) (3.85000,0.47945) (3.90000,0.46720) (3.93333,0.45913) (3.97136,0.45000) (4.01667,0.43924) (4.05000,0.43140) (4.10000,0.41976) (4.13333,0.41208) (4.18333,0.40068) (4.21667,0.39316) (4.26061,0.38333) (4.30000,0.37461) (4.33618,0.36667) (4.38333,0.35641) (4.41667,0.34923) (4.46667,0.33856) (4.50000,0.33152) (4.55000,0.32104) (4.58333,0.31413) (4.63333,0.30384) (4.66667,0.29705) (4.71667,0.28694) (4.75000,0.28027) (4.80000,0.27034) (4.83333,0.26378) (4.88333,0.25402) (4.91667,0.24757) (4.96667,0.23798) (5.00000,0.23163)};
\draw[main,line width=.35pt] plot coordinates {(-5.00000,0.55020) (-4.95000,0.55999) (-4.91623,0.56667) (-4.86667,0.57655) (-4.83292,0.58333) (-4.78333,0.59339) (-4.75000,0.60022) (-4.70000,0.61054) (-4.66667,0.61749) (-4.61667,0.62801) (-4.58333,0.63509) (-4.53333,0.64580) (-4.50000,0.65302) (-4.45000,0.66394) (-4.41667,0.67129) (-4.36667,0.68243) (-4.33333,0.68993) (-4.28901,0.70000) (-4.25000,0.70895) (-4.21665,0.71667) (-4.16667,0.72835) (-4.13333,0.73623) (-4.08333,0.74817) (-4.05000,0.75622) (-4.00715,0.76667) (-3.96667,0.77664) (-3.93333,0.78494) (-3.88333,0.79752) (-3.85000,0.80601) (-3.80858,0.81667) (-3.76667,0.82757) (-3.73333,0.83634) (-3.68333,0.84965) (-3.65000,0.85862) (-3.61667,0.86769) (-3.56667,0.88146) (-3.53333,0.89076) (-3.50000,0.90015) (-3.45000,0.91442) (-3.41667,0.92406) (-3.38333,0.93380) (-3.33333,0.94860) (-3.30000,0.95861) (-3.26667,0.96873) (-3.21920,0.98333) (-3.18333,0.99452) (-3.15000,1.00505) (-3.11366,1.01667) (-3.06667,1.03191) (-3.03333,1.04288) (-3.00000,1.05399) (-2.96245,1.06667) (-2.91667,1.08236) (-2.88333,1.09397) (-2.85000,1.10572) (-2.81667,1.11764) (-2.77343,1.13333) (-2.73333,1.14813) (-2.70000,1.16062) (-2.66667,1.17328) (-2.63333,1.18613) (-2.59790,1.20000) (-2.55602,1.21667) (-2.51667,1.23261) (-2.48333,1.24634) (-2.45000,1.26029) (-2.41667,1.27446) (-2.38333,1.28885) (-2.35000,1.30348) (-2.31667,1.31835) (-2.28333,1.33348) (-2.24755,1.35000) (-2.21211,1.36667) (-2.17731,1.38333) (-2.14315,1.40000) (-2.10960,1.41667) (-2.07667,1.43333) (-2.04432,1.45000) (-2.01255,1.46667) (-1.98135,1.48333) (-1.95000,1.50038) (-1.91667,1.51885) (-1.88333,1.53770) (-1.85000,1.55692) (-1.81667,1.57653) (-1.78333,1.59655) (-1.75052,1.61667) (-1.72384,1.63333) (-1.69756,1.65000) (-1.66667,1.66993) (-1.63333,1.69188) (-1.60000,1.71430) (-1.57223,1.73333) (-1.54827,1.75000) (-1.51667,1.77233) (-1.48333,1.79636) (-1.45569,1.81667) (-1.43326,1.83333) (-1.40000,1.85842) (-1.36758,1.88333) (-1.34617,1.90000) (-1.31667,1.92318) (-1.28333,1.94977) (-1.26238,1.96667) (-1.23333,1.99023) (-1.20111,2.01667) (-1.18090,2.03333) (-1.15000,2.05890) (-1.12067,2.08333) (-1.10000,2.10053) (-1.06667,2.12826) (-1.04055,2.15000) (-1.01667,2.16971) (-0.98333,2.19709) (-0.95934,2.21667) (-0.93333,2.23758) (-0.90000,2.26406) (-0.87536,2.28333) (-0.85000,2.30277) (-0.81667,2.32780) (-0.78635,2.35000) (-0.76297,2.36667) (-0.73333,2.38719) (-0.70000,2.40949) (-0.66667,2.43091) (-0.63564,2.45000) (-0.60736,2.46667) (-0.57777,2.48333) (-0.54664,2.50000) (-0.51369,2.51667) (-0.47850,2.53333) (-0.44053,2.55000) (-0.40000,2.56628) (-0.36667,2.57849) (-0.33333,2.58963) (-0.29886,2.60000) (-0.25000,2.61279) (-0.21667,2.62019) (-0.16667,2.62930) (-0.13333,2.63403) (-0.08333,2.63919) (-0.03333,2.64196) (0.01667,2.64236) (0.06667,2.64038) (0.11667,2.63602) (0.15000,2.63180) (0.20000,2.62350) (0.23333,2.61661) (0.28333,2.60433) (0.31667,2.59480) (0.35251,2.58333) (0.39898,2.56667) (0.43333,2.55302) (0.46667,2.53869) (0.50000,2.52329) (0.53333,2.50685) (0.56667,2.48937) (0.60000,2.47087) (0.63333,2.45138) (0.66282,2.43333) (0.68903,2.41667) (0.71667,2.39846) (0.75000,2.37570) (0.78333,2.35215) (0.80924,2.33333) (0.83333,2.31538) (0.86667,2.29000) (0.89672,2.26667) (0.91780,2.25000) (0.95000,2.22417) (0.97982,2.20000) (1.00014,2.18333) (1.03333,2.15593) (1.06062,2.13333) (1.08333,2.11440) (1.11667,2.08663) (1.14071,2.06667) (1.16667,2.04508) (1.20000,2.01757) (1.22143,2.00000) (1.25000,1.97666) (1.28305,1.95000) (1.30392,1.93333) (1.33333,1.91003) (1.36667,1.88404) (1.38925,1.86667) (1.41667,1.84579) (1.45000,1.82085) (1.47838,1.80000) (1.50135,1.78333) (1.53333,1.76049) (1.56667,1.73716) (1.59654,1.71667) (1.62121,1.70000) (1.65000,1.68085) (1.68333,1.65912) (1.71667,1.63783) (1.75000,1.61699) (1.77769,1.60000) (1.80530,1.58333) (1.83336,1.56667) (1.86667,1.54726) (1.90000,1.52823) (1.93333,1.50956) (1.96667,1.49127) (2.00000,1.47332) (2.03333,1.45572) (2.06667,1.43844) (2.10000,1.42148) (2.13333,1.40483) (2.16667,1.38848) (2.20000,1.37242) (2.23333,1.35664) (2.26667,1.34113) (2.30000,1.32588) (2.33333,1.31088) (2.36667,1.29614) (2.40000,1.28163) (2.43494,1.26667) (2.47456,1.25000) (2.51491,1.23333) (2.55000,1.21908) (2.58333,1.20576) (2.61667,1.19262) (2.65000,1.17969) (2.68403,1.16667) (2.72832,1.15000) (2.76667,1.13581) (2.80000,1.12365) (2.83333,1.11166) (2.86667,1.09983) (2.91388,1.08333) (2.95000,1.07090) (2.98333,1.05959) (3.01667,1.04842) (3.06233,1.03333) (3.10000,1.02107) (3.13333,1.01036) (3.16667,0.99977) (3.21667,0.98412) (3.25000,0.97383) (3.28333,0.96365) (3.32867,0.95000) (3.36667,0.93870) (3.40000,0.92891) (3.44221,0.91667) (3.48333,0.90488) (3.51667,0.89544) (3.55994,0.88333) (3.60000,0.87226) (3.63333,0.86315) (3.68202,0.85000) (3.71667,0.84075) (3.75000,0.83194) (3.80000,0.81889) (3.83333,0.81028) (3.87359,0.80000) (3.91667,0.78911) (3.95000,0.78078) (4.00000,0.76842) (4.03333,0.76027) (4.07575,0.75000) (4.11667,0.74019) (4.15000,0.73228) (4.20000,0.72054) (4.23333,0.71280) (4.28333,0.70129) (4.31667,0.69370) (4.36266,0.68333) (4.40000,0.67499) (4.43762,0.66667) (4.48333,0.65664) (4.51667,0.64940) (4.56667,0.63865) (4.60000,0.63154) (4.65000,0.62098) (4.68333,0.61401) (4.73333,0.60365) (4.76667,0.59680) (4.81667,0.58662) (4.85000,0.57989) (4.90000,0.56989) (4.93333,0.56328) (4.98333,0.55345) (5.00000,0.55020)};
\draw[main,line width=.35pt] plot coordinates {(-5.00000,0.86546) (-4.96667,0.87211) (-4.91667,0.88217) (-4.88333,0.88894) (-4.83333,0.89918) (-4.80000,0.90606) (-4.75000,0.91649) (-4.71667,0.92351) (-4.67040,0.93333) (-4.63333,0.94127) (-4.59299,0.95000) (-4.55000,0.95938) (-4.51667,0.96672) (-4.46667,0.97783) (-4.43333,0.98532) (-4.38333,0.99666) (-4.35000,1.00429) (-4.30000,1.01586) (-4.26667,1.02365) (-4.22561,1.03333) (-4.18333,1.04341) (-4.15000,1.05143) (-4.10000,1.06359) (-4.06667,1.07178) (-4.02017,1.08333) (-3.98333,1.09258) (-3.95000,1.10103) (-3.90000,1.11385) (-3.86667,1.12250) (-3.82531,1.13333) (-3.78333,1.14446) (-3.75000,1.15339) (-3.70101,1.16667) (-3.66667,1.17609) (-3.63333,1.18532) (-3.58333,1.19935) (-3.55000,1.20881) (-3.51667,1.21838) (-3.46667,1.23290) (-3.43333,1.24271) (-3.40000,1.25263) (-3.35341,1.26667) (-3.31667,1.27788) (-3.28333,1.28817) (-3.24549,1.30000) (-3.20000,1.31441) (-3.16667,1.32511) (-3.13333,1.33593) (-3.09059,1.35000) (-3.05000,1.36354) (-3.01667,1.37482) (-2.98333,1.38623) (-2.94365,1.40000) (-2.90000,1.41537) (-2.86667,1.42728) (-2.83333,1.43934) (-2.80000,1.45155) (-2.75936,1.46667) (-2.71667,1.48279) (-2.68333,1.49557) (-2.65000,1.50853) (-2.61667,1.52166) (-2.58333,1.53497) (-2.54626,1.55000) (-2.50579,1.56667) (-2.46667,1.58305) (-2.43333,1.59722) (-2.40000,1.61160) (-2.36667,1.62619) (-2.33333,1.64100) (-2.30000,1.65603) (-2.26667,1.67128) (-2.23333,1.68676) (-2.20000,1.70248) (-2.16667,1.71843) (-2.13333,1.73464) (-2.10000,1.75108) (-2.06667,1.76779) (-2.03333,1.78475) (-2.00000,1.80197) (-1.96667,1.81946) (-1.93333,1.83722) (-1.90000,1.85524) (-1.86667,1.87355) (-1.83333,1.89213) (-1.80000,1.91098) (-1.76667,1.93012) (-1.73333,1.94953) (-1.70429,1.96667) (-1.67637,1.98333) (-1.64873,2.00000) (-1.61667,2.01955) (-1.58333,2.04013) (-1.55000,2.06096) (-1.51667,2.08202) (-1.48849,2.10000) (-1.46254,2.11667) (-1.43333,2.13554) (-1.40000,2.15724) (-1.36667,2.17910) (-1.33497,2.20000) (-1.30977,2.21667) (-1.28333,2.23416) (-1.25000,2.25625) (-1.21667,2.27836) (-1.18399,2.30000) (-1.15875,2.31667) (-1.13333,2.33335) (-1.10000,2.35510) (-1.06667,2.37668) (-1.03333,2.39803) (-1.00388,2.41667) (-0.97719,2.43333) (-0.95000,2.45005) (-0.91667,2.47021) (-0.88333,2.48993) (-0.85000,2.50918) (-0.81667,2.52793) (-0.78333,2.54612) (-0.75000,2.56374) (-0.71667,2.58074) (-0.68333,2.59711) (-0.65000,2.61280) (-0.61667,2.62781) (-0.58333,2.64212) (-0.55000,2.65570) (-0.51667,2.66853) (-0.47547,2.68333) (-0.43333,2.69729) (-0.40000,2.70743) (-0.36667,2.71677) (-0.31667,2.72932) (-0.28333,2.73667) (-0.23333,2.74618) (-0.20000,2.75149) (-0.15000,2.75796) (-0.10000,2.76257) (-0.05000,2.76533) (0.00000,2.76625) (0.05000,2.76533) (0.10000,2.76257) (0.15000,2.75796) (0.20000,2.75149) (0.23333,2.74618) (0.28333,2.73667) (0.31667,2.72932) (0.36667,2.71677) (0.40000,2.70743) (0.43333,2.69729) (0.47547,2.68333) (0.51667,2.66853) (0.55000,2.65570) (0.58333,2.64212) (0.61667,2.62781) (0.65000,2.61280) (0.68333,2.59711) (0.71667,2.58074) (0.75000,2.56374) (0.78333,2.54612) (0.81667,2.52793) (0.85000,2.50918) (0.88333,2.48993) (0.91667,2.47021) (0.95000,2.45005) (0.97719,2.43333) (1.00388,2.41667) (1.03333,2.39803) (1.06667,2.37668) (1.10000,2.35510) (1.13333,2.33335) (1.15875,2.31667) (1.18399,2.30000) (1.21667,2.27836) (1.25000,2.25625) (1.28333,2.23416) (1.30977,2.21667) (1.33497,2.20000) (1.36667,2.17910) (1.40000,2.15724) (1.43333,2.13554) (1.46254,2.11667) (1.48849,2.10000) (1.51667,2.08202) (1.55000,2.06096) (1.58333,2.04013) (1.61667,2.01955) (1.64873,2.00000) (1.67637,1.98333) (1.70429,1.96667) (1.73333,1.94953) (1.76667,1.93012) (1.80000,1.91098) (1.83333,1.89213) (1.86667,1.87355) (1.90000,1.85524) (1.93333,1.83722) (1.96667,1.81946) (2.00000,1.80197) (2.03333,1.78475) (2.06667,1.76779) (2.10000,1.75108) (2.13333,1.73464) (2.16667,1.71843) (2.20000,1.70248) (2.23333,1.68676) (2.26667,1.67128) (2.30000,1.65603) (2.33333,1.64100) (2.36667,1.62619) (2.40000,1.61160) (2.43333,1.59722) (2.46667,1.58305) (2.50579,1.56667) (2.54626,1.55000) (2.58333,1.53497) (2.61667,1.52166) (2.65000,1.50853) (2.68333,1.49557) (2.71667,1.48279) (2.75936,1.46667) (2.80000,1.45155) (2.83333,1.43934) (2.86667,1.42728) (2.90000,1.41537) (2.94365,1.40000) (2.98333,1.38623) (3.01667,1.37482) (3.05000,1.36354) (3.09059,1.35000) (3.13333,1.33593) (3.16667,1.32511) (3.20000,1.31441) (3.24549,1.30000) (3.28333,1.28817) (3.31667,1.27788) (3.35341,1.26667) (3.40000,1.25263) (3.43333,1.24271) (3.46667,1.23290) (3.51667,1.21838) (3.55000,1.20881) (3.58333,1.19935) (3.63333,1.18532) (3.66667,1.17609) (3.70101,1.16667) (3.75000,1.15339) (3.78333,1.14446) (3.82531,1.13333) (3.86667,1.12250) (3.90000,1.11385) (3.95000,1.10103) (3.98333,1.09258) (4.02017,1.08333) (4.06667,1.07178) (4.10000,1.06359) (4.15000,1.05143) (4.18333,1.04341) (4.22561,1.03333) (4.26667,1.02365) (4.30000,1.01586) (4.35000,1.00429) (4.38333,0.99666) (4.43333,0.98532) (4.46667,0.97783) (4.51667,0.96672) (4.55000,0.95938) (4.59299,0.95000) (4.63333,0.94127) (4.67040,0.93333) (4.71667,0.92351) (4.75000,0.91649) (4.80000,0.90606) (4.83333,0.89918) (4.88333,0.88894) (4.91667,0.88217) (4.96667,0.87211) (5.00000,0.86546)};
\draw[main,line width=.35pt] plot coordinates {(-5.00000,1.38301) (-4.96667,1.38976) (-4.91667,1.39998) (-4.88333,1.40685) (-4.83612,1.41667) (-4.80000,1.42424) (-4.75703,1.43333) (-4.71667,1.44195) (-4.67926,1.45000) (-4.63333,1.45998) (-4.60000,1.46728) (-4.55000,1.47834) (-4.51667,1.48579) (-4.46667,1.49706) (-4.43333,1.50464) (-4.38333,1.51613) (-4.35000,1.52387) (-4.30958,1.53333) (-4.26667,1.54347) (-4.23333,1.55143) (-4.18333,1.56348) (-4.15000,1.57159) (-4.10225,1.58333) (-4.06667,1.59217) (-4.03333,1.60053) (-3.98333,1.61319) (-3.95000,1.62172) (-3.90509,1.63333) (-3.86667,1.64337) (-3.83333,1.65216) (-3.78333,1.66550) (-3.75000,1.67449) (-3.71667,1.68356) (-3.66667,1.69732) (-3.63333,1.70660) (-3.59751,1.71667) (-3.55000,1.73017) (-3.51667,1.73975) (-3.48138,1.75000) (-3.43333,1.76411) (-3.40000,1.77402) (-3.36667,1.78402) (-3.31667,1.79921) (-3.28333,1.80946) (-3.25000,1.81981) (-3.20695,1.83333) (-3.16667,1.84614) (-3.13333,1.85686) (-3.10000,1.86769) (-3.05243,1.88333) (-3.01667,1.89524) (-2.98333,1.90646) (-2.95000,1.91781) (-2.90492,1.93333) (-2.86667,1.94667) (-2.83333,1.95843) (-2.80000,1.97032) (-2.76392,1.98333) (-2.71829,2.00000) (-2.68333,2.01293) (-2.65000,2.02539) (-2.61667,2.03799) (-2.58333,2.05073) (-2.54211,2.06667) (-2.50000,2.08316) (-2.46667,2.09637) (-2.43333,2.10972) (-2.40000,2.12321) (-2.36667,2.13684) (-2.33333,2.15062) (-2.29494,2.16667) (-2.25552,2.18333) (-2.21667,2.19994) (-2.18333,2.21435) (-2.15000,2.22889) (-2.11667,2.24357) (-2.08333,2.25839) (-2.05000,2.27334) (-2.01667,2.28842) (-1.98333,2.30363) (-1.95000,2.31896) (-1.91667,2.33442) (-1.88329,2.35000) (-1.84786,2.36667) (-1.81267,2.38333) (-1.77769,2.40000) (-1.74291,2.41667) (-1.70830,2.43333) (-1.67384,2.45000) (-1.63950,2.46667) (-1.60526,2.48333) (-1.57109,2.50000) (-1.53696,2.51667) (-1.50285,2.53333) (-1.46872,2.55000) (-1.43453,2.56667) (-1.40025,2.58333) (-1.36667,2.59960) (-1.33333,2.61568) (-1.30000,2.63166) (-1.26667,2.64752) (-1.23333,2.66326) (-1.20000,2.67885) (-1.16667,2.69426) (-1.13333,2.70949) (-1.10000,2.72451) (-1.06667,2.73930) (-1.03333,2.75385) (-1.00000,2.76813) (-0.96374,2.78333) (-0.92295,2.80000) (-0.88333,2.81571) (-0.85000,2.82856) (-0.81667,2.84103) (-0.78333,2.85313) (-0.74464,2.86667) (-0.70000,2.88159) (-0.66667,2.89224) (-0.63333,2.90243) (-0.58402,2.91667) (-0.55000,2.92590) (-0.51667,2.93444) (-0.46667,2.94634) (-0.43333,2.95363) (-0.38333,2.96362) (-0.35000,2.96962) (-0.30000,2.97764) (-0.25863,2.98333) (-0.21667,2.98831) (-0.16667,2.99308) (-0.11667,2.99661) (-0.06667,2.99889) (-0.01667,2.99993) (0.00000,3.00000) (0.05000,2.99938) (0.10000,2.99751) (0.15000,2.99440) (0.20000,2.99004) (0.25000,2.98444) (0.28333,2.98004) (0.33333,2.97243) (0.36673,2.96667) (0.41667,2.95710) (0.45019,2.95000) (0.50000,2.93853) (0.53333,2.93023) (0.58333,2.91686) (0.61667,2.90736) (0.65000,2.89739) (0.69463,2.88333) (0.73333,2.87052) (0.76667,2.85903) (0.80000,2.84713) (0.83736,2.83333) (0.88089,2.81667) (0.91667,2.80252) (0.95000,2.78899) (0.98333,2.77516) (1.01667,2.76102) (1.05000,2.74661) (1.08333,2.73194) (1.11749,2.71667) (1.15418,2.70000) (1.19037,2.68333) (1.22610,2.66667) (1.26146,2.65000) (1.29651,2.63333) (1.33129,2.61667) (1.36585,2.60000) (1.40000,2.58345) (1.43333,2.56724) (1.46667,2.55099) (1.50000,2.53472) (1.53333,2.51843) (1.56667,2.50215) (1.60000,2.48588) (1.63333,2.46965) (1.66667,2.45346) (1.70000,2.43732) (1.73333,2.42125) (1.76667,2.40526) (1.80000,2.38934) (1.83333,2.37352) (1.86667,2.35780) (1.90000,2.34218) (1.93333,2.32667) (1.96667,2.31128) (2.00000,2.29601) (2.03333,2.28087) (2.06667,2.26586) (2.10219,2.25000) (2.13991,2.23333) (2.17802,2.21667) (2.21654,2.20000) (2.25000,2.18567) (2.28333,2.17154) (2.31667,2.15756) (2.35000,2.14371) (2.38333,2.13001) (2.41667,2.11645) (2.45758,2.10000) (2.49955,2.08333) (2.53333,2.07008) (2.56667,2.05714) (2.60000,2.04434) (2.63333,2.03168) (2.67331,2.01667) (2.71667,2.00059) (2.75000,1.98839) (2.78333,1.97631) (2.81667,1.96436) (2.85723,1.95000) (2.90000,1.93504) (2.93333,1.92352) (2.96667,1.91212) (3.00251,1.90000) (3.05000,1.88414) (3.08333,1.87314) (3.11667,1.86226) (3.15464,1.85000) (3.20000,1.83553) (3.23333,1.82502) (3.26667,1.81462) (3.31409,1.80000) (3.35000,1.78906) (3.38333,1.77901) (3.42472,1.76667) (3.46667,1.75430) (3.50000,1.74458) (3.53898,1.73333) (3.58333,1.72068) (3.61667,1.71127) (3.65702,1.70000) (3.70000,1.68813) (3.73333,1.67901) (3.77900,1.66667) (3.81667,1.65659) (3.85000,1.64776) (3.90000,1.63465) (3.93333,1.62601) (3.96971,1.61667) (4.01667,1.60473) (4.05000,1.59634) (4.10000,1.58389) (4.13333,1.57567) (4.17020,1.56667) (4.21667,1.55543) (4.25000,1.54744) (4.30000,1.53559) (4.33333,1.52776) (4.38104,1.51667) (4.41667,1.50846) (4.45372,1.50000) (4.50000,1.48953) (4.53333,1.48206) (4.58333,1.47095) (4.61667,1.46362) (4.66667,1.45272) (4.70000,1.44553) (4.75000,1.43483) (4.78333,1.42776) (4.83333,1.41725) (4.86667,1.41030) (4.91656,1.40000) (4.95000,1.39315) (4.99838,1.38333) (5.00000,1.38301)};
\draw[main,line width=.35pt] plot coordinates {(-5.00000,1.89032) (-4.95203,1.90000) (-4.91667,1.90719) (-4.87046,1.91667) (-4.83333,1.92434) (-4.79020,1.93333) (-4.75000,1.94178) (-4.71121,1.95000) (-4.66667,1.95952) (-4.63333,1.96670) (-4.58333,1.97756) (-4.55000,1.98487) (-4.50000,1.99592) (-4.46667,2.00336) (-4.41667,2.01461) (-4.38333,2.02218) (-4.33464,2.03333) (-4.30000,2.04134) (-4.26283,2.05000) (-4.21667,2.06085) (-4.18333,2.06876) (-4.13333,2.08072) (-4.10000,2.08878) (-4.05397,2.10000) (-4.01667,2.10918) (-3.98333,2.11745) (-3.93333,2.12997) (-3.90000,2.13839) (-3.85452,2.15000) (-3.81667,2.15975) (-3.78333,2.16840) (-3.73333,2.18152) (-3.70000,2.19034) (-3.66385,2.20000) (-3.61667,2.21272) (-3.58333,2.22179) (-3.54133,2.23333) (-3.50000,2.24479) (-3.46667,2.25412) (-3.42229,2.26667) (-3.38333,2.27778) (-3.35000,2.28737) (-3.30655,2.30000) (-3.26667,2.31170) (-3.23333,2.32156) (-3.19393,2.33333) (-3.15000,2.34658) (-3.11667,2.35672) (-3.08333,2.36696) (-3.03333,2.38245) (-3.00000,2.39288) (-2.96667,2.40340) (-2.92499,2.41667) (-2.88333,2.43005) (-2.85000,2.44085) (-2.81667,2.45175) (-2.77141,2.46667) (-2.73333,2.47933) (-2.70000,2.49051) (-2.66667,2.50177) (-2.62296,2.51667) (-2.58333,2.53028) (-2.55000,2.54182) (-2.51667,2.55345) (-2.47904,2.56667) (-2.43333,2.58284) (-2.40000,2.59473) (-2.36667,2.60669) (-2.33333,2.61872) (-2.29313,2.63333) (-2.25000,2.64910) (-2.21667,2.66135) (-2.18333,2.67367) (-2.15000,2.68605) (-2.11260,2.70000) (-2.06811,2.71667) (-2.03333,2.72974) (-2.00000,2.74232) (-1.96667,2.75493) (-1.93333,2.76757) (-1.89185,2.78333) (-1.85000,2.79925) (-1.81667,2.81193) (-1.78333,2.82462) (-1.75000,2.83731) (-1.71662,2.85000) (-1.67270,2.86667) (-1.63333,2.88155) (-1.60000,2.89411) (-1.56667,2.90662) (-1.53333,2.91908) (-1.49496,2.93333) (-1.45000,2.94989) (-1.41667,2.96207) (-1.38333,2.97414) (-1.35000,2.98610) (-1.31084,3.00000) (-1.26667,3.01545) (-1.23333,3.02694) (-1.20000,3.03827) (-1.16495,3.05000) (-1.11667,3.06583) (-1.08333,3.07653) (-1.05000,3.08701) (-1.00774,3.10000) (-0.96667,3.11227) (-0.93333,3.12196) (-0.89306,3.13333) (-0.85000,3.14508) (-0.81667,3.15385) (-0.76667,3.16648) (-0.73333,3.17454) (-0.69540,3.18333) (-0.65000,3.19334) (-0.61667,3.20031) (-0.56667,3.21015) (-0.53111,3.21667) (-0.48333,3.22484) (-0.43333,3.23261) (-0.40000,3.23735) (-0.35000,3.24377) (-0.30000,3.24935) (-0.26667,3.25261) (-0.21667,3.25680) (-0.16667,3.26013) (-0.11667,3.26260) (-0.06667,3.26419) (-0.01667,3.26492) (0.03333,3.26477) (0.08333,3.26376) (0.13333,3.26187) (0.18333,3.25912) (0.23333,3.25550) (0.28333,3.25103) (0.31667,3.24758) (0.36667,3.24172) (0.41667,3.23502) (0.45000,3.23011) (0.50000,3.22207) (0.53333,3.21627) (0.58333,3.20695) (0.61815,3.20000) (0.66667,3.18974) (0.70000,3.18229) (0.75000,3.17055) (0.78333,3.16235) (0.83141,3.15000) (0.86667,3.14059) (0.90000,3.13140) (0.95000,3.11715) (0.98333,3.10734) (1.01667,3.09729) (1.06177,3.08333) (1.10000,3.07120) (1.13333,3.06041) (1.16667,3.04943) (1.21458,3.03333) (1.25000,3.02121) (1.28333,3.00965) (1.31667,2.99794) (1.35775,2.98333) (1.40000,2.96812) (1.43333,2.95599) (1.46667,2.94377) (1.50000,2.93146) (1.53981,2.91667) (1.58333,2.90038) (1.61667,2.88783) (1.65000,2.87525) (1.68333,2.86263) (1.71667,2.84998) (1.76047,2.83333) (1.80000,2.81828) (1.83333,2.80559) (1.86667,2.79290) (1.90000,2.78023) (1.93572,2.76667) (1.97970,2.75000) (2.01667,2.73603) (2.05000,2.72347) (2.08333,2.71095) (2.11667,2.69847) (2.15732,2.68333) (2.20000,2.66751) (2.23333,2.65522) (2.26667,2.64299) (2.30000,2.63082) (2.33904,2.61667) (2.38333,2.60070) (2.41667,2.58878) (2.45000,2.57692) (2.48333,2.56515) (2.52655,2.55000) (2.56667,2.53604) (2.60000,2.52454) (2.63333,2.51311) (2.67191,2.50000) (2.71667,2.48491) (2.75000,2.47377) (2.78333,2.46272) (2.82201,2.45000) (2.86667,2.43544) (2.90000,2.42468) (2.93333,2.41400) (2.97743,2.40000) (3.01667,2.38766) (3.05000,2.37726) (3.08428,2.36667) (3.13333,2.35164) (3.16667,2.34153) (3.20000,2.33151) (3.24985,2.31667) (3.28333,2.30679) (3.31667,2.29704) (3.36401,2.28333) (3.40000,2.27301) (3.43333,2.26353) (3.48138,2.25000) (3.51667,2.24016) (3.55000,2.23094) (3.60000,2.21725) (3.63333,2.20821) (3.66667,2.19924) (3.71667,2.18592) (3.75000,2.17713) (3.79001,2.16667) (3.83333,2.15544) (3.86667,2.14688) (3.91667,2.13417) (3.95000,2.12578) (3.98647,2.11667) (4.03333,2.10507) (4.06667,2.09689) (4.11667,2.08474) (4.15000,2.07672) (4.19213,2.06667) (4.23333,2.05692) (4.26667,2.04910) (4.31667,2.03748) (4.35000,2.02980) (4.40000,2.01839) (4.43333,2.01084) (4.48169,2.00000) (4.51667,1.99222) (4.55699,1.98333) (4.60000,1.97393) (4.63348,1.96667) (4.68333,1.95594) (4.71667,1.94884) (4.76667,1.93827) (4.80000,1.93128) (4.85000,1.92089) (4.88333,1.91402) (4.93333,1.90379) (4.96667,1.89704) (5.00000,1.89032)};
\draw[main,line width=.35pt] plot coordinates {(-5.00000,2.38757) (-4.95000,2.39740) (-4.91667,2.40400) (-4.86667,2.41398) (-4.83333,2.42069) (-4.78333,2.43082) (-4.75000,2.43763) (-4.70000,2.44792) (-4.66667,2.45483) (-4.61667,2.46527) (-4.58333,2.47229) (-4.53333,2.48290) (-4.50000,2.49003) (-4.45372,2.50000) (-4.41667,2.50804) (-4.37723,2.51667) (-4.33333,2.52634) (-4.30000,2.53374) (-4.25000,2.54493) (-4.21667,2.55244) (-4.16667,2.56381) (-4.13333,2.57145) (-4.08333,2.58300) (-4.05000,2.59075) (-4.01057,2.60000) (-3.96667,2.61037) (-3.93333,2.61831) (-3.88333,2.63030) (-3.85000,2.63836) (-3.80227,2.65000) (-3.76667,2.65874) (-3.73333,2.66699) (-3.68333,2.67945) (-3.65000,2.68782) (-3.60191,2.70000) (-3.56667,2.70899) (-3.53333,2.71755) (-3.48333,2.73049) (-3.45000,2.73918) (-3.40882,2.75000) (-3.36667,2.76115) (-3.33333,2.77003) (-3.28378,2.78333) (-3.25000,2.79246) (-3.21667,2.80153) (-3.16667,2.81523) (-3.13333,2.82443) (-3.10000,2.83368) (-3.05000,2.84766) (-3.01667,2.85703) (-2.98263,2.86667) (-2.93333,2.88070) (-2.90000,2.89026) (-2.86619,2.90000) (-2.81667,2.91435) (-2.78333,2.92407) (-2.75000,2.93383) (-2.70000,2.94856) (-2.66667,2.95842) (-2.63333,2.96834) (-2.58333,2.98327) (-2.55000,2.99327) (-2.51667,3.00331) (-2.47250,3.01667) (-2.43333,3.02854) (-2.40000,3.03869) (-2.36293,3.05000) (-2.31667,3.06415) (-2.28333,3.07436) (-2.25000,3.08460) (-2.20000,3.09998) (-2.16667,3.11023) (-2.13333,3.12050) (-2.09167,3.13333) (-2.05000,3.14616) (-2.01667,3.15641) (-1.98327,3.16667) (-1.93333,3.18197) (-1.90000,3.19216) (-1.86667,3.20233) (-1.81945,3.21667) (-1.78333,3.22758) (-1.75000,3.23760) (-1.70857,3.25000) (-1.66667,3.26244) (-1.63333,3.27227) (-1.59553,3.28333) (-1.55000,3.29652) (-1.51667,3.30609) (-1.47941,3.31667) (-1.43333,3.32958) (-1.40000,3.33880) (-1.35897,3.35000) (-1.31667,3.36136) (-1.28333,3.37016) (-1.23333,3.38312) (-1.20000,3.39158) (-1.16624,3.40000) (-1.11667,3.41208) (-1.08333,3.41999) (-1.03333,3.43155) (-1.00000,3.43903) (-0.95000,3.44991) (-0.91667,3.45692) (-0.86865,3.46667) (-0.83333,3.47357) (-0.78333,3.48293) (-0.75000,3.48890) (-0.70000,3.49743) (-0.66667,3.50283) (-0.61667,3.51049) (-0.57344,3.51667) (-0.53333,3.52205) (-0.48333,3.52823) (-0.43796,3.53333) (-0.40000,3.53725) (-0.35000,3.54189) (-0.30000,3.54591) (-0.25000,3.54933) (-0.21667,3.55127) (-0.16667,3.55367) (-0.11667,3.55545) (-0.06667,3.55660) (-0.01667,3.55712) (0.03333,3.55701) (0.08333,3.55628) (0.13333,3.55492) (0.18333,3.55294) (0.23333,3.55033) (0.26667,3.54826) (0.31667,3.54464) (0.36667,3.54041) (0.41667,3.53558) (0.45000,3.53203) (0.50000,3.52623) (0.55000,3.51986) (0.58333,3.51530) (0.63333,3.50800) (0.68333,3.50016) (0.71667,3.49465) (0.76667,3.48594) (0.80000,3.47987) (0.85000,3.47034) (0.88333,3.46373) (0.93333,3.45344) (0.96667,3.44633) (1.01667,3.43531) (1.05000,3.42774) (1.09741,3.41667) (1.13333,3.40806) (1.16667,3.39989) (1.21667,3.38737) (1.25000,3.37883) (1.29662,3.36667) (1.33333,3.35691) (1.36667,3.34791) (1.41667,3.33421) (1.45000,3.32493) (1.48333,3.31556) (1.53333,3.30132) (1.56667,3.29171) (1.60000,3.28203) (1.65000,3.26736) (1.68333,3.25750) (1.71667,3.24758) (1.76423,3.23333) (1.80000,3.22255) (1.83333,3.21245) (1.87432,3.20000) (1.91667,3.18707) (1.95000,3.17687) (1.98333,3.16665) (2.03333,3.15128) (2.06667,3.14103) (2.10000,3.13076) (2.14579,3.11667) (2.18333,3.10510) (2.21667,3.09485) (2.25413,3.08333) (2.30000,3.06925) (2.33333,3.05904) (2.36667,3.04886) (2.41667,3.03361) (2.45000,3.02348) (2.48333,3.01338) (2.52767,3.00000) (2.56667,2.98827) (2.60000,2.97828) (2.63895,2.96667) (2.68333,2.95349) (2.71667,2.94364) (2.75170,2.93333) (2.80000,2.91920) (2.83333,2.90951) (2.86667,2.89986) (2.91667,2.88547) (2.95000,2.87594) (2.98333,2.86647) (3.03333,2.85234) (3.06667,2.84298) (3.10126,2.83333) (3.15000,2.81983) (3.18333,2.81065) (3.22230,2.80000) (3.26667,2.78795) (3.30000,2.77896) (3.34595,2.76667) (3.38333,2.75673) (3.41667,2.74793) (3.46667,2.73483) (3.50000,2.72616) (3.53677,2.71667) (3.58333,2.70473) (3.61667,2.69625) (3.66667,2.68363) (3.70000,2.67528) (3.73463,2.66667) (3.78333,2.65464) (3.81667,2.64648) (3.86667,2.63433) (3.90000,2.62629) (3.94022,2.61667) (3.98333,2.60642) (4.01667,2.59856) (4.06667,2.58687) (4.10000,2.57913) (4.15000,2.56762) (4.18333,2.56001) (4.22749,2.55000) (4.26667,2.54119) (4.30183,2.53333) (4.35000,2.52266) (4.38333,2.51533) (4.43333,2.50442) (4.46667,2.49720) (4.51667,2.48646) (4.55000,2.47935) (4.60000,2.46878) (4.63333,2.46178) (4.68333,2.45137) (4.71667,2.44448) (4.76667,2.43422) (4.80000,2.42743) (4.85000,2.41733) (4.88333,2.41065) (4.93333,2.40070) (4.96667,2.39411) (5.00000,2.38757)};
\draw[main,line width=.35pt] plot coordinates {(-5.00000,2.87546) (-4.95821,2.88333) (-4.91667,2.89120) (-4.87058,2.90000) (-4.83333,2.90715) (-4.78414,2.91667) (-4.75000,2.92331) (-4.70000,2.93311) (-4.66667,2.93968) (-4.61667,2.94961) (-4.58333,2.95627) (-4.53333,2.96632) (-4.50000,2.97307) (-4.45000,2.98326) (-4.41667,2.99010) (-4.36870,3.00000) (-4.33333,3.00735) (-4.28876,3.01667) (-4.25000,3.02482) (-4.20980,3.03333) (-4.16667,3.04252) (-4.13179,3.05000) (-4.08333,3.06045) (-4.05000,3.06769) (-4.00000,3.07861) (-3.96667,3.08594) (-3.91667,3.09700) (-3.88333,3.10442) (-3.83333,3.11561) (-3.80000,3.12312) (-3.75495,3.13333) (-3.71667,3.14206) (-3.68202,3.15000) (-3.63333,3.16122) (-3.60000,3.16895) (-3.55000,3.18060) (-3.51667,3.18842) (-3.46753,3.20000) (-3.43333,3.20810) (-3.39737,3.21667) (-3.35000,3.22800) (-3.31667,3.23602) (-3.26667,3.24810) (-3.23333,3.25619) (-3.19041,3.26667) (-3.15000,3.27656) (-3.11667,3.28476) (-3.06667,3.29711) (-3.03333,3.30538) (-2.98801,3.31667) (-2.95000,3.32616) (-2.91667,3.33452) (-2.86667,3.34709) (-2.83333,3.35549) (-2.78917,3.36667) (-2.75000,3.37659) (-2.71667,3.38506) (-2.66667,3.39779) (-2.63333,3.40629) (-2.59273,3.41667) (-2.55000,3.42759) (-2.51667,3.43612) (-2.46667,3.44892) (-2.43333,3.45745) (-2.39736,3.46667) (-2.35000,3.47878) (-2.31667,3.48731) (-2.26695,3.50000) (-2.23333,3.50856) (-2.20000,3.51704) (-2.15000,3.52973) (-2.11667,3.53815) (-2.06964,3.55000) (-2.03333,3.55910) (-2.00000,3.56743) (-1.95000,3.57985) (-1.91667,3.58808) (-1.86806,3.60000) (-1.83333,3.60845) (-1.79936,3.61667) (-1.75000,3.62850) (-1.71667,3.63641) (-1.66667,3.64815) (-1.63333,3.65590) (-1.58641,3.66667) (-1.55000,3.67492) (-1.51238,3.68333) (-1.46667,3.69341) (-1.43333,3.70065) (-1.38333,3.71132) (-1.35000,3.71830) (-1.30000,3.72857) (-1.26667,3.73528) (-1.21667,3.74512) (-1.18333,3.75152) (-1.13333,3.76089) (-1.10000,3.76697) (-1.05000,3.77583) (-1.00615,3.78333) (-0.96667,3.78988) (-0.91667,3.79786) (-0.88333,3.80299) (-0.83333,3.81038) (-0.78863,3.81667) (-0.75000,3.82187) (-0.70000,3.82825) (-0.65765,3.83333) (-0.61667,3.83800) (-0.56667,3.84330) (-0.51667,3.84818) (-0.48333,3.85120) (-0.43333,3.85537) (-0.38333,3.85910) (-0.33333,3.86238) (-0.28333,3.86522) (-0.25000,3.86686) (-0.20000,3.86894) (-0.15000,3.87056) (-0.10000,3.87172) (-0.05000,3.87242) (0.00000,3.87265) (0.05000,3.87242) (0.10000,3.87172) (0.15000,3.87056) (0.20000,3.86894) (0.25000,3.86686) (0.28333,3.86522) (0.33333,3.86238) (0.38333,3.85910) (0.43333,3.85537) (0.48333,3.85120) (0.51667,3.84818) (0.56667,3.84330) (0.61667,3.83800) (0.65765,3.83333) (0.70000,3.82825) (0.75000,3.82187) (0.78863,3.81667) (0.83333,3.81038) (0.88333,3.80299) (0.91667,3.79786) (0.96667,3.78988) (1.00615,3.78333) (1.05000,3.77583) (1.10000,3.76697) (1.13333,3.76089) (1.18333,3.75152) (1.21667,3.74512) (1.26667,3.73528) (1.30000,3.72857) (1.35000,3.71830) (1.38333,3.71132) (1.43333,3.70065) (1.46667,3.69341) (1.51238,3.68333) (1.55000,3.67492) (1.58641,3.66667) (1.63333,3.65590) (1.66667,3.64815) (1.71667,3.63641) (1.75000,3.62850) (1.79936,3.61667) (1.83333,3.60845) (1.86806,3.60000) (1.91667,3.58808) (1.95000,3.57985) (2.00000,3.56743) (2.03333,3.55910) (2.06964,3.55000) (2.11667,3.53815) (2.15000,3.52973) (2.20000,3.51704) (2.23333,3.50856) (2.26695,3.50000) (2.31667,3.48731) (2.35000,3.47878) (2.39736,3.46667) (2.43333,3.45745) (2.46667,3.44892) (2.51667,3.43612) (2.55000,3.42759) (2.59273,3.41667) (2.63333,3.40629) (2.66667,3.39779) (2.71667,3.38506) (2.75000,3.37659) (2.78917,3.36667) (2.83333,3.35549) (2.86667,3.34709) (2.91667,3.33452) (2.95000,3.32616) (2.98801,3.31667) (3.03333,3.30538) (3.06667,3.29711) (3.11667,3.28476) (3.15000,3.27656) (3.19041,3.26667) (3.23333,3.25619) (3.26667,3.24810) (3.31667,3.23602) (3.35000,3.22800) (3.39737,3.21667) (3.43333,3.20810) (3.46753,3.20000) (3.51667,3.18842) (3.55000,3.18060) (3.60000,3.16895) (3.63333,3.16122) (3.68202,3.15000) (3.71667,3.14206) (3.75495,3.13333) (3.80000,3.12312) (3.83333,3.11561) (3.88333,3.10442) (3.91667,3.09700) (3.96667,3.08594) (4.00000,3.07861) (4.05000,3.06769) (4.08333,3.06045) (4.13179,3.05000) (4.16667,3.04252) (4.20980,3.03333) (4.25000,3.02482) (4.28876,3.01667) (4.33333,3.00735) (4.36870,3.00000) (4.41667,2.99010) (4.45000,2.98326) (4.50000,2.97307) (4.53333,2.96632) (4.58333,2.95627) (4.61667,2.94961) (4.66667,2.93968) (4.70000,2.93311) (4.75000,2.92331) (4.78414,2.91667) (4.83333,2.90715) (4.87058,2.90000) (4.91667,2.89120) (4.95821,2.88333) (5.00000,2.87546)};
\draw[main,line width=.35pt] plot coordinates {(-5.00000,3.35512) (-4.95000,3.36402) (-4.91667,3.36998) (-4.86667,3.37898) (-4.83333,3.38501) (-4.78333,3.39410) (-4.75000,3.40019) (-4.70000,3.40938) (-4.66062,3.41667) (-4.61667,3.42483) (-4.57121,3.43333) (-4.53333,3.44045) (-4.48333,3.44990) (-4.45000,3.45623) (-4.40000,3.46578) (-4.36667,3.47217) (-4.31667,3.48182) (-4.28333,3.48828) (-4.23333,3.49802) (-4.20000,3.50455) (-4.15000,3.51439) (-4.11667,3.52098) (-4.06667,3.53091) (-4.03333,3.53756) (-3.98333,3.54759) (-3.95000,3.55430) (-3.90000,3.56442) (-3.86667,3.57120) (-3.81667,3.58140) (-3.78333,3.58824) (-3.73333,3.59853) (-3.70000,3.60542) (-3.65000,3.61579) (-3.61667,3.62274) (-3.56667,3.63319) (-3.53333,3.64019) (-3.48673,3.65000) (-3.45000,3.65776) (-3.40796,3.66667) (-3.36667,3.67544) (-3.32964,3.68333) (-3.28333,3.69323) (-3.25000,3.70037) (-3.20000,3.71111) (-3.16667,3.71828) (-3.11667,3.72907) (-3.08333,3.73627) (-3.03333,3.74710) (-3.00000,3.75432) (-2.95000,3.76517) (-2.91667,3.77241) (-2.86667,3.78329) (-2.83333,3.79054) (-2.78984,3.80000) (-2.75000,3.80866) (-2.71321,3.81667) (-2.66667,3.82678) (-2.63333,3.83402) (-2.58333,3.84486) (-2.55000,3.85208) (-2.50000,3.86289) (-2.46667,3.87008) (-2.41667,3.88083) (-2.38333,3.88797) (-2.33333,3.89865) (-2.30000,3.90575) (-2.25000,3.91634) (-2.21667,3.92337) (-2.16914,3.93333) (-2.13333,3.94080) (-2.08894,3.95000) (-2.05000,3.95801) (-2.00765,3.96667) (-1.96667,3.97497) (-1.92507,3.98333) (-1.88333,3.99164) (-1.84094,4.00000) (-1.80000,4.00798) (-1.75495,4.01667) (-1.71667,4.02396) (-1.66675,4.03333) (-1.63333,4.03953) (-1.58333,4.04865) (-1.55000,4.05465) (-1.50000,4.06349) (-1.46667,4.06928) (-1.41667,4.07781) (-1.38333,4.08338) (-1.33333,4.09158) (-1.28333,4.09955) (-1.25000,4.10475) (-1.20000,4.11235) (-1.16667,4.11729) (-1.11667,4.12450) (-1.06667,4.13145) (-1.03333,4.13594) (-0.98333,4.14247) (-0.93333,4.14871) (-0.90000,4.15272) (-0.85000,4.15851) (-0.80000,4.16400) (-0.76667,4.16749) (-0.71667,4.17248) (-0.66667,4.17717) (-0.61667,4.18154) (-0.58333,4.18427) (-0.53333,4.18811) (-0.48333,4.19162) (-0.43333,4.19480) (-0.38333,4.19764) (-0.33615,4.20000) (-0.30000,4.20161) (-0.25000,4.20354) (-0.20000,4.20513) (-0.15000,4.20636) (-0.10000,4.20724) (-0.05000,4.20777) (0.00000,4.20794) (0.05000,4.20777) (0.10000,4.20724) (0.15000,4.20636) (0.20000,4.20513) (0.25000,4.20354) (0.30000,4.20161) (0.33615,4.20000) (0.38333,4.19764) (0.43333,4.19480) (0.48333,4.19162) (0.53333,4.18811) (0.58333,4.18427) (0.61667,4.18154) (0.66667,4.17717) (0.71667,4.17248) (0.76667,4.16749) (0.80000,4.16400) (0.85000,4.15851) (0.90000,4.15272) (0.93333,4.14871) (0.98333,4.14247) (1.03333,4.13594) (1.06667,4.13145) (1.11667,4.12450) (1.16667,4.11729) (1.20000,4.11235) (1.25000,4.10475) (1.28333,4.09955) (1.33333,4.09158) (1.38333,4.08338) (1.41667,4.07781) (1.46667,4.06928) (1.50000,4.06349) (1.55000,4.05465) (1.58333,4.04865) (1.63333,4.03953) (1.66675,4.03333) (1.71667,4.02396) (1.75495,4.01667) (1.80000,4.00798) (1.84094,4.00000) (1.88333,3.99164) (1.92507,3.98333) (1.96667,3.97497) (2.00765,3.96667) (2.05000,3.95801) (2.08894,3.95000) (2.13333,3.94080) (2.16914,3.93333) (2.21667,3.92337) (2.25000,3.91634) (2.30000,3.90575) (2.33333,3.89865) (2.38333,3.88797) (2.41667,3.88083) (2.46667,3.87008) (2.50000,3.86289) (2.55000,3.85208) (2.58333,3.84486) (2.63333,3.83402) (2.66667,3.82678) (2.71321,3.81667) (2.75000,3.80866) (2.78984,3.80000) (2.83333,3.79054) (2.86667,3.78329) (2.91667,3.77241) (2.95000,3.76517) (3.00000,3.75432) (3.03333,3.74710) (3.08333,3.73627) (3.11667,3.72907) (3.16667,3.71828) (3.20000,3.71111) (3.25000,3.70037) (3.28333,3.69323) (3.32964,3.68333) (3.36667,3.67544) (3.40796,3.66667) (3.45000,3.65776) (3.48673,3.65000) (3.53333,3.64019) (3.56667,3.63319) (3.61667,3.62274) (3.65000,3.61579) (3.70000,3.60542) (3.73333,3.59853) (3.78333,3.58824) (3.81667,3.58140) (3.86667,3.57120) (3.90000,3.56442) (3.95000,3.55430) (3.98333,3.54759) (4.03333,3.53756) (4.06667,3.53091) (4.11667,3.52098) (4.15000,3.51439) (4.20000,3.50455) (4.23333,3.49802) (4.28333,3.48828) (4.31667,3.48182) (4.36667,3.47217) (4.40000,3.46578) (4.45000,3.45623) (4.48333,3.44990) (4.53333,3.44045) (4.57121,3.43333) (4.61667,3.42483) (4.66062,3.41667) (4.70000,3.40938) (4.75000,3.40019) (4.78333,3.39410) (4.83333,3.38501) (4.86667,3.37898) (4.91667,3.36998) (4.95000,3.36402) (5.00000,3.35512)};
\draw[black,->](-3,-2)--(-2.52482,-1.93212);
\draw[black,->](3,-2)--(3.47518,-2.06788);
\draw[black,->](-3,0)--(-2.53691,0.12630);
\draw[black,->](3,0)--(3.46309,-0.12630);
\draw[black,->](-3,3)--(-2.53691,3.12630);
\draw[black,->](3,3)--(3.46309,2.87370);
\draw[black,->](-2,4)--(-1.52932,4.09414);
\draw[black,->](2,4)--(2.47068,3.90586);
\draw[black,->](0,-1.7)--(0.48000,-1.70000);
\end{scope}
\draw[structurecolor](0,2) circle (2pt) node[above right,fill=white,inner sep=1pt]{$2\ii$};
\fill (0,1) circle (2pt) node[left]{$\ii$};
\node at (0,-3.8){$(a)$ 均匀流与涡};
\end{tikzpicture}
\begin{tikzpicture}[x=.49cm,y=.49cm,>=Stealth]
\draw[->](-5.2,0)--(5.5,0) node[right]{$x$};\draw[->](0,-3.2)--(0,5.9) node[above]{$y$};
\begin{scope}\clip (-5,-3) rectangle (5,5.5);
\draw[main,line width=.35pt] plot coordinates {(-5.00000,-2.78173) (-4.95000,-2.77660) (-4.90000,-2.77147) (-4.85320,-2.76667) (-4.81667,-2.76292) (-4.76667,-2.75780) (-4.71667,-2.75268) (-4.68333,-2.74927) (-4.63333,-2.74416) (-4.58333,-2.73906) (-4.53333,-2.73397) (-4.50000,-2.73058) (-4.45000,-2.72550) (-4.40000,-2.72044) (-4.36263,-2.71667) (-4.31667,-2.71203) (-4.26667,-2.70701) (-4.21667,-2.70200) (-4.18333,-2.69867) (-4.13333,-2.69369) (-4.08333,-2.68874) (-4.03333,-2.68380) (-4.00000,-2.68053) (-3.95000,-2.67563) (-3.90000,-2.67077) (-3.85758,-2.66667) (-3.81667,-2.66273) (-3.76667,-2.65795) (-3.71667,-2.65320) (-3.68276,-2.65000) (-3.63333,-2.64537) (-3.58333,-2.64072) (-3.53333,-2.63612) (-3.50000,-2.63308) (-3.45000,-2.62856) (-3.40000,-2.62409) (-3.35000,-2.61967) (-3.31570,-2.61667) (-3.26667,-2.61242) (-3.21667,-2.60816) (-3.16667,-2.60396) (-3.11886,-2.60000) (-3.08333,-2.59710) (-3.03333,-2.59308) (-2.98333,-2.58914) (-2.93333,-2.58527) (-2.90000,-2.58274) (-2.85000,-2.57901) (-2.80000,-2.57537) (-2.75000,-2.57182) (-2.70000,-2.56837) (-2.66667,-2.56612) (-2.61667,-2.56282) (-2.56667,-2.55964) (-2.51667,-2.55656) (-2.46667,-2.55359) (-2.41667,-2.55073) (-2.38333,-2.54889) (-2.33333,-2.54623) (-2.28333,-2.54369) (-2.23333,-2.54127) (-2.18333,-2.53899) (-2.13333,-2.53683) (-2.08333,-2.53481) (-2.04464,-2.53333) (-2.00000,-2.53173) (-1.95000,-2.53006) (-1.90000,-2.52854) (-1.85000,-2.52715) (-1.80000,-2.52589) (-1.75000,-2.52478) (-1.70000,-2.52380) (-1.65000,-2.52296) (-1.60000,-2.52225) (-1.55000,-2.52167) (-1.50000,-2.52122) (-1.45000,-2.52089) (-1.40000,-2.52068) (-1.35000,-2.52059) (-1.30000,-2.52061) (-1.25000,-2.52072) (-1.20000,-2.52094) (-1.15000,-2.52124) (-1.10000,-2.52163) (-1.05000,-2.52208) (-1.00000,-2.52260) (-0.95000,-2.52318) (-0.90000,-2.52380) (-0.85000,-2.52446) (-0.80000,-2.52515) (-0.75000,-2.52585) (-0.70000,-2.52657) (-0.65000,-2.52728) (-0.60000,-2.52798) (-0.55000,-2.52866) (-0.50000,-2.52931) (-0.45000,-2.52993) (-0.40000,-2.53051) (-0.35000,-2.53103) (-0.30000,-2.53150) (-0.25000,-2.53190) (-0.20000,-2.53224) (-0.15000,-2.53251) (-0.10000,-2.53270) (-0.05000,-2.53282) (0.00000,-2.53286) (0.05000,-2.53282) (0.10000,-2.53270) (0.15000,-2.53251) (0.20000,-2.53224) (0.25000,-2.53190) (0.30000,-2.53150) (0.35000,-2.53103) (0.40000,-2.53051) (0.45000,-2.52993) (0.50000,-2.52931) (0.55000,-2.52866) (0.60000,-2.52798) (0.65000,-2.52728) (0.70000,-2.52657) (0.75000,-2.52585) (0.80000,-2.52515) (0.85000,-2.52446) (0.90000,-2.52380) (0.95000,-2.52318) (1.00000,-2.52260) (1.05000,-2.52208) (1.10000,-2.52163) (1.15000,-2.52124) (1.20000,-2.52094) (1.25000,-2.52072) (1.30000,-2.52061) (1.35000,-2.52059) (1.40000,-2.52068) (1.45000,-2.52089) (1.50000,-2.52122) (1.55000,-2.52167) (1.60000,-2.52225) (1.65000,-2.52296) (1.70000,-2.52380) (1.75000,-2.52478) (1.80000,-2.52589) (1.85000,-2.52715) (1.90000,-2.52854) (1.95000,-2.53006) (2.00000,-2.53173) (2.04464,-2.53333) (2.08333,-2.53481) (2.13333,-2.53683) (2.18333,-2.53899) (2.23333,-2.54127) (2.28333,-2.54369) (2.33333,-2.54623) (2.38333,-2.54889) (2.41667,-2.55073) (2.46667,-2.55359) (2.51667,-2.55656) (2.56667,-2.55964) (2.61667,-2.56282) (2.66667,-2.56612) (2.70000,-2.56837) (2.75000,-2.57182) (2.80000,-2.57537) (2.85000,-2.57901) (2.90000,-2.58274) (2.93333,-2.58527) (2.98333,-2.58914) (3.03333,-2.59308) (3.08333,-2.59710) (3.11886,-2.60000) (3.16667,-2.60396) (3.21667,-2.60816) (3.26667,-2.61242) (3.31570,-2.61667) (3.35000,-2.61967) (3.40000,-2.62409) (3.45000,-2.62856) (3.50000,-2.63308) (3.53333,-2.63612) (3.58333,-2.64072) (3.63333,-2.64537) (3.68276,-2.65000) (3.71667,-2.65320) (3.76667,-2.65795) (3.81667,-2.66273) (3.85758,-2.66667) (3.90000,-2.67077) (3.95000,-2.67563) (4.00000,-2.68053) (4.03333,-2.68380) (4.08333,-2.68874) (4.13333,-2.69369) (4.18333,-2.69867) (4.21667,-2.70200) (4.26667,-2.70701) (4.31667,-2.71203) (4.36263,-2.71667) (4.40000,-2.72044) (4.45000,-2.72550) (4.50000,-2.73058) (4.53333,-2.73397) (4.58333,-2.73906) (4.63333,-2.74416) (4.68333,-2.74927) (4.71667,-2.75268) (4.76667,-2.75780) (4.81667,-2.76292) (4.85320,-2.76667) (4.90000,-2.77147) (4.95000,-2.77660) (5.00000,-2.78173)};
\draw[main,line width=.35pt] plot coordinates {(0.06667,2.02758) (0.07460,1.98333) (0.06104,1.95000) (0.03333,1.92595) (-0.01667,1.91999) (-0.05000,1.93759) (-0.07007,1.96667) (-0.07177,2.01667) (-0.06667,2.02758)};
\draw[main,line width=.35pt] plot coordinates {(-5.00000,-2.21971) (-4.96667,-2.21578) (-4.91667,-2.20988) (-4.86667,-2.20396) (-4.83324,-2.20000) (-4.78333,-2.19408) (-4.73333,-2.18814) (-4.69293,-2.18333) (-4.65000,-2.17822) (-4.60000,-2.17226) (-4.55307,-2.16667) (-4.51667,-2.16232) (-4.46667,-2.15635) (-4.41667,-2.15038) (-4.38333,-2.14640) (-4.33333,-2.14043) (-4.28333,-2.13446) (-4.25000,-2.13048) (-4.20000,-2.12451) (-4.15000,-2.11855) (-4.11667,-2.11458) (-4.06667,-2.10864) (-4.01667,-2.10270) (-3.98333,-2.09875) (-3.93333,-2.09283) (-3.88333,-2.08693) (-3.85000,-2.08300) (-3.80000,-2.07713) (-3.75000,-2.07128) (-3.71042,-2.06667) (-3.66667,-2.06158) (-3.61667,-2.05580) (-3.56667,-2.05005) (-3.53333,-2.04624) (-3.48333,-2.04055) (-3.43333,-2.03490) (-3.40000,-2.03116) (-3.35000,-2.02558) (-3.30000,-2.02006) (-3.26667,-2.01640) (-3.21667,-2.01097) (-3.16667,-2.00560) (-3.11667,-2.00029) (-3.08333,-1.99680) (-3.03333,-1.99161) (-2.98333,-1.98651) (-2.95000,-1.98315) (-2.90000,-1.97820) (-2.85000,-1.97333) (-2.80000,-1.96858) (-2.76667,-1.96546) (-2.71667,-1.96089) (-2.66667,-1.95643) (-2.61667,-1.95211) (-2.58333,-1.94930) (-2.53333,-1.94520) (-2.48333,-1.94125) (-2.43333,-1.93746) (-2.38333,-1.93384) (-2.35000,-1.93152) (-2.30000,-1.92819) (-2.25000,-1.92505) (-2.20000,-1.92210) (-2.15000,-1.91937) (-2.10000,-1.91685) (-2.06667,-1.91529) (-2.01667,-1.91314) (-1.96667,-1.91124) (-1.91667,-1.90958) (-1.86667,-1.90817) (-1.81667,-1.90702) (-1.76667,-1.90614) (-1.71667,-1.90553) (-1.66667,-1.90518) (-1.61667,-1.90512) (-1.56667,-1.90532) (-1.51667,-1.90580) (-1.46667,-1.90655) (-1.41667,-1.90756) (-1.36667,-1.90883) (-1.31667,-1.91035) (-1.26667,-1.91210) (-1.21667,-1.91407) (-1.16667,-1.91625) (-1.13333,-1.91781) (-1.08333,-1.92028) (-1.03333,-1.92290) (-0.98333,-1.92565) (-0.93333,-1.92850) (-0.88333,-1.93143) (-0.85000,-1.93341) (-0.80000,-1.93639) (-0.75000,-1.93937) (-0.70000,-1.94233) (-0.65000,-1.94522) (-0.60000,-1.94803) (-0.56365,-1.95000) (-0.51667,-1.95244) (-0.46667,-1.95489) (-0.41667,-1.95716) (-0.36667,-1.95923) (-0.31667,-1.96108) (-0.26667,-1.96269) (-0.21667,-1.96404) (-0.16667,-1.96513) (-0.11667,-1.96594) (-0.06667,-1.96647) (-0.03333,-1.96667) (0.01667,-1.96672) (0.05000,-1.96659) (0.10000,-1.96615) (0.15000,-1.96543) (0.20000,-1.96443) (0.25000,-1.96317) (0.30000,-1.96164) (0.35000,-1.95987) (0.40000,-1.95788) (0.45000,-1.95567) (0.50000,-1.95328) (0.55000,-1.95072) (0.58333,-1.94894) (0.63333,-1.94617) (0.68333,-1.94330) (0.73333,-1.94036) (0.78333,-1.93739) (0.83333,-1.93440) (0.86667,-1.93242) (0.91667,-1.92947) (0.96667,-1.92659) (1.01667,-1.92381) (1.06667,-1.92114) (1.11667,-1.91861) (1.15769,-1.91667) (1.20000,-1.91478) (1.25000,-1.91274) (1.30000,-1.91091) (1.35000,-1.90931) (1.40000,-1.90796) (1.45000,-1.90686) (1.50000,-1.90602) (1.55000,-1.90545) (1.60000,-1.90515) (1.65000,-1.90513) (1.70000,-1.90538) (1.75000,-1.90591) (1.80000,-1.90670) (1.85000,-1.90776) (1.90000,-1.90908) (1.95000,-1.91066) (2.00000,-1.91248) (2.05000,-1.91455) (2.09620,-1.91667) (2.13333,-1.91850) (2.18333,-1.92117) (2.23333,-1.92404) (2.28333,-1.92712) (2.33333,-1.93039) (2.37611,-1.93333) (2.41667,-1.93624) (2.46667,-1.93997) (2.51667,-1.94387) (2.56667,-1.94792) (2.60000,-1.95070) (2.65000,-1.95498) (2.70000,-1.95939) (2.75000,-1.96393) (2.78333,-1.96701) (2.83333,-1.97174) (2.88333,-1.97656) (2.93333,-1.98149) (2.96667,-1.98483) (3.01667,-1.98990) (3.06667,-1.99506) (3.11387,-2.00000) (3.15000,-2.00382) (3.20000,-2.00917) (3.25000,-2.01458) (3.28333,-2.01823) (3.33333,-2.02373) (3.38333,-2.02929) (3.41940,-2.03333) (3.46667,-2.03866) (3.51667,-2.04434) (3.56620,-2.05000) (3.60000,-2.05388) (3.65000,-2.05965) (3.70000,-2.06545) (3.73333,-2.06934) (3.78333,-2.07518) (3.83333,-2.08104) (3.86667,-2.08496) (3.91667,-2.09086) (3.96667,-2.09677) (4.00000,-2.10072) (4.05000,-2.10666) (4.10000,-2.11260) (4.13415,-2.11667) (4.18333,-2.12253) (4.23333,-2.12849) (4.27390,-2.13333) (4.31667,-2.13844) (4.36667,-2.14441) (4.41346,-2.15000) (4.45000,-2.15436) (4.50000,-2.16033) (4.55000,-2.16630) (4.58333,-2.17028) (4.63333,-2.17624) (4.68333,-2.18219) (4.71667,-2.18616) (4.76667,-2.19210) (4.81667,-2.19803) (4.85000,-2.20199) (4.90000,-2.20791) (4.95000,-2.21382) (4.98333,-2.21775) (5.00000,-2.21971)};
\draw[main,line width=.35pt] plot coordinates {(-0.03333,1.86042) (-0.06667,1.87569) (-0.09393,1.90000) (-0.11436,1.93333) (-0.12340,1.96667) (-0.11989,2.01667) (-0.10553,2.05000) (-0.08333,2.07673) (-0.05000,2.09877) (-0.01667,2.10853) (0.03333,2.10497) (0.06667,2.08977) (0.09325,2.06667) (0.11409,2.03333) (0.12334,2.00000) (0.12006,1.95000) (0.10596,1.91667) (0.08333,1.88878) (0.05000,1.86676) (0.01667,1.85688) (-0.03333,1.86042)};
\draw[main,line width=.35pt] plot coordinates {(-5.00000,-1.65703) (-4.95000,-1.65024) (-4.91667,-1.64570) (-4.86667,-1.63887) (-4.82634,-1.63333) (-4.78333,-1.62741) (-4.73333,-1.62050) (-4.70000,-1.61588) (-4.65000,-1.60892) (-4.60000,-1.60194) (-4.56667,-1.59727) (-4.51667,-1.59024) (-4.46774,-1.58333) (-4.43333,-1.57846) (-4.38333,-1.57137) (-4.35000,-1.56662) (-4.30000,-1.55948) (-4.25000,-1.55232) (-4.21667,-1.54753) (-4.16667,-1.54033) (-4.11818,-1.53333) (-4.08333,-1.52829) (-4.03333,-1.52104) (-4.00000,-1.51619) (-3.95000,-1.50891) (-3.90000,-1.50161) (-3.86667,-1.49674) (-3.81667,-1.48942) (-3.77514,-1.48333) (-3.73333,-1.47720) (-3.68333,-1.46985) (-3.65000,-1.46495) (-3.60000,-1.45759) (-3.55000,-1.45024) (-3.51667,-1.44533) (-3.46667,-1.43798) (-3.43333,-1.43308) (-3.38333,-1.42574) (-3.33333,-1.41840) (-3.30000,-1.41352) (-3.25000,-1.40622) (-3.20731,-1.40000) (-3.16667,-1.39409) (-3.11667,-1.38686) (-3.08333,-1.38205) (-3.03333,-1.37488) (-2.98333,-1.36775) (-2.95000,-1.36303) (-2.90000,-1.35599) (-2.85707,-1.35000) (-2.81667,-1.34440) (-2.76667,-1.33756) (-2.73333,-1.33304) (-2.68333,-1.32634) (-2.63333,-1.31975) (-2.60000,-1.31543) (-2.55000,-1.30904) (-2.50000,-1.30279) (-2.46667,-1.29871) (-2.41667,-1.29273) (-2.36667,-1.28693) (-2.33333,-1.28318) (-2.28333,-1.27772) (-2.23333,-1.27250) (-2.18333,-1.26754) (-2.15000,-1.26439) (-2.10000,-1.25991) (-2.05000,-1.25576) (-2.00000,-1.25197) (-1.96667,-1.24966) (-1.91667,-1.24653) (-1.86667,-1.24385) (-1.81667,-1.24165) (-1.76667,-1.23998) (-1.71667,-1.23887) (-1.66667,-1.23836) (-1.61667,-1.23849) (-1.56667,-1.23929) (-1.51667,-1.24081) (-1.46667,-1.24307) (-1.41667,-1.24610) (-1.36667,-1.24992) (-1.33333,-1.25290) (-1.28333,-1.25803) (-1.23333,-1.26394) (-1.20000,-1.26831) (-1.15000,-1.27544) (-1.10000,-1.28326) (-1.06667,-1.28880) (-1.01667,-1.29757) (-0.98333,-1.30368) (-0.93333,-1.31317) (-0.90000,-1.31966) (-0.85000,-1.32956) (-0.81667,-1.33623) (-0.76667,-1.34624) (-0.73333,-1.35287) (-0.68333,-1.36267) (-0.65000,-1.36907) (-0.60000,-1.37838) (-0.56667,-1.38435) (-0.51667,-1.39290) (-0.47239,-1.40000) (-0.43333,-1.40587) (-0.38333,-1.41277) (-0.35000,-1.41697) (-0.30000,-1.42263) (-0.25000,-1.42749) (-0.20000,-1.43150) (-0.16667,-1.43370) (-0.11667,-1.43627) (-0.06667,-1.43794) (-0.01667,-1.43870) (0.03333,-1.43855) (0.08333,-1.43748) (0.13333,-1.43551) (0.17246,-1.43333) (0.21667,-1.43026) (0.26667,-1.42596) (0.31667,-1.42083) (0.35243,-1.41667) (0.40000,-1.41055) (0.45000,-1.40342) (0.48333,-1.39829) (0.53333,-1.39011) (0.57241,-1.38333) (0.61667,-1.37532) (0.66262,-1.36667) (0.70000,-1.35943) (0.74781,-1.35000) (0.78333,-1.34290) (0.83115,-1.33333) (0.86667,-1.32624) (0.91532,-1.31667) (0.95000,-1.30996) (1.00000,-1.30061) (1.03333,-1.29459) (1.08333,-1.28599) (1.11667,-1.28058) (1.16667,-1.27298) (1.21234,-1.26667) (1.25000,-1.26188) (1.30000,-1.25623) (1.35000,-1.25137) (1.38333,-1.24856) (1.43333,-1.24500) (1.48333,-1.24223) (1.53333,-1.24022) (1.58333,-1.23894) (1.63333,-1.23837) (1.68333,-1.23846) (1.73333,-1.23917) (1.78333,-1.24048) (1.83333,-1.24233) (1.88333,-1.24469) (1.93333,-1.24752) (1.97171,-1.25000) (2.01667,-1.25319) (2.06667,-1.25710) (2.11667,-1.26136) (2.16667,-1.26595) (2.20000,-1.26916) (2.25000,-1.27421) (2.30000,-1.27951) (2.33473,-1.28333) (2.38333,-1.28884) (2.43333,-1.29470) (2.47723,-1.30000) (2.51667,-1.30486) (2.56667,-1.31115) (2.60957,-1.31667) (2.65000,-1.32194) (2.70000,-1.32856) (2.73549,-1.33333) (2.78333,-1.33983) (2.83333,-1.34671) (2.86667,-1.35133) (2.91667,-1.35833) (2.96667,-1.36539) (3.00000,-1.37012) (3.05000,-1.37727) (3.09221,-1.38333) (3.13333,-1.38927) (3.18333,-1.39651) (3.21667,-1.40136) (3.26667,-1.40865) (3.31667,-1.41596) (3.35000,-1.42085) (3.40000,-1.42818) (3.43506,-1.43333) (3.48333,-1.44043) (3.53333,-1.44779) (3.56667,-1.45269) (3.61667,-1.46005) (3.66167,-1.46667) (3.70000,-1.47230) (3.75000,-1.47964) (3.78333,-1.48454) (3.83333,-1.49186) (3.88333,-1.49918) (3.91667,-1.50405) (3.96667,-1.51134) (4.00325,-1.51667) (4.05000,-1.52346) (4.10000,-1.53070) (4.13333,-1.53552) (4.18333,-1.54274) (4.23333,-1.54993) (4.26667,-1.55471) (4.31667,-1.56186) (4.35032,-1.56667) (4.40000,-1.57373) (4.45000,-1.58082) (4.48333,-1.58554) (4.53333,-1.59258) (4.58333,-1.59960) (4.61667,-1.60427) (4.66667,-1.61124) (4.70567,-1.61667) (4.75000,-1.62281) (4.80000,-1.62971) (4.83333,-1.63429) (4.88333,-1.64115) (4.93333,-1.64797) (4.96667,-1.65251) (5.00000,-1.65703)};
\draw[main,line width=.35pt] plot coordinates {(-0.05000,1.73300) (-0.08333,1.74522) (-0.11667,1.76529) (-0.13797,1.78333) (-0.16647,1.81667) (-0.18333,1.84513) (-0.19844,1.88333) (-0.20582,1.91667) (-0.20771,1.96667) (-0.20000,2.01069) (-0.18559,2.05000) (-0.16667,2.08178) (-0.15000,2.10238) (-0.11667,2.13180) (-0.08536,2.15000) (-0.05000,2.16289) (-0.01667,2.16851) (0.03080,2.16667) (0.06667,2.15783) (0.10000,2.14249) (0.13333,2.11867) (0.15223,2.10000) (0.17670,2.06667) (0.19292,2.03333) (0.20286,2.00000) (0.20828,1.95000) (0.20275,1.90000) (0.19285,1.86667) (0.17703,1.83333) (0.15364,1.80000) (0.13333,1.77885) (0.10000,1.75415) (0.06667,1.73818) (0.03333,1.72916) (-0.01667,1.72691) (-0.05000,1.73300)};
\draw[main,line width=.35pt] plot coordinates {(-5.00000,-1.09527) (-4.95000,-1.08751) (-4.91667,-1.08231) (-4.86667,-1.07446) (-4.81739,-1.06667) (-4.78333,-1.06125) (-4.73333,-1.05326) (-4.70000,-1.04789) (-4.65000,-1.03980) (-4.61029,-1.03333) (-4.56667,-1.02619) (-4.51667,-1.01794) (-4.48333,-1.01241) (-4.43333,-1.00407) (-4.40000,-0.99847) (-4.35000,-0.99002) (-4.31067,-0.98333) (-4.26667,-0.97580) (-4.21667,-0.96719) (-4.18333,-0.96141) (-4.13333,-0.95269) (-4.10000,-0.94684) (-4.05000,-0.93801) (-4.01667,-0.93209) (-3.96667,-0.92315) (-3.93064,-0.91667) (-3.88333,-0.90810) (-3.83892,-0.90000) (-3.80000,-0.89286) (-3.75000,-0.88362) (-3.71667,-0.87742) (-3.66667,-0.86806) (-3.63333,-0.86178) (-3.58333,-0.85230) (-3.55000,-0.84594) (-3.50000,-0.83634) (-3.46667,-0.82990) (-3.41667,-0.82017) (-3.38333,-0.81365) (-3.33333,-0.80379) (-3.30000,-0.79718) (-3.25000,-0.78719) (-3.21667,-0.78050) (-3.16667,-0.77038) (-3.13333,-0.76359) (-3.08333,-0.75334) (-3.05000,-0.74647) (-3.00000,-0.73608) (-2.96667,-0.72912) (-2.91667,-0.71860) (-2.88333,-0.71154) (-2.83333,-0.70088) (-2.80000,-0.69373) (-2.75186,-0.68333) (-2.71667,-0.67568) (-2.67541,-0.66667) (-2.63333,-0.65741) (-2.59986,-0.65000) (-2.55000,-0.63890) (-2.51667,-0.63143) (-2.46667,-0.62015) (-2.43333,-0.61258) (-2.38333,-0.60117) (-2.35000,-0.59351) (-2.30594,-0.58333) (-2.26667,-0.57420) (-2.23333,-0.56642) (-2.18333,-0.55466) (-2.15000,-0.54679) (-2.10000,-0.53490) (-2.06667,-0.52693) (-2.02390,-0.51667) (-1.98333,-0.50686) (-1.95000,-0.49878) (-1.90000,-0.48658) (-1.86667,-0.47841) (-1.81893,-0.46667) (-1.78333,-0.45785) (-1.75000,-0.44958) (-1.70000,-0.43711) (-1.66667,-0.42877) (-1.61847,-0.41667) (-1.58333,-0.40779) (-1.55000,-0.39937) (-1.50000,-0.38667) (-1.46667,-0.37819) (-1.42138,-0.36667) (-1.38333,-0.35690) (-1.35000,-0.34837) (-1.30000,-0.33553) (-1.26667,-0.32694) (-1.22645,-0.31667) (-1.18333,-0.30546) (-1.15000,-0.29687) (-1.10000,-0.28400) (-1.06667,-0.27526) (-1.03265,-0.26667) (-0.98333,-0.25320)};
\draw[main,line width=.35pt] plot coordinates {(0.98333,-0.25320) (1.03265,-0.26667) (1.06667,-0.27526) (1.10000,-0.28400) (1.15000,-0.29687) (1.18333,-0.30546) (1.22645,-0.31667) (1.26667,-0.32694) (1.30000,-0.33553) (1.35000,-0.34837) (1.38333,-0.35690) (1.42138,-0.36667) (1.46667,-0.37819) (1.50000,-0.38667) (1.55000,-0.39937) (1.58333,-0.40779) (1.61847,-0.41667) (1.66667,-0.42877) (1.70000,-0.43711) (1.75000,-0.44958) (1.78333,-0.45785) (1.81893,-0.46667) (1.86667,-0.47841) (1.90000,-0.48658) (1.95000,-0.49878) (1.98333,-0.50686) (2.02390,-0.51667) (2.06667,-0.52693) (2.10000,-0.53490) (2.15000,-0.54679) (2.18333,-0.55466) (2.23333,-0.56642) (2.26667,-0.57420) (2.30594,-0.58333) (2.35000,-0.59351) (2.38333,-0.60117) (2.43333,-0.61258) (2.46667,-0.62015) (2.51667,-0.63143) (2.55000,-0.63890) (2.59986,-0.65000) (2.63333,-0.65741) (2.67541,-0.66667) (2.71667,-0.67568) (2.75186,-0.68333) (2.80000,-0.69373) (2.83333,-0.70088) (2.88333,-0.71154) (2.91667,-0.71860) (2.96667,-0.72912) (3.00000,-0.73608) (3.05000,-0.74647) (3.08333,-0.75334) (3.13333,-0.76359) (3.16667,-0.77038) (3.21667,-0.78050) (3.25000,-0.78719) (3.30000,-0.79718) (3.33333,-0.80379) (3.38333,-0.81365) (3.41667,-0.82017) (3.46667,-0.82990) (3.50000,-0.83634) (3.55000,-0.84594) (3.58333,-0.85230) (3.63333,-0.86178) (3.66667,-0.86806) (3.71667,-0.87742) (3.75000,-0.88362) (3.80000,-0.89286) (3.83892,-0.90000) (3.88333,-0.90810) (3.93064,-0.91667) (3.96667,-0.92315) (4.01667,-0.93209) (4.05000,-0.93801) (4.10000,-0.94684) (4.13333,-0.95269) (4.18333,-0.96141) (4.21667,-0.96719) (4.26667,-0.97580) (4.31067,-0.98333) (4.35000,-0.99002) (4.40000,-0.99847) (4.43333,-1.00407) (4.48333,-1.01241) (4.51667,-1.01794) (4.56667,-1.02619) (4.61029,-1.03333) (4.65000,-1.03980) (4.70000,-1.04789) (4.73333,-1.05326) (4.78333,-1.06125) (4.81739,-1.06667) (4.86667,-1.07446) (4.91667,-1.08231) (4.95000,-1.08751) (5.00000,-1.09527)};
\draw[main,line width=.35pt] plot coordinates {(-0.00565,1.01667) (-0.02877,1.05000) (-0.05000,1.08262) (-0.06667,1.10843) (-0.08333,1.13522) (-0.10306,1.16667) (-0.12426,1.20000) (-0.14543,1.23333) (-0.16632,1.26667) (-0.18333,1.29437) (-0.20000,1.32253) (-0.21667,1.35173) (-0.23402,1.38333) (-0.25147,1.41667) (-0.26794,1.45000) (-0.28333,1.48334) (-0.30000,1.52265) (-0.31651,1.56667) (-0.32751,1.60000) (-0.33718,1.63333) (-0.34898,1.68333) (-0.35500,1.71667) (-0.36106,1.76667) (-0.36336,1.81667) (-0.36163,1.86667) (-0.35546,1.91667) (-0.34857,1.95000) (-0.33366,2.00000) (-0.32027,2.03333) (-0.30348,2.06667) (-0.28333,2.09893) (-0.26667,2.12136) (-0.24151,2.15000) (-0.21667,2.17312) (-0.18333,2.19830) (-0.15230,2.21667) (-0.11667,2.23279) (-0.08333,2.24365) (-0.05000,2.25060) (0.00000,2.25457) (0.05000,2.25060) (0.08333,2.24365) (0.11667,2.23279) (0.15230,2.21667) (0.18333,2.19830) (0.21667,2.17312) (0.24151,2.15000) (0.26667,2.12136) (0.28333,2.09893) (0.30348,2.06667) (0.32027,2.03333) (0.33366,2.00000) (0.34857,1.95000) (0.35546,1.91667) (0.36163,1.86667) (0.36336,1.81667) (0.36106,1.76667) (0.35500,1.71667) (0.34898,1.68333) (0.33718,1.63333) (0.32751,1.60000) (0.31651,1.56667) (0.30000,1.52265) (0.28333,1.48334) (0.26794,1.45000) (0.25147,1.41667) (0.23402,1.38333) (0.21667,1.35173) (0.20000,1.32253) (0.18333,1.29437) (0.16632,1.26667) (0.14543,1.23333) (0.12426,1.20000) (0.10306,1.16667) (0.08333,1.13522) (0.06667,1.10843) (0.05000,1.08262) (0.02877,1.05000) (0.00565,1.01667)};
\draw[main,line width=.35pt] plot coordinates {(-5.00000,-0.75959) (-4.95000,-0.75123) (-4.91667,-0.74562) (-4.86667,-0.73714) (-4.83333,-0.73145) (-4.78333,-0.72285) (-4.74761,-0.71667) (-4.70000,-0.70836) (-4.65248,-0.70000) (-4.61667,-0.69365) (-4.56667,-0.68473) (-4.53333,-0.67873) (-4.48333,-0.66966) (-4.45000,-0.66357) (-4.40000,-0.65437) (-4.36667,-0.64818) (-4.31667,-0.63883) (-4.28333,-0.63255) (-4.23333,-0.62304) (-4.20000,-0.61666) (-4.15000,-0.60700) (-4.11410,-0.60000) (-4.06667,-0.59068) (-4.02961,-0.58333) (-3.98333,-0.57408) (-3.94657,-0.56667) (-3.90000,-0.55719) (-3.86496,-0.55000) (-3.81667,-0.54000) (-3.78333,-0.53303) (-3.73333,-0.52248) (-3.70000,-0.51538) (-3.65000,-0.50463) (-3.61667,-0.49739) (-3.56667,-0.48643) (-3.53333,-0.47904) (-3.48333,-0.46785) (-3.45000,-0.46031) (-3.40486,-0.45000) (-3.36667,-0.44118) (-3.33299,-0.43333) (-3.28333,-0.42163) (-3.25000,-0.41368) (-3.20000,-0.40162) (-3.16667,-0.39349) (-3.12553,-0.38333) (-3.08333,-0.37279) (-3.05000,-0.36437) (-3.00000,-0.35156) (-2.96667,-0.34291) (-2.93018,-0.33333) (-2.88333,-0.32086) (-2.85000,-0.31186) (-2.80666,-0.30000) (-2.76667,-0.28888) (-2.73333,-0.27949) (-2.68854,-0.26667) (-2.65000,-0.25545) (-2.61667,-0.24560) (-2.57584,-0.23333) (-2.53333,-0.22032) (-2.50000,-0.20994) (-2.46667,-0.19938) (-2.41706,-0.18333) (-2.38333,-0.17218) (-2.35000,-0.16095) (-2.31667,-0.14949) (-2.27076,-0.13333) (-2.23333,-0.11981) (-2.20000,-0.10746) (-2.16667,-0.09483) (-2.13333,-0.08187) (-2.09529,-0.06667) (-2.05491,-0.05000) (-2.01667,-0.03366) (-1.98333,-0.01893) (-1.95000,-0.00371) (-1.91667,0.01205) (-1.88333,0.02841) (-1.85000,0.04540) (-1.81667,0.06311) (-1.78333,0.08161) (-1.75163,0.10000) (-1.72410,0.11667) (-1.69764,0.13333) (-1.66667,0.15373) (-1.63333,0.17685) (-1.60175,0.20000) (-1.58000,0.21667) (-1.55000,0.24069) (-1.51928,0.26667) (-1.50000,0.28367) (-1.46667,0.31441) (-1.44705,0.33333) (-1.41667,0.36376) (-1.39785,0.38333) (-1.36683,0.41667) (-1.35000,0.43527) (-1.32232,0.46667) (-1.30000,0.49248) (-1.27950,0.51667) (-1.25157,0.55000) (-1.23333,0.57195) (-1.21024,0.60000) (-1.18333,0.63269) (-1.16667,0.65296) (-1.14168,0.68333) (-1.11667,0.71354) (-1.10000,0.73361) (-1.07237,0.76667) (-1.05000,0.79321) (-1.03014,0.81667) (-1.00171,0.85000) (-0.98333,0.87144) (-0.95879,0.90000) (-0.93333,0.92959) (-0.91580,0.95000) (-0.88732,0.98333) (-0.86667,1.00779) (-0.84535,1.03333) (-0.81818,1.06667) (-0.80000,1.08962) (-0.77919,1.11667) (-0.75470,1.15000) (-0.73333,1.18082) (-0.71667,1.20628) (-0.69997,1.23333) (-0.68092,1.26667) (-0.66355,1.30000) (-0.64780,1.33333) (-0.63333,1.36730) (-0.61667,1.41156) (-0.60396,1.45000) (-0.59408,1.48333) (-0.58333,1.52336) (-0.57281,1.56667) (-0.56522,1.60000) (-0.55432,1.65000) (-0.54710,1.68333) (-0.53598,1.73333) (-0.52815,1.76667) (-0.51667,1.81150) (-0.50591,1.85000) (-0.49553,1.88333) (-0.48333,1.91845) (-0.46667,1.96098) (-0.45000,1.99834) (-0.43333,2.03129) (-0.41667,2.06065) (-0.40000,2.08716) (-0.37934,2.11667) (-0.35251,2.15000) (-0.33333,2.17098) (-0.30339,2.20000) (-0.28333,2.21702) (-0.25000,2.24170) (-0.21667,2.26238) (-0.18333,2.27946) (-0.15000,2.29331) (-0.11667,2.30411) (-0.06667,2.31502) (-0.03333,2.31895) (0.01667,2.31994) (0.05451,2.31667) (0.10000,2.30846) (0.13333,2.29904) (0.17458,2.28333) (0.20883,2.26667) (0.23733,2.25000) (0.26667,2.22989) (0.30000,2.20296) (0.32121,2.18333) (0.35000,2.15280) (0.36667,2.13296) (0.39140,2.10000) (0.41315,2.06667) (0.43229,2.03333) (0.44925,2.00000) (0.46436,1.96667) (0.47786,1.93333) (0.48995,1.90000) (0.50081,1.86667) (0.51533,1.81667) (0.52405,1.78333) (0.53333,1.74457) (0.54347,1.70000) (0.55070,1.66667) (0.56155,1.61667) (0.56897,1.58333) (0.58083,1.53333) (0.58947,1.50000) (0.60000,1.46296) (0.61489,1.41667) (0.62702,1.38333) (0.64050,1.35000) (0.65547,1.31667) (0.67202,1.28333) (0.69023,1.25000) (0.71010,1.21667) (0.73163,1.18333) (0.75000,1.15665) (0.76679,1.13333) (0.79191,1.10000) (0.81667,1.06857) (0.83333,1.04797) (0.85921,1.01667) (0.88333,0.98805) (0.90153,0.96667) (0.93011,0.93333) (0.95000,0.91021) (0.97312,0.88333) (1.00000,0.85200) (1.01667,0.83249) (1.04427,0.80000) (1.06667,0.77343) (1.08633,0.75000) (1.11409,0.71667) (1.13333,0.69340) (1.15541,0.66667) (1.18281,0.63333) (1.20000,0.61240) (1.22398,0.58333) (1.25000,0.55187) (1.26667,0.53190) (1.29363,0.50000) (1.31667,0.47312) (1.33693,0.45000) (1.36667,0.41684) (1.38333,0.39875) (1.41386,0.36667) (1.43333,0.34688) (1.46431,0.31667) (1.48333,0.29882) (1.51667,0.26892) (1.53883,0.25000) (1.56667,0.22718) (1.60000,0.20131) (1.62434,0.18333) (1.65000,0.16514) (1.68333,0.14262) (1.71667,0.12127) (1.75000,0.10096) (1.78031,0.08333) (1.81016,0.06667) (1.84123,0.05000) (1.87354,0.03333) (1.90715,0.01667) (1.94206,0.00000) (1.97832,-0.01667) (2.01593,-0.03333) (2.05000,-0.04793) (2.08333,-0.06179) (2.11667,-0.07527) (2.15000,-0.08839) (2.18333,-0.10118) (2.22479,-0.11667) (2.26667,-0.13187) (2.30000,-0.14368) (2.33333,-0.15525) (2.36690,-0.16667) (2.41667,-0.18320) (2.45000,-0.19404) (2.48333,-0.20468) (2.52154,-0.21667) (2.56667,-0.23055) (2.60000,-0.24062) (2.63333,-0.25054) (2.68333,-0.26516) (2.71667,-0.27475) (2.75000,-0.28420) (2.80000,-0.29816) (2.83333,-0.30732) (2.86775,-0.31667) (2.91667,-0.32976) (2.95000,-0.33855) (2.99395,-0.35000) (3.03333,-0.36012) (3.06667,-0.36859) (3.11667,-0.38113) (3.15000,-0.38939) (3.19334,-0.40000) (3.23333,-0.40968) (3.26667,-0.41766) (3.31667,-0.42950) (3.35000,-0.43731) (3.40000,-0.44888) (3.43333,-0.45652) (3.47808,-0.46667) (3.51667,-0.47533) (3.55268,-0.48333) (3.60000,-0.49375) (3.63333,-0.50102) (3.68333,-0.51181) (3.71667,-0.51894) (3.76667,-0.52953) (3.80000,-0.53652) (3.85000,-0.54691) (3.88333,-0.55378) (3.93333,-0.56398) (3.96667,-0.57073) (4.01667,-0.58076) (4.05000,-0.58738) (4.10000,-0.59724) (4.13333,-0.60375) (4.18333,-0.61345) (4.21667,-0.61985) (4.26667,-0.62939) (4.30000,-0.63569) (4.35000,-0.64507) (4.38333,-0.65128) (4.43333,-0.66051) (4.46691,-0.66667) (4.51667,-0.67572) (4.55892,-0.68333) (4.60000,-0.69069) (4.65000,-0.69956) (4.68333,-0.70544) (4.73333,-0.71418) (4.76667,-0.71997) (4.81667,-0.72859) (4.85000,-0.73430) (4.90000,-0.74280) (4.94270,-0.75000) (4.98333,-0.75681) (5.00000,-0.75959)};
\draw[main,line width=.35pt] plot coordinates {(-5.00000,-0.42563) (-4.95000,-0.41669) (-4.91667,-0.41068) (-4.86667,-0.40160) (-4.83333,-0.39549) (-4.78333,-0.38626) (-4.75000,-0.38005) (-4.70000,-0.37066) (-4.66667,-0.36434) (-4.61667,-0.35479) (-4.58333,-0.34836) (-4.53333,-0.33864) (-4.50000,-0.33209) (-4.45000,-0.32219) (-4.41667,-0.31553) (-4.36667,-0.30545) (-4.33333,-0.29866) (-4.28333,-0.28838) (-4.25000,-0.28146) (-4.20000,-0.27098) (-4.16667,-0.26392) (-4.11667,-0.25323) (-4.08333,-0.24603) (-4.03333,-0.23511) (-4.00000,-0.22775) (-3.95029,-0.21667) (-3.91667,-0.20909) (-3.87677,-0.20000) (-3.83333,-0.19000) (-3.80000,-0.18224) (-3.75000,-0.17047) (-3.71667,-0.16253) (-3.66667,-0.15047) (-3.63333,-0.14234) (-3.59685,-0.13333) (-3.55000,-0.12163) (-3.51667,-0.11320) (-3.46667,-0.10038) (-3.43333,-0.09171) (-3.40000,-0.08295) (-3.35000,-0.06963) (-3.31667,-0.06062) (-3.27791,-0.05000) (-3.23333,-0.03761) (-3.20000,-0.02821) (-3.15966,-0.01667) (-3.11667,-0.00417) (-3.08333,0.00567) (-3.04660,0.01667) (-3.00000,0.03086) (-2.96667,0.04119) (-2.93333,0.05168) (-2.88653,0.06667) (-2.85000,0.07859) (-2.81667,0.08966) (-2.78333,0.10091) (-2.73759,0.11667) (-2.70000,0.12989) (-2.66667,0.14184) (-2.63333,0.15402) (-2.59937,0.16667) (-2.55562,0.18333) (-2.51667,0.19854) (-2.48333,0.21185) (-2.45000,0.22545) (-2.41667,0.23936) (-2.38333,0.25358) (-2.35000,0.26813) (-2.31603,0.28333) (-2.27973,0.30000) (-2.24440,0.31667) (-2.21002,0.33333) (-2.17656,0.35000) (-2.14400,0.36667) (-2.11230,0.38333) (-2.08144,0.40000) (-2.05000,0.41745) (-2.01667,0.43648) (-1.98333,0.45611) (-1.95000,0.47634) (-1.91667,0.49722) (-1.88652,0.51667) (-1.86135,0.53333) (-1.83333,0.55233) (-1.80000,0.57560) (-1.76667,0.59962) (-1.74362,0.61667) (-1.71667,0.63702) (-1.68333,0.66289) (-1.65769,0.68333) (-1.63333,0.70311) (-1.60000,0.73085) (-1.57750,0.75000) (-1.55000,0.77380) (-1.52039,0.80000) (-1.50000,0.81832) (-1.46667,0.84881) (-1.44748,0.86667) (-1.41667,0.89572) (-1.39481,0.91667) (-1.36667,0.94397) (-1.34362,0.96667) (-1.31667,0.99352) (-1.29377,1.01667) (-1.26667,1.04438) (-1.24520,1.06667) (-1.21667,1.09666) (-1.19793,1.11667) (-1.16717,1.15000) (-1.15000,1.16892) (-1.12230,1.20000) (-1.10000,1.22550) (-1.07903,1.25000) (-1.05115,1.28333) (-1.03333,1.30517) (-1.01092,1.33333) (-0.98518,1.36667) (-0.96667,1.39137) (-0.94825,1.41667) (-0.92475,1.45000) (-0.90211,1.48333) (-0.88333,1.51193) (-0.86667,1.53811) (-0.84904,1.56667) (-0.82907,1.60000) (-0.80970,1.63333) (-0.79082,1.66667) (-0.77235,1.70000) (-0.75416,1.73333) (-0.73615,1.76667) (-0.71819,1.80000) (-0.70014,1.83333) (-0.68333,1.86400) (-0.66667,1.89390) (-0.65000,1.92314) (-0.63333,1.95162) (-0.61411,1.98333) (-0.59293,2.01667) (-0.57051,2.05000) (-0.55000,2.07867) (-0.53333,2.10091) (-0.50747,2.13333) (-0.48333,2.16116) (-0.46289,2.18333) (-0.43333,2.21270) (-0.41084,2.23333) (-0.38333,2.25638) (-0.35000,2.28157) (-0.32276,2.30000) (-0.29537,2.31667) (-0.26458,2.33333) (-0.22893,2.35000) (-0.18573,2.36667) (-0.15000,2.37771) (-0.11667,2.38574) (-0.06667,2.39399) (-0.01667,2.39768) (0.03333,2.39695) (0.08333,2.39176) (0.12728,2.38333) (0.16667,2.37289) (0.20000,2.36163) (0.23333,2.34812) (0.26667,2.33228) (0.30000,2.31398) (0.33333,2.29306) (0.36667,2.26935) (0.39126,2.25000) (0.41667,2.22809) (0.44650,2.20000) (0.46667,2.17928) (0.49329,2.15000) (0.51667,2.12199) (0.53406,2.10000) (0.55880,2.06667) (0.58186,2.03333) (0.60000,2.00558) (0.61667,1.97909) (0.63433,1.95000) (0.65381,1.91667) (0.67267,1.88333) (0.69109,1.85000) (0.70921,1.81667) (0.72719,1.78333) (0.74515,1.75000) (0.76322,1.71667) (0.78153,1.68333) (0.80000,1.65031) (0.81667,1.62116) (0.83333,1.59273) (0.85000,1.56507) (0.86970,1.53333) (0.89111,1.50000) (0.91332,1.46667) (0.93333,1.43765) (0.95000,1.41422) (0.97265,1.38333) (0.99794,1.35000) (1.01667,1.32602) (1.03753,1.30000) (1.06499,1.26667) (1.08333,1.24492) (1.10769,1.21667) (1.13333,1.18753) (1.15205,1.16667) (1.18247,1.13333) (1.20000,1.11443) (1.22929,1.08333) (1.25000,1.06164) (1.27744,1.03333) (1.30000,1.01032) (1.32686,0.98333) (1.35000,0.96034) (1.37759,0.93333) (1.40000,0.91165) (1.42975,0.88333) (1.45000,0.86429) (1.48333,0.83349) (1.50183,0.81667) (1.53333,0.78847) (1.55821,0.76667) (1.58333,0.74499) (1.61667,0.71690) (1.63715,0.70000) (1.66667,0.67611) (1.69984,0.65000) (1.72152,0.63333) (1.75000,0.61189) (1.78333,0.58752) (1.81270,0.56667) (1.83675,0.55000) (1.86667,0.52977) (1.90000,0.50789) (1.93333,0.48670) (1.96583,0.46667) (1.99362,0.45000) (2.02213,0.43333) (2.05139,0.41667) (2.08333,0.39896) (2.11667,0.38100) (2.15000,0.36355) (2.18333,0.34658) (2.21667,0.33007) (2.25000,0.31399) (2.28333,0.29832) (2.31667,0.28304) (2.35334,0.26667) (2.39166,0.25000) (2.43103,0.23333) (2.46667,0.21862) (2.50000,0.20516) (2.53333,0.19199) (2.56667,0.17908) (2.60000,0.16643) (2.64428,0.15000) (2.68333,0.13584) (2.71667,0.12399) (2.75000,0.11235) (2.78602,0.10000) (2.83333,0.08411) (2.86667,0.07312) (2.90000,0.06232) (2.93864,0.05000) (2.98333,0.03601) (3.01667,0.02575) (3.05000,0.01564) (3.10000,0.00073) (3.13333,-0.00904) (3.16667,-0.01869) (3.21667,-0.03292) (3.25000,-0.04227) (3.28333,-0.05150) (3.33333,-0.06514) (3.36667,-0.07410) (3.40144,-0.08333) (3.45000,-0.09606) (3.48333,-0.10467) (3.53035,-0.11667) (3.56667,-0.12581) (3.60000,-0.13412) (3.65000,-0.14641) (3.68333,-0.15451) (3.73333,-0.16651) (3.76667,-0.17441) (3.80467,-0.18333) (3.85000,-0.19385) (3.88333,-0.20150) (3.93333,-0.21285) (3.96667,-0.22034) (4.01667,-0.23144) (4.05000,-0.23876) (4.10000,-0.24963) (4.13333,-0.25681) (4.17960,-0.26667) (4.21667,-0.27449) (4.25900,-0.28333) (4.30000,-0.29182) (4.33991,-0.30000) (4.38333,-0.30882) (4.42233,-0.31667) (4.46667,-0.32551) (4.50630,-0.33333) (4.55000,-0.34189) (4.59182,-0.35000) (4.63333,-0.35798) (4.67892,-0.36667) (4.71667,-0.37380) (4.76667,-0.38316) (4.80000,-0.38935) (4.85000,-0.39855) (4.88333,-0.40464) (4.93333,-0.41369) (4.96667,-0.41968) (5.00000,-0.42563)};
\draw[main,line width=.35pt] plot coordinates {(-5.00000,-0.09402) (-4.95000,-0.08454) (-4.91667,-0.07816) (-4.86667,-0.06851) (-4.83333,-0.06202) (-4.78333,-0.05219) (-4.75000,-0.04558) (-4.70000,-0.03556) (-4.66667,-0.02882) (-4.61667,-0.01861) (-4.58333,-0.01174) (-4.53333,-0.00133) (-4.50000,0.00568) (-4.45000,0.01631) (-4.41667,0.02346) (-4.37115,0.03333) (-4.33333,0.04162) (-4.29545,0.05000) (-4.25000,0.06016) (-4.21667,0.06770) (-4.16667,0.07913) (-4.13333,0.08683) (-4.08333,0.09853) (-4.05000,0.10641) (-4.00715,0.11667) (-3.96667,0.12647) (-3.93333,0.13463) (-3.88333,0.14702) (-3.85000,0.15539) (-3.80563,0.16667) (-3.76667,0.17669) (-3.73333,0.18538) (-3.68333,0.19858) (-3.65000,0.20750) (-3.61616,0.21667) (-3.56667,0.23026) (-3.53333,0.23955) (-3.49633,0.25000) (-3.45000,0.26327) (-3.41667,0.27297) (-3.38150,0.28333) (-3.33333,0.29775) (-3.30000,0.30790) (-3.26667,0.31818) (-3.21834,0.33333) (-3.18333,0.34450) (-3.15000,0.35529) (-3.11541,0.36667) (-3.06667,0.38298) (-3.03333,0.39435) (-3.00000,0.40589) (-2.96667,0.41763) (-2.92286,0.43333) (-2.88333,0.44779) (-2.85000,0.46022) (-2.81667,0.47285) (-2.78333,0.48572) (-2.74700,0.50000) (-2.70549,0.51667) (-2.66667,0.53260) (-2.63333,0.54656) (-2.60000,0.56080) (-2.56667,0.57532) (-2.53333,0.59013) (-2.50000,0.60524) (-2.46667,0.62067) (-2.43333,0.63643) (-2.40000,0.65251) (-2.36667,0.66895) (-2.33333,0.68575) (-2.30000,0.70292) (-2.26667,0.72047) (-2.23333,0.73842) (-2.20000,0.75678) (-2.16667,0.77556) (-2.13333,0.79478) (-2.10000,0.81444) (-2.06867,0.83333) (-2.04158,0.85000) (-2.01496,0.86667) (-1.98333,0.88687) (-1.95000,0.90864) (-1.91667,0.93090) (-1.88867,0.95000) (-1.86463,0.96667) (-1.83333,0.98875) (-1.80000,1.01278) (-1.77207,1.03333) (-1.74973,1.05000) (-1.71667,1.07510) (-1.68455,1.10000) (-1.66339,1.11667) (-1.63333,1.14067) (-1.60142,1.16667) (-1.58126,1.18333) (-1.55000,1.20954) (-1.52215,1.23333) (-1.50000,1.25250) (-1.46667,1.28184) (-1.44637,1.30000) (-1.41667,1.32693) (-1.39169,1.35000) (-1.36667,1.37341) (-1.33878,1.40000) (-1.31667,1.42137) (-1.28759,1.45000) (-1.26667,1.47091) (-1.23810,1.50000) (-1.21667,1.52215) (-1.19024,1.55000) (-1.16667,1.57521) (-1.14394,1.60000) (-1.11667,1.63019) (-1.09908,1.65000) (-1.06994,1.68333) (-1.05000,1.70644) (-1.02719,1.73333) (-1.00000,1.76572) (-0.98333,1.78580) (-0.95798,1.81667) (-0.93333,1.84679) (-0.91667,1.86731) (-0.89017,1.90000) (-0.86667,1.92885) (-0.84945,1.95000) (-0.82209,1.98333) (-0.80000,2.00981) (-0.78023,2.03333) (-0.75157,2.06667) (-0.73333,2.08736) (-0.70700,2.11667) (-0.68333,2.14202) (-0.65973,2.16667) (-0.63333,2.19300) (-0.60875,2.21667) (-0.58333,2.23990) (-0.55258,2.26667) (-0.53236,2.28333) (-0.50000,2.30852) (-0.46667,2.33263) (-0.44085,2.35000) (-0.41431,2.36667) (-0.38333,2.38467) (-0.35000,2.40231) (-0.31667,2.41820) (-0.28090,2.43333) (-0.23486,2.45000) (-0.20000,2.46059) (-0.16667,2.46903) (-0.11667,2.47879) (-0.08316,2.48333) (-0.03333,2.48733) (0.01667,2.48790) (0.06667,2.48504) (0.10000,2.48124) (0.15000,2.47268) (0.18333,2.46499) (0.23333,2.45051) (0.26667,2.43885) (0.30000,2.42552) (0.33333,2.41048) (0.36667,2.39372) (0.40000,2.37518) (0.43333,2.35482) (0.46566,2.33333) (0.48907,2.31667) (0.51667,2.29576) (0.55000,2.26880) (0.57200,2.25000) (0.60000,2.22473) (0.62623,2.20000) (0.65000,2.17645) (0.67584,2.15000) (0.70000,2.12418) (0.72211,2.10000) (0.75000,2.06843) (0.76667,2.04918) (0.79432,2.01667) (0.81667,1.98980) (0.83581,1.96667) (0.86307,1.93333) (0.88333,1.90834) (0.90369,1.88333) (0.93076,1.85000) (0.95000,1.82635) (0.97166,1.80000) (0.99923,1.76667) (1.01667,1.74578) (1.04132,1.71667) (1.06667,1.68708) (1.08444,1.66667) (1.11388,1.63333) (1.13333,1.61165) (1.15921,1.58333) (1.18333,1.55732) (1.20602,1.53333) (1.23333,1.50487) (1.25441,1.48333) (1.28333,1.45421) (1.30446,1.43333) (1.33333,1.40521) (1.35622,1.38333) (1.38333,1.35776) (1.40972,1.33333) (1.43333,1.31175) (1.46500,1.28333) (1.48383,1.26667) (1.51667,1.23804) (1.54163,1.21667) (1.56667,1.19550) (1.60000,1.16783) (1.62183,1.15000) (1.65000,1.12730) (1.68333,1.10095) (1.70600,1.08333) (1.73333,1.06238) (1.76667,1.03733) (1.79470,1.01667) (1.81766,1.00000) (1.85000,0.97693) (1.88333,0.95367) (1.91309,0.93333) (1.93792,0.91667) (1.96667,0.89769) (2.00000,0.87616) (2.03333,0.85512) (2.06667,0.83455) (2.09628,0.81667) (2.12443,0.80000) (2.15312,0.78333) (2.18333,0.76613) (2.21667,0.74755) (2.25000,0.72940) (2.28333,0.71165) (2.31667,0.69428) (2.35000,0.67730) (2.38333,0.66069) (2.41667,0.64443) (2.45000,0.62851) (2.48333,0.61292) (2.51667,0.59765) (2.55000,0.58269) (2.58647,0.56667) (2.62525,0.55000) (2.66490,0.53333) (2.70000,0.51889) (2.73333,0.50544) (2.76667,0.49223) (2.80000,0.47926) (2.83333,0.46651) (2.87739,0.45000) (2.91667,0.43558) (2.95000,0.42356) (2.98333,0.41174) (3.01696,0.40000) (3.06563,0.38333) (3.10000,0.37178) (3.13333,0.36075) (3.16667,0.34988) (3.21667,0.33386) (3.25000,0.32337) (3.28333,0.31302) (3.32593,0.30000) (3.36667,0.28774) (3.40000,0.27786) (3.43830,0.26667) (3.48333,0.25370) (3.51667,0.24424) (3.55561,0.23333) (3.60000,0.22108) (3.63333,0.21200) (3.67801,0.20000) (3.71667,0.18975) (3.75000,0.18102) (3.80000,0.16811) (3.83333,0.15961) (3.87144,0.15000) (3.91667,0.13874) (3.95000,0.13054) (4.00000,0.11839) (4.03333,0.11039) (4.07709,0.10000) (4.11667,0.09071) (4.15000,0.08297) (4.20000,0.07149) (4.23333,0.06392) (4.28333,0.05270) (4.31667,0.04529) (4.36667,0.03431) (4.40000,0.02706) (4.44832,0.01667) (4.48333,0.00921) (4.52701,0.00000) (4.56667,-0.00828) (4.60722,-0.01667) (4.65000,-0.02543) (4.68897,-0.03333) (4.73333,-0.04225) (4.77228,-0.05000) (4.81667,-0.05875) (4.85718,-0.06667) (4.90000,-0.07496) (4.94368,-0.08333) (4.98333,-0.09087) (5.00000,-0.09402)};
\draw[main,line width=.35pt] plot coordinates {(-5.00000,0.01589) (-4.96667,0.02231) (-4.91667,0.03203) (-4.88333,0.03857) (-4.83333,0.04847) (-4.80000,0.05513) (-4.75000,0.06522) (-4.71667,0.07201) (-4.66667,0.08230) (-4.63333,0.08922) (-4.58333,0.09971) (-4.55000,0.10678) (-4.50383,0.11667) (-4.46667,0.12470) (-4.42717,0.13333) (-4.38333,0.14301) (-4.35000,0.15044) (-4.30000,0.16172) (-4.26667,0.16932) (-4.21667,0.18085) (-4.18333,0.18862) (-4.13511,0.20000) (-4.10000,0.20838) (-4.06564,0.21667) (-4.01667,0.22862) (-3.98333,0.23686) (-3.93333,0.24937) (-3.90000,0.25782) (-3.86545,0.26667) (-3.81667,0.27933) (-3.78333,0.28809) (-3.73865,0.30000) (-3.70000,0.31044) (-3.66667,0.31955) (-3.61698,0.33333) (-3.58333,0.34280) (-3.55000,0.35230) (-3.50032,0.36667) (-3.46667,0.37655) (-3.43333,0.38646) (-3.38850,0.40000) (-3.35000,0.41181) (-3.31667,0.42219) (-3.28136,0.43333) (-3.23333,0.44875) (-3.20000,0.45964) (-3.16667,0.47069) (-3.12910,0.48333) (-3.08333,0.49902) (-3.05000,0.51065) (-3.01667,0.52247) (-2.98333,0.53447) (-2.94100,0.55000) (-2.90000,0.56534) (-2.86667,0.57804) (-2.83333,0.59096) (-2.80000,0.60411) (-2.76667,0.61749) (-2.72794,0.63333) (-2.68805,0.65000) (-2.65000,0.66624) (-2.61667,0.68076) (-2.58333,0.69555) (-2.55000,0.71063) (-2.51667,0.72601) (-2.48333,0.74169) (-2.45000,0.75768) (-2.41667,0.77401) (-2.38333,0.79066) (-2.35000,0.80766) (-2.31667,0.82502) (-2.28333,0.84274) (-2.25000,0.86083) (-2.21667,0.87931) (-2.18333,0.89819) (-2.15137,0.91667) (-2.12309,0.93333) (-2.09529,0.95000) (-2.06667,0.96746) (-2.03333,0.98820) (-2.00000,1.00938) (-1.96667,1.03101) (-1.93797,1.05000) (-1.91317,1.06667) (-1.88333,1.08704) (-1.85000,1.11026) (-1.81753,1.13333) (-1.79445,1.15000) (-1.76667,1.17034) (-1.73333,1.19519) (-1.70507,1.21667) (-1.68333,1.23338) (-1.65000,1.25943) (-1.61999,1.28333) (-1.59933,1.30000) (-1.56667,1.32674) (-1.53877,1.35000) (-1.51667,1.36865) (-1.48333,1.39724) (-1.46106,1.41667) (-1.43333,1.44114) (-1.40494,1.46667) (-1.38333,1.48632) (-1.35055,1.51667) (-1.33280,1.53333) (-1.30000,1.56459) (-1.28065,1.58333) (-1.25000,1.61341) (-1.23003,1.63333) (-1.20000,1.66367) (-1.18085,1.68333) (-1.15000,1.71537) (-1.13294,1.73333) (-1.10164,1.76667) (-1.08333,1.78636) (-1.05549,1.81667) (-1.03333,1.84090) (-1.01003,1.86667) (-0.98333,1.89622) (-0.96498,1.91667) (-0.93501,1.95000) (-0.91667,1.97033) (-0.88988,2.00000) (-0.86667,2.02539) (-0.84409,2.05000) (-0.81667,2.07930) (-0.79709,2.10000) (-0.76667,2.13138) (-0.74825,2.15000) (-0.71667,2.18101) (-0.69679,2.20000) (-0.66667,2.22772) (-0.64160,2.25000) (-0.61667,2.27119) (-0.58333,2.29826) (-0.55950,2.31667) (-0.53333,2.33586) (-0.50000,2.35891) (-0.46667,2.38036) (-0.43369,2.40000) (-0.40340,2.41667) (-0.37041,2.43333) (-0.33395,2.45000) (-0.30000,2.46388) (-0.26667,2.47600) (-0.23333,2.48662) (-0.18333,2.49980) (-0.15000,2.50684) (-0.10000,2.51468) (-0.06667,2.51815) (-0.01667,2.52078) (0.03333,2.52025) (0.08236,2.51667) (0.11667,2.51243) (0.16667,2.50351) (0.20000,2.49579) (0.24401,2.48333) (0.28333,2.47012) (0.31667,2.45727) (0.35000,2.44290) (0.38333,2.42699) (0.41667,2.40953) (0.45000,2.39048) (0.48333,2.36984) (0.51320,2.35000) (0.53688,2.33333) (0.56667,2.31118) (0.60000,2.28493) (0.62212,2.26667) (0.65000,2.24258) (0.67887,2.21667) (0.70000,2.19692) (0.73144,2.16667) (0.75000,2.14822) (0.78105,2.11667) (0.80000,2.09690) (0.82858,2.06667) (0.85000,2.04352) (0.87472,2.01667) (0.90000,1.98877) (0.92001,1.96667) (0.94999,1.93333) (0.96667,1.91476) (0.99499,1.88333) (1.01667,1.85927) (1.04027,1.83333) (1.06667,1.80443) (1.08616,1.78333) (1.11667,1.75059) (1.13333,1.73291) (1.16474,1.70000) (1.18333,1.68075) (1.21349,1.65000) (1.23333,1.63000) (1.26361,1.60000) (1.28333,1.58070) (1.31524,1.55000) (1.33333,1.53283) (1.36667,1.50168) (1.38662,1.48333) (1.41667,1.45606) (1.44216,1.43333) (1.46667,1.41174) (1.49948,1.38333) (1.51902,1.36667) (1.55000,1.34058) (1.57892,1.31667) (1.60000,1.29945) (1.63333,1.27265) (1.66202,1.25000) (1.68340,1.23333) (1.71667,1.20780) (1.74917,1.18333) (1.77167,1.16667) (1.80000,1.14596) (1.83333,1.12204) (1.86467,1.10000) (1.88875,1.08333) (1.91667,1.06429) (1.95000,1.04199) (1.98333,1.02014) (2.01470,1.00000) (2.04112,0.98333) (2.06797,0.96667) (2.10000,0.94714) (2.13333,0.92725) (2.16667,0.90777) (2.20000,0.88870) (2.23333,0.87003) (2.26667,0.85174) (2.30000,0.83384) (2.33264,0.81667) (2.36496,0.80000) (2.39793,0.78333) (2.43159,0.76667) (2.46594,0.75000) (2.50000,0.73381) (2.53333,0.71828) (2.56667,0.70305) (2.60000,0.68812) (2.63333,0.67346) (2.66667,0.65908) (2.70000,0.64496) (2.73333,0.63110) (2.76870,0.61667) (2.81037,0.60000) (2.85000,0.58448) (2.88333,0.57166) (2.91667,0.55906) (2.95000,0.54667) (2.98649,0.53333) (3.03298,0.51667) (3.06667,0.50482) (3.10000,0.49327) (3.13333,0.48190) (3.17876,0.46667) (3.21667,0.45418) (3.25000,0.44336) (3.28333,0.43270) (3.33333,0.41698) (3.36667,0.40667) (3.40000,0.39650) (3.44381,0.38333) (3.48333,0.37164) (3.51667,0.36191) (3.55804,0.35000) (3.60000,0.33810) (3.63333,0.32877) (3.67718,0.31667) (3.71667,0.30592) (3.75000,0.29696) (3.80000,0.28370) (3.83333,0.27498) (3.86667,0.26635) (3.91667,0.25358) (3.95000,0.24518) (3.99756,0.23333) (4.03333,0.22453) (4.06667,0.21642) (4.11667,0.20439) (4.15000,0.19647) (4.20000,0.18473) (4.23333,0.17699) (4.27827,0.16667) (4.31667,0.15794) (4.35199,0.15000) (4.40000,0.13932) (4.43333,0.13198) (4.48333,0.12109) (4.51667,0.11391) (4.56667,0.10324) (4.60000,0.09620) (4.65000,0.08575) (4.68333,0.07885) (4.73333,0.06861) (4.76667,0.06184) (4.81667,0.05179) (4.85000,0.04515) (4.90000,0.03529) (4.93333,0.02878) (4.98333,0.01909) (5.00000,0.01589)};
\draw[main,line width=.35pt] plot coordinates {(-5.00000,0.34344) (-4.96667,0.35016) (-4.91667,0.36034) (-4.88333,0.36719) (-4.83333,0.37757) (-4.80000,0.38456) (-4.75000,0.39514) (-4.71667,0.40227) (-4.66667,0.41308) (-4.63333,0.42036) (-4.58333,0.43140) (-4.55000,0.43883) (-4.50049,0.45000) (-4.46667,0.45771) (-4.42778,0.46667) (-4.38333,0.47701) (-4.35000,0.48486) (-4.30000,0.49676) (-4.26667,0.50479) (-4.21796,0.51667) (-4.18333,0.52521) (-4.15000,0.53352) (-4.10000,0.54613) (-4.06667,0.55464) (-4.02019,0.56667) (-3.98333,0.57632) (-3.95000,0.58515) (-3.90000,0.59857) (-3.86667,0.60764) (-3.83333,0.61682) (-3.78333,0.63076) (-3.75000,0.64019) (-3.71574,0.65000) (-3.66667,0.66425) (-3.63333,0.67407) (-3.60000,0.68401) (-3.55000,0.69915) (-3.51667,0.70940) (-3.48333,0.71978) (-3.44044,0.73333) (-3.40000,0.74631) (-3.36667,0.75717) (-3.33333,0.76817) (-3.28812,0.78333) (-3.25000,0.79633) (-3.21667,0.80787) (-3.18333,0.81957) (-3.14474,0.83333) (-3.10000,0.84957) (-3.06667,0.86187) (-3.03333,0.87436) (-3.00000,0.88704) (-2.96645,0.90000) (-2.92405,0.91667) (-2.88333,0.93297) (-2.85000,0.94656) (-2.81667,0.96037) (-2.78333,0.97439) (-2.75000,0.98865) (-2.71667,1.00314) (-2.68333,1.01788) (-2.64895,1.03333) (-2.61252,1.05000) (-2.57675,1.06667) (-2.54160,1.08333) (-2.50708,1.10000) (-2.47315,1.11667) (-2.43980,1.13333) (-2.40702,1.15000) (-2.37477,1.16667) (-2.34306,1.18333) (-2.31185,1.20000) (-2.28113,1.21667) (-2.25000,1.23382) (-2.21667,1.25250) (-2.18333,1.27150) (-2.15000,1.29083) (-2.11667,1.31048) (-2.08333,1.33047) (-2.05131,1.35000) (-2.02437,1.36667) (-1.99777,1.38333) (-1.96667,1.40309) (-1.93333,1.42460) (-1.90000,1.44647) (-1.86970,1.46667) (-1.84499,1.48333) (-1.81667,1.50267) (-1.78333,1.52576) (-1.75000,1.54921) (-1.72553,1.56667) (-1.70000,1.58505) (-1.66667,1.60938) (-1.63435,1.63333) (-1.61212,1.65000) (-1.58333,1.67179) (-1.55000,1.69737) (-1.52519,1.71667) (-1.50000,1.73640) (-1.46667,1.76285) (-1.44117,1.78333) (-1.41667,1.80314) (-1.38333,1.83040) (-1.35963,1.85000) (-1.33333,1.87184) (-1.30000,1.89983) (-1.28009,1.91667) (-1.25000,1.94220) (-1.22142,1.96667) (-1.20000,1.98503) (-1.16667,2.01372) (-1.14396,2.03333) (-1.11667,2.05685) (-1.08599,2.08333) (-1.06660,2.10000) (-1.03333,2.12847) (-1.00810,2.15000) (-0.98333,2.17088) (-0.95000,2.19875) (-0.92834,2.21667) (-0.90000,2.23972) (-0.86667,2.26640) (-0.84512,2.28333) (-0.81667,2.30519) (-0.78333,2.33018) (-0.75614,2.35000) (-0.73261,2.36667) (-0.70000,2.38903) (-0.66667,2.41100) (-0.63333,2.43201) (-0.60345,2.45000) (-0.57444,2.46667) (-0.54394,2.48333) (-0.51168,2.50000) (-0.47726,2.51667) (-0.44018,2.53333) (-0.40000,2.54987) (-0.36667,2.56239) (-0.33333,2.57381) (-0.30000,2.58411) (-0.25000,2.59754) (-0.21667,2.60513) (-0.16667,2.61446) (-0.13333,2.61932) (-0.08333,2.62460) (-0.03333,2.62743) (0.01667,2.62783) (0.06667,2.62581) (0.11667,2.62135) (0.15224,2.61667) (0.20000,2.60852) (0.23956,2.60000) (0.28333,2.58888) (0.31667,2.57910) (0.35449,2.56667) (0.39967,2.55000) (0.43333,2.53627) (0.46667,2.52157) (0.50000,2.50579) (0.53333,2.48892) (0.56667,2.47099) (0.60000,2.45201) (0.63119,2.43333) (0.65786,2.41667) (0.68356,2.40000) (0.71667,2.37770) (0.75000,2.35436) (0.77909,2.33333) (0.80152,2.31667) (0.83333,2.29242) (0.86634,2.26667) (0.88728,2.25000) (0.91667,2.22618) (0.94852,2.20000) (0.96852,2.18333) (1.00000,2.15681) (1.02769,2.13333) (1.05000,2.11423) (1.08333,2.08559) (1.10534,2.06667) (1.13333,2.04246) (1.16328,2.01667) (1.18333,1.99937) (1.21667,1.97071) (1.24091,1.95000) (1.26667,1.92802) (1.29980,1.90000) (1.31964,1.88333) (1.35000,1.85795) (1.37980,1.83333) (1.40009,1.81667) (1.43333,1.78962) (1.46192,1.76667) (1.48333,1.74959) (1.51667,1.72330) (1.54663,1.70000) (1.56826,1.68333) (1.60000,1.65913) (1.63333,1.63409) (1.65682,1.61667) (1.68333,1.59717) (1.71667,1.57301) (1.74890,1.55000) (1.77254,1.53333) (1.80000,1.51417) (1.83333,1.49125) (1.86667,1.46869) (1.89470,1.45000) (1.91999,1.43333) (1.95000,1.41381) (1.98333,1.39246) (2.01667,1.37146) (2.05000,1.35080) (2.07863,1.33333) (2.10632,1.31667) (2.13439,1.30000) (2.16667,1.28113) (2.20000,1.26196) (2.23333,1.24312) (2.26667,1.22459) (2.30000,1.20638) (2.33333,1.18848) (2.36667,1.17089) (2.40000,1.15359) (2.43333,1.13659) (2.46667,1.11987) (2.50000,1.10344) (2.53333,1.08729) (2.56667,1.07141) (2.60000,1.05579) (2.63333,1.04043) (2.66667,1.02533) (2.70000,1.01048) (2.73333,0.99587) (2.76667,0.98150) (2.80164,0.96667) (2.84167,0.95000) (2.88245,0.93333) (2.91667,0.91959) (2.95000,0.90642) (2.98333,0.89345) (3.01667,0.88068) (3.05383,0.86667) (3.09883,0.85000) (3.13333,0.83744) (3.16667,0.82548) (3.20000,0.81370) (3.23936,0.80000) (3.28333,0.78495) (3.31667,0.77372) (3.35000,0.76265) (3.38864,0.75000) (3.43333,0.73560) (3.46667,0.72502) (3.50000,0.71457) (3.54723,0.70000) (3.58333,0.68903) (3.61667,0.67903) (3.65844,0.66667) (3.70000,0.65454) (3.73333,0.64495) (3.77422,0.63333) (3.81667,0.62144) (3.85000,0.61222) (3.89475,0.60000) (3.93333,0.58960) (3.96667,0.58072) (4.01667,0.56758) (4.05000,0.55893) (4.08481,0.55000) (4.13333,0.53770) (4.16667,0.52935) (4.21667,0.51698) (4.25000,0.50884) (4.28654,0.50000) (4.33333,0.48881) (4.36667,0.48093) (4.41667,0.46924) (4.45000,0.46154) (4.50000,0.45011) (4.53333,0.44258) (4.57463,0.43333) (4.61667,0.42402) (4.65021,0.41667) (4.70000,0.40586) (4.73333,0.39870) (4.78333,0.38807) (4.81667,0.38106) (4.86667,0.37064) (4.90000,0.36376) (4.95000,0.35354) (4.98333,0.34680) (5.00000,0.34344)};
\draw[main,line width=.35pt] plot coordinates {(-5.00000,0.66735) (-4.95000,0.67781) (-4.91667,0.68485) (-4.86667,0.69552) (-4.83333,0.70271) (-4.78333,0.71359) (-4.75000,0.72093) (-4.70000,0.73204) (-4.66667,0.73953) (-4.62057,0.75000) (-4.58333,0.75854) (-4.54828,0.76667) (-4.50000,0.77797) (-4.46667,0.78587) (-4.41667,0.79785) (-4.38333,0.80592) (-4.33949,0.81667) (-4.30000,0.82645) (-4.26667,0.83480) (-4.21667,0.84747) (-4.18333,0.85603) (-4.14234,0.86667) (-4.10000,0.87778) (-4.06667,0.88664) (-4.01705,0.90000) (-3.98333,0.90919) (-3.95000,0.91838) (-3.90000,0.93234) (-3.86667,0.94177) (-3.83333,0.95131) (-3.78333,0.96581) (-3.75000,0.97562) (-3.71667,0.98553) (-3.66872,1.00000) (-3.63333,1.01083) (-3.60000,1.02115) (-3.56118,1.03333) (-3.51667,1.04751) (-3.48333,1.05827) (-3.45000,1.06917) (-3.40729,1.08333) (-3.36667,1.09701) (-3.33333,1.10839) (-3.30000,1.11992) (-3.26178,1.13333) (-3.21667,1.14941) (-3.18333,1.16147) (-3.15000,1.17370) (-3.11667,1.18609) (-3.07977,1.20000) (-3.03628,1.21667) (-3.00000,1.23079) (-2.96667,1.24396) (-2.93333,1.25731) (-2.90000,1.27085) (-2.86667,1.28459) (-2.82981,1.30000) (-2.79058,1.31667) (-2.75198,1.33333) (-2.71667,1.34883) (-2.68333,1.36366) (-2.65000,1.37872) (-2.61667,1.39399) (-2.58333,1.40948) (-2.55000,1.42519) (-2.51667,1.44113) (-2.48333,1.45730) (-2.45000,1.47370) (-2.41667,1.49033) (-2.38333,1.50720) (-2.35000,1.52431) (-2.31667,1.54166) (-2.28333,1.55925) (-2.25000,1.57708) (-2.21667,1.59517) (-2.18333,1.61350) (-2.15000,1.63208) (-2.11826,1.65000) (-2.08910,1.66667) (-2.06025,1.68333) (-2.03172,1.70000) (-2.00000,1.71873) (-1.96667,1.73867) (-1.93333,1.75885) (-1.90000,1.77928) (-1.86667,1.79996) (-1.84001,1.81667) (-1.81363,1.83333) (-1.78333,1.85265) (-1.75000,1.87413) (-1.71667,1.89583) (-1.68499,1.91667) (-1.65986,1.93333) (-1.63333,1.95102) (-1.60000,1.97343) (-1.56667,1.99603) (-1.53647,2.01667) (-1.51220,2.03333) (-1.48333,2.05324) (-1.45000,2.07635) (-1.41667,2.09960) (-1.39228,2.11667) (-1.36667,2.13460) (-1.33333,2.15799) (-1.30000,2.18142) (-1.27359,2.20000) (-1.24981,2.21667) (-1.21667,2.23986) (-1.18333,2.26310) (-1.15420,2.28333) (-1.13002,2.30000) (-1.10000,2.32053) (-1.06667,2.34312) (-1.03333,2.36545) (-1.00629,2.38333) (-0.98073,2.40000) (-0.95000,2.41972) (-0.91667,2.44072) (-0.88333,2.46125) (-0.85000,2.48127) (-0.81798,2.50000) (-0.78867,2.51667) (-0.75849,2.53333) (-0.72732,2.55000) (-0.69500,2.56667) (-0.66134,2.58333) (-0.62611,2.60000) (-0.58901,2.61667) (-0.55000,2.63319) (-0.51667,2.64647) (-0.48333,2.65896) (-0.45000,2.67065) (-0.41087,2.68333) (-0.36667,2.69634) (-0.33333,2.70518) (-0.28422,2.71667) (-0.25000,2.72364) (-0.20000,2.73218) (-0.16667,2.73684) (-0.11667,2.74224) (-0.06667,2.74573) (-0.01667,2.74731) (0.03333,2.74699) (0.08333,2.74478) (0.13333,2.74065) (0.18333,2.73461) (0.21667,2.72954) (0.26667,2.72036) (0.30000,2.71318) (0.35000,2.70086) (0.38333,2.69162) (0.41667,2.68154) (0.46158,2.66667) (0.50000,2.65281) (0.53333,2.63993) (0.56667,2.62626) (0.60000,2.61183) (0.63333,2.59665) (0.66667,2.58074) (0.70000,2.56412) (0.73333,2.54682) (0.76667,2.52886) (0.80000,2.51026) (0.83333,2.49107) (0.86667,2.47132) (0.90000,2.45105) (0.92851,2.43333) (0.95482,2.41667) (0.98333,2.39830) (1.01667,2.37647) (1.05000,2.35432) (1.08119,2.33333) (1.10571,2.31667) (1.13333,2.29771) (1.16667,2.27467) (1.20000,2.25150) (1.22604,2.23333) (1.25000,2.21654) (1.28333,2.19312) (1.31667,2.16971) (1.34475,2.15000) (1.36848,2.13333) (1.40000,2.11124) (1.43333,2.08796) (1.46398,2.06667) (1.48804,2.05000) (1.51667,2.03024) (1.55000,2.00739) (1.58333,1.98471) (1.61006,1.96667) (1.63486,1.95000) (1.66667,1.92879) (1.70000,1.90676) (1.73333,1.88496) (1.76157,1.86667) (1.78749,1.85000) (1.81667,1.83140) (1.85000,1.81037) (1.88333,1.78958) (1.91667,1.76904) (1.94793,1.75000) (1.97558,1.73333) (2.00350,1.71667) (2.03333,1.69905) (2.06667,1.67960) (2.10000,1.66040) (2.13333,1.64145) (2.16667,1.62275) (2.20000,1.60430) (2.23333,1.58610) (2.26667,1.56814) (2.30000,1.55043) (2.33262,1.53333) (2.36485,1.51667) (2.39752,1.50000) (2.43065,1.48333) (2.46425,1.46667) (2.49834,1.45000) (2.53293,1.43333) (2.56667,1.41731) (2.60000,1.40171) (2.63333,1.38633) (2.66667,1.37116) (2.70000,1.35622) (2.73333,1.34148) (2.76667,1.32695) (2.80000,1.31263) (2.83333,1.29851) (2.86970,1.28333) (2.91028,1.26667) (2.95000,1.25062) (2.98333,1.23735) (3.01667,1.22427) (3.05000,1.21137) (3.08333,1.19864) (3.12405,1.18333) (3.16667,1.16757) (3.20000,1.15542) (3.23333,1.14343) (3.26667,1.13160) (3.30939,1.11667) (3.35000,1.10268) (3.38333,1.09137) (3.41667,1.08020) (3.45763,1.06667) (3.50000,1.05287) (3.53333,1.04217) (3.56667,1.03160) (3.61444,1.01667) (3.65000,1.00571) (3.68333,0.99556) (3.72404,0.98333) (3.76667,0.97070) (3.80000,0.96095) (3.83790,0.95000) (3.88333,0.93704) (3.91667,0.92766) (3.95619,0.91667) (4.00000,0.90463) (4.03333,0.89559) (4.07909,0.88333) (4.11667,0.87339) (4.15000,0.86467) (4.20000,0.85174) (4.23333,0.84323) (4.27251,0.83333) (4.31667,0.82231) (4.35000,0.81408) (4.40000,0.80188) (4.43333,0.79384) (4.47735,0.78333) (4.51667,0.77405) (4.55000,0.76627) (4.60000,0.75471) (4.63333,0.74709) (4.68333,0.73578) (4.71667,0.72832) (4.76667,0.71725) (4.80000,0.70995) (4.84586,0.70000) (4.88333,0.69195) (4.92384,0.68333) (4.96667,0.67431) (5.00000,0.66735)};
\draw[main,line width=.35pt] plot coordinates {(-5.00000,0.98721) (-4.95000,0.99793) (-4.91667,1.00514) (-4.86667,1.01607) (-4.83333,1.02343) (-4.78894,1.03333) (-4.75000,1.04210) (-4.71527,1.05000) (-4.66667,1.06117) (-4.63333,1.06891) (-4.58333,1.08064) (-4.55000,1.08855) (-4.50226,1.10000) (-4.46667,1.10863) (-4.43333,1.11679) (-4.38333,1.12916) (-4.35000,1.13751) (-4.30070,1.15000) (-4.26667,1.15872) (-4.23333,1.16734) (-4.18333,1.18043) (-4.15000,1.18927) (-4.10994,1.20000) (-4.06667,1.21173) (-4.03333,1.22087) (-3.98844,1.23333) (-3.95000,1.24413) (-3.91667,1.25361) (-3.87127,1.26667) (-3.83333,1.27772) (-3.80000,1.28754) (-3.75824,1.30000) (-3.71667,1.31256) (-3.68333,1.32276) (-3.64915,1.33333) (-3.60000,1.34875) (-3.56667,1.35934) (-3.53333,1.37006) (-3.49256,1.38333) (-3.45000,1.39738) (-3.41667,1.40852) (-3.38333,1.41980) (-3.34386,1.43333) (-3.30000,1.44857) (-3.26667,1.46031) (-3.23333,1.47220) (-3.20000,1.48422) (-3.15687,1.50000) (-3.11667,1.51491) (-3.08333,1.52745) (-3.05000,1.54013) (-3.01667,1.55297) (-2.98156,1.56667) (-2.93943,1.58333) (-2.90000,1.59915) (-2.86667,1.61270) (-2.83333,1.62642) (-2.80000,1.64031) (-2.76667,1.65436) (-2.73333,1.66859) (-2.69921,1.68333) (-2.66114,1.70000) (-2.62356,1.71667) (-2.58648,1.73333) (-2.55000,1.74994) (-2.51667,1.76530) (-2.48333,1.78084) (-2.45000,1.79657) (-2.41667,1.81247) (-2.38333,1.82855) (-2.35000,1.84482) (-2.31667,1.86126) (-2.28333,1.87788) (-2.25000,1.89468) (-2.21667,1.91166) (-2.18333,1.92882) (-2.15000,1.94614) (-2.11667,1.96365) (-2.08333,1.98132) (-2.05000,1.99916) (-2.01757,2.01667) (-1.98697,2.03333) (-1.95660,2.05000) (-1.92645,2.06667) (-1.89651,2.08333) (-1.86667,2.10006) (-1.83333,2.11887) (-1.80000,2.13780) (-1.76667,2.15686) (-1.73333,2.17604) (-1.70000,2.19531) (-1.66667,2.21468) (-1.63471,2.23333) (-1.60624,2.25000) (-1.57783,2.26667) (-1.54946,2.28333) (-1.51667,2.30263) (-1.48333,2.32227) (-1.45000,2.34190) (-1.41667,2.36152) (-1.38333,2.38110) (-1.35106,2.40000) (-1.32250,2.41667) (-1.29380,2.43333) (-1.26494,2.45000) (-1.23333,2.46813) (-1.20000,2.48709) (-1.16667,2.50586) (-1.13333,2.52442) (-1.10000,2.54274) (-1.06667,2.56080) (-1.03333,2.57856) (-1.00000,2.59601) (-0.96667,2.61311) (-0.93333,2.62985) (-0.90000,2.64619) (-0.86667,2.66211) (-0.83333,2.67760) (-0.80000,2.69263) (-0.76667,2.70718) (-0.73333,2.72123) (-0.70000,2.73476) (-0.66077,2.75000) (-0.61667,2.76622) (-0.58333,2.77782) (-0.55000,2.78883) (-0.51414,2.80000) (-0.46667,2.81372) (-0.43333,2.82261) (-0.38936,2.83333) (-0.35000,2.84204) (-0.30951,2.85000) (-0.26667,2.85744) (-0.21667,2.86470) (-0.18333,2.86872) (-0.13333,2.87351) (-0.08333,2.87678) (-0.03333,2.87854) (0.01667,2.87879) (0.06667,2.87754) (0.11667,2.87477) (0.16667,2.87048) (0.20093,2.86667) (0.25000,2.86003) (0.30000,2.85176) (0.33333,2.84544) (0.38333,2.83474) (0.41667,2.82681) (0.45582,2.81667) (0.50000,2.80423) (0.53333,2.79411) (0.56685,2.78333) (0.61541,2.76667) (0.65000,2.75406) (0.68333,2.74133) (0.71667,2.72806) (0.75000,2.71427) (0.78333,2.69997) (0.82076,2.68333) (0.85698,2.66667) (0.89212,2.65000) (0.92631,2.63333) (0.95967,2.61667) (0.99231,2.60000) (1.02431,2.58333) (1.05575,2.56667) (1.08669,2.55000) (1.11720,2.53333) (1.15000,2.51518) (1.18333,2.49650) (1.21667,2.47762) (1.25000,2.45857) (1.28333,2.43937) (1.31667,2.42004) (1.35000,2.40062) (1.37954,2.38333) (1.40794,2.36667) (1.43627,2.35000) (1.46667,2.33209) (1.50000,2.31245) (1.53333,2.29281) (1.56667,2.27321) (1.60000,2.25364) (1.63333,2.23413) (1.66328,2.21667) (1.69195,2.20000) (1.72073,2.18333) (1.75000,2.16645) (1.78333,2.14732) (1.81667,2.12832) (1.85000,2.10944) (1.88333,2.09069) (1.91667,2.07208) (1.95000,2.05362) (1.98333,2.03531) (2.01667,2.01716) (2.04844,2.00000) (2.07957,1.98333) (2.11097,1.96667) (2.14266,1.95000) (2.17464,1.93333) (2.20693,1.91667) (2.23955,1.90000) (2.27250,1.88333) (2.30581,1.86667) (2.33948,1.85000) (2.37352,1.83333) (2.40796,1.81667) (2.44279,1.80000) (2.47805,1.78333) (2.51373,1.76667) (2.54987,1.75000) (2.58333,1.73475) (2.61667,1.71974) (2.65000,1.70491) (2.68333,1.69025) (2.71667,1.67576) (2.75000,1.66145) (2.78333,1.64732) (2.81671,1.63333) (2.85702,1.61667) (2.89791,1.60000) (2.93333,1.58576) (2.96667,1.57252) (3.00000,1.55945) (3.03333,1.54653) (3.06783,1.53333) (3.11200,1.51667) (3.15000,1.50253) (3.18333,1.49029) (3.21667,1.47819) (3.25000,1.46624) (3.29594,1.45000) (3.33333,1.43696) (3.36667,1.42549) (3.40000,1.41415) (3.44214,1.40000) (3.48333,1.38636) (3.51667,1.37546) (3.55000,1.36468) (3.59604,1.35000) (3.63333,1.33826) (3.66667,1.32790) (3.70321,1.31667) (3.75000,1.30247) (3.78333,1.29249) (3.81667,1.28262) (3.86667,1.26800) (3.90000,1.25838) (3.93333,1.24886) (3.98333,1.23476) (4.01667,1.22548) (4.05000,1.21629) (4.10000,1.20268) (4.13333,1.19371) (4.17237,1.18333) (4.21667,1.17168) (4.25000,1.16302) (4.30000,1.15018) (4.33333,1.14171) (4.36667,1.13333) (4.41667,1.12089) (4.45000,1.11270) (4.50000,1.10054) (4.53333,1.09253) (4.57196,1.08333) (4.61667,1.07280) (4.65000,1.06503) (4.70000,1.05349) (4.73333,1.04588) (4.78333,1.03459) (4.81667,1.02714) (4.86397,1.01667) (4.90000,1.00877) (4.94041,1.00000) (4.98333,0.99077) (5.00000,0.98721)};
\draw[main,line width=.35pt] plot coordinates {(-5.00000,1.30276) (-4.95000,1.31362) (-4.91667,1.32093) (-4.86667,1.33201) (-4.83333,1.33947) (-4.78672,1.35000) (-4.75000,1.35838) (-4.71400,1.36667) (-4.66667,1.37768) (-4.63333,1.38551) (-4.58333,1.39738) (-4.55000,1.40538) (-4.50344,1.41667) (-4.46667,1.42567) (-4.43333,1.43392) (-4.38333,1.44642) (-4.35000,1.45484) (-4.30372,1.46667) (-4.26667,1.47623) (-4.23333,1.48493) (-4.18333,1.49812) (-4.15000,1.50701) (-4.11417,1.51667) (-4.06667,1.52961) (-4.03333,1.53880) (-3.99313,1.55000) (-3.95000,1.56215) (-3.91667,1.57165) (-3.87612,1.58333) (-3.83333,1.59580) (-3.80000,1.60563) (-3.76297,1.61667) (-3.71667,1.63063) (-3.68333,1.64080) (-3.65000,1.65108) (-3.60005,1.66667) (-3.56667,1.67722) (-3.53333,1.68786) (-3.49574,1.70000) (-3.45000,1.71495) (-3.41667,1.72598) (-3.38333,1.73712) (-3.34526,1.75000) (-3.30000,1.76549) (-3.26667,1.77705) (-3.23333,1.78872) (-3.20000,1.80052) (-3.15493,1.81667) (-3.11667,1.83055) (-3.08333,1.84278) (-3.05000,1.85513) (-3.01667,1.86762) (-2.97522,1.88333) (-2.93333,1.89940) (-2.90000,1.91234) (-2.86667,1.92541) (-2.83333,1.93862) (-2.80000,1.95196) (-2.76365,1.96667) (-2.72291,1.98333) (-2.68333,1.99971) (-2.65000,2.01364) (-2.61667,2.02772) (-2.58333,2.04193) (-2.55000,2.05627) (-2.51667,2.07075) (-2.48333,2.08536) (-2.45000,2.10010) (-2.41288,2.11667) (-2.37589,2.13333) (-2.33924,2.15000) (-2.30290,2.16667) (-2.26687,2.18333) (-2.23333,2.19898) (-2.20000,2.21464) (-2.16667,2.23041) (-2.13333,2.24629) (-2.10000,2.26228) (-2.06667,2.27837) (-2.03333,2.29455) (-2.00000,2.31083) (-1.96667,2.32719) (-1.93333,2.34363) (-1.90000,2.36014) (-1.86667,2.37671) (-1.83333,2.39335) (-1.80000,2.41003) (-1.76667,2.42676) (-1.73333,2.44351) (-1.70000,2.46030) (-1.66667,2.47709) (-1.63333,2.49388) (-1.60000,2.51067) (-1.56667,2.52743) (-1.53333,2.54416) (-1.50000,2.56084) (-1.46667,2.57746) (-1.43333,2.59401) (-1.40000,2.61046) (-1.36667,2.62681) (-1.33333,2.64303) (-1.30000,2.65912) (-1.26667,2.67506) (-1.23333,2.69082) (-1.20000,2.70640) (-1.16667,2.72178) (-1.13333,2.73693) (-1.10000,2.75185) (-1.06632,2.76667) (-1.02765,2.78333) (-0.98803,2.80000) (-0.95000,2.81558) (-0.91667,2.82889) (-0.88333,2.84186) (-0.85000,2.85447) (-0.81667,2.86670) (-0.76954,2.88333) (-0.73333,2.89557) (-0.70000,2.90639) (-0.66667,2.91678) (-0.61667,2.93154) (-0.58333,2.94081) (-0.54847,2.95000) (-0.50000,2.96193) (-0.46667,2.96952) (-0.41667,2.97999) (-0.38333,2.98634) (-0.33333,2.99492) (-0.29979,3.00000) (-0.25000,3.00660) (-0.20000,3.01203) (-0.15000,3.01626) (-0.11667,3.01842) (-0.06667,3.02065) (-0.01667,3.02166) (0.03333,3.02146) (0.08333,3.02004) (0.13333,3.01740) (0.16667,3.01499) (0.21667,3.01036) (0.26667,3.00453) (0.30000,2.99997) (0.35000,2.99219) (0.39943,2.98333) (0.43333,2.97663) (0.47940,2.96667) (0.51667,2.95795) (0.55000,2.94961) (0.60000,2.93623) (0.63333,2.92674) (0.66704,2.91667) (0.71667,2.90103) (0.75000,2.89000) (0.78333,2.87855) (0.81676,2.86667) (0.86192,2.85000) (0.90000,2.83542) (0.93333,2.82228) (0.96667,2.80880) (1.00000,2.79500) (1.03333,2.78090) (1.06667,2.76651) (1.10418,2.75000) (1.14132,2.73333) (1.17782,2.71667) (1.21377,2.70000) (1.24923,2.68333) (1.28333,2.66712) (1.31667,2.65110) (1.35000,2.63494) (1.38333,2.61865) (1.41667,2.60225) (1.45000,2.58575) (1.48333,2.56916) (1.51667,2.55251) (1.55000,2.53580) (1.58333,2.51906) (1.61667,2.50228) (1.65000,2.48549) (1.68333,2.46870) (1.71667,2.45191) (1.75000,2.43514) (1.78333,2.41839) (1.81667,2.40169) (1.85000,2.38503) (1.88333,2.36842) (1.91667,2.35188) (1.95000,2.33540) (1.98333,2.31900) (2.01667,2.30268) (2.05000,2.28645) (2.08333,2.27031) (2.11667,2.25427) (2.15000,2.23834) (2.18333,2.22251) (2.21667,2.20679) (2.25000,2.19118) (2.28333,2.17569) (2.31667,2.16032) (2.35000,2.14508) (2.38333,2.12996) (2.41667,2.11496) (2.45022,2.10000) (2.48795,2.08333) (2.52606,2.06667) (2.56456,2.05000) (2.60000,2.03481) (2.63333,2.02066) (2.66667,2.00665) (2.70000,1.99278) (2.73333,1.97904) (2.76667,1.96544) (2.80490,1.95000) (2.84666,1.93333) (2.88333,1.91886) (2.91667,1.90585) (2.95000,1.89298) (2.98333,1.88024) (3.01922,1.86667) (3.06382,1.85000) (3.10000,1.83665) (3.13333,1.82448) (3.16667,1.81244) (3.20147,1.80000) (3.24869,1.78333) (3.28333,1.77125) (3.31667,1.75976) (3.35000,1.74839) (3.39465,1.73333) (3.43333,1.72045) (3.46667,1.70948) (3.50000,1.69861) (3.54748,1.68333) (3.58333,1.67194) (3.61667,1.66145) (3.65348,1.65000) (3.70000,1.63570) (3.73333,1.62558) (3.76667,1.61556) (3.81667,1.60071) (3.85000,1.59093) (3.88333,1.58124) (3.93333,1.56689) (3.96667,1.55744) (4.00000,1.54807) (4.05000,1.53419) (4.08333,1.52505) (4.11667,1.51599) (4.16667,1.50255) (4.20000,1.49370) (4.23944,1.48333) (4.28333,1.47192) (4.31667,1.46334) (4.36667,1.45062) (4.40000,1.44223) (4.43569,1.43333) (4.48333,1.42158) (4.51667,1.41344) (4.56667,1.40137) (4.60000,1.39340) (4.64257,1.38333) (4.68333,1.37378) (4.71667,1.36605) (4.76667,1.35456) (4.80000,1.34699) (4.85000,1.33573) (4.88333,1.32830) (4.93333,1.31727) (4.96667,1.30999) (5.00000,1.30276)};
\draw[main,line width=.35pt] plot coordinates {(-5.00000,1.61386) (-4.96667,1.62112) (-4.91667,1.63210) (-4.88333,1.63949) (-4.83640,1.65000) (-4.80000,1.65822) (-4.76297,1.66667) (-4.71667,1.67732) (-4.68333,1.68506) (-4.63333,1.69679) (-4.60000,1.70469) (-4.55000,1.71666) (-4.51667,1.72472) (-4.48134,1.73333) (-4.43333,1.74516) (-4.40000,1.75345) (-4.35000,1.76602) (-4.31667,1.77448) (-4.28214,1.78333) (-4.23333,1.79597) (-4.20000,1.80469) (-4.15468,1.81667) (-4.11667,1.82682) (-4.08333,1.83580) (-4.03333,1.84943) (-4.00000,1.85861) (-3.96667,1.86787) (-3.91667,1.88192) (-3.88333,1.89139) (-3.85000,1.90094) (-3.80000,1.91543) (-3.76667,1.92520) (-3.73333,1.93505) (-3.68334,1.95000) (-3.65000,1.96008) (-3.61667,1.97025) (-3.57422,1.98333) (-3.53333,1.99607) (-3.50000,2.00656) (-3.46667,2.01715) (-3.41667,2.03321) (-3.38333,2.04403) (-3.35000,2.05496) (-3.31462,2.06667) (-3.26667,2.08271) (-3.23333,2.09397) (-3.20000,2.10535) (-3.16667,2.11682) (-3.11920,2.13333) (-3.08333,2.14594) (-3.05000,2.15777) (-3.01667,2.16970) (-2.97893,2.18333) (-2.93333,2.19997) (-2.90000,2.21226) (-2.86667,2.22465) (-2.83333,2.23714) (-2.79930,2.25000) (-2.75561,2.26667) (-2.71667,2.28165) (-2.68333,2.29459) (-2.65000,2.30763) (-2.61667,2.32077) (-2.58333,2.33400) (-2.54336,2.35000) (-2.50204,2.36667) (-2.46667,2.38104) (-2.43333,2.39468) (-2.40000,2.40840) (-2.36667,2.42221) (-2.33333,2.43610) (-2.30000,2.45006) (-2.26060,2.46667) (-2.22127,2.48333) (-2.18333,2.49949) (-2.15000,2.51376) (-2.11667,2.52808) (-2.08333,2.54245) (-2.05000,2.55687) (-2.01667,2.57133) (-1.98333,2.58583) (-1.95000,2.60036) (-1.91265,2.61667) (-1.87452,2.63333) (-1.83641,2.65000) (-1.80000,2.66592) (-1.76667,2.68048) (-1.73333,2.69503) (-1.70000,2.70956) (-1.66667,2.72405) (-1.63333,2.73850) (-1.60000,2.75290) (-1.56667,2.76724) (-1.52904,2.78333) (-1.48980,2.80000) (-1.45023,2.81667) (-1.41667,2.83068) (-1.38333,2.84448) (-1.35000,2.85815) (-1.31667,2.87168) (-1.28333,2.88505) (-1.24556,2.90000) (-1.20273,2.91667) (-1.16667,2.93045) (-1.13333,2.94297) (-1.10000,2.95527) (-1.06667,2.96734) (-1.02142,2.98333) (-0.98333,2.99643) (-0.95000,3.00759) (-0.91667,3.01848) (-0.86962,3.03333) (-0.83333,3.04438) (-0.80000,3.05419) (-0.75589,3.06667) (-0.71667,3.07727) (-0.68333,3.08590) (-0.63333,3.09816) (-0.60000,3.10588) (-0.55032,3.11667) (-0.51667,3.12350) (-0.46667,3.13289) (-0.43333,3.13867) (-0.38333,3.14655) (-0.35000,3.15128) (-0.30000,3.15762) (-0.25000,3.16299) (-0.20881,3.16667) (-0.16667,3.16979) (-0.11667,3.17258) (-0.06667,3.17439) (-0.01667,3.17521) (0.03333,3.17505) (0.08333,3.17390) (0.13333,3.17176) (0.18333,3.16865) (0.21667,3.16603) (0.26667,3.16130) (0.31667,3.15561) (0.35922,3.15000) (0.40000,3.14403) (0.45000,3.13583) (0.48333,3.12986) (0.53333,3.12017) (0.56667,3.11321) (0.61667,3.10206) (0.65000,3.09417) (0.69335,3.08333) (0.73333,3.07283) (0.76667,3.06368) (0.81437,3.05000) (0.85000,3.03936) (0.88333,3.02907) (0.92226,3.01667) (0.96667,3.00204) (1.00000,2.99074) (1.03333,2.97917) (1.06855,2.96667) (1.11436,2.95000) (1.15000,2.93673) (1.18333,2.92411) (1.21667,2.91127) (1.25000,2.89825) (1.28765,2.88333) (1.32908,2.86667) (1.36667,2.85134) (1.40000,2.83760) (1.43333,2.82373) (1.46667,2.80975) (1.50000,2.79567) (1.53333,2.78149) (1.56800,2.76667) (1.60674,2.75000) (1.64529,2.73333) (1.68333,2.71682) (1.71667,2.70230) (1.75000,2.68776) (1.78333,2.67320) (1.81667,2.65862) (1.85000,2.64405) (1.88333,2.62947) (1.91667,2.61491) (1.95083,2.60000) (1.98909,2.58333) (2.02743,2.56667) (2.06588,2.55000) (2.10000,2.53526) (2.13333,2.52091) (2.16667,2.50661) (2.20000,2.49238) (2.23333,2.47820) (2.26667,2.46410) (2.30015,2.45000) (2.33997,2.43333) (2.38004,2.41667) (2.41667,2.40153) (2.45000,2.38785) (2.48333,2.37425) (2.51667,2.36074) (2.55000,2.34733) (2.58502,2.33333) (2.62707,2.31667) (2.66667,2.30110) (2.70000,2.28811) (2.73333,2.27522) (2.76667,2.26243) (2.80000,2.24974) (2.84348,2.23333) (2.88333,2.21844) (2.91667,2.20610) (2.95000,2.19386) (2.98333,2.18173) (3.02513,2.16667) (3.06667,2.15184) (3.10000,2.14007) (3.13333,2.12839) (3.16711,2.11667) (3.21565,2.10000) (3.25000,2.08833) (3.28333,2.07710) (3.31667,2.06599) (3.36511,2.05000) (3.40000,2.03861) (3.43333,2.02783) (3.46819,2.01667) (3.51667,2.00130) (3.55000,1.99086) (3.58333,1.98051) (3.62838,1.96667) (3.66667,1.95503) (3.70000,1.94500) (3.73915,1.93333) (3.78333,1.92030) (3.81667,1.91058) (3.85329,1.90000) (3.90000,1.88665) (3.93333,1.87722) (3.97101,1.86667) (4.01667,1.85401) (4.05000,1.84487) (4.09248,1.83333) (4.13333,1.82235) (4.16667,1.81348) (4.21667,1.80032) (4.25000,1.79163) (4.28333,1.78302) (4.33333,1.77024) (4.36667,1.76181) (4.41384,1.75000) (4.45000,1.74103) (4.48333,1.73284) (4.53333,1.72068) (4.56667,1.71265) (4.61667,1.70074) (4.65000,1.69287) (4.69077,1.68333) (4.73333,1.67347) (4.76667,1.66582) (4.81667,1.65445) (4.85000,1.64694) (4.90000,1.63579) (4.93333,1.62843) (4.98333,1.61748) (5.00000,1.61386)};
\draw[main,line width=.35pt] plot coordinates {(-5.00000,2.12256) (-4.95000,2.13331) (-4.91667,2.14054) (-4.87342,2.15000) (-4.83333,2.15884) (-4.79811,2.16667) (-4.75000,2.17745) (-4.71667,2.18499) (-4.66667,2.19639) (-4.63333,2.20406) (-4.58333,2.21567) (-4.55000,2.22348) (-4.50826,2.23333) (-4.46667,2.24324) (-4.43333,2.25125) (-4.38333,2.26336) (-4.35000,2.27151) (-4.30208,2.28333) (-4.26667,2.29215) (-4.23333,2.30051) (-4.18333,2.31316) (-4.15000,2.32167) (-4.10473,2.33333) (-4.06667,2.34322) (-4.03333,2.35195) (-3.98333,2.36517) (-3.95000,2.37406) (-3.91550,2.38333) (-3.86667,2.39657) (-3.83333,2.40569) (-3.79355,2.41667) (-3.75000,2.42878) (-3.71667,2.43813) (-3.67471,2.45000) (-3.63333,2.46180) (-3.60000,2.47139) (-3.55882,2.48333) (-3.51667,2.49565) (-3.48333,2.50548) (-3.44566,2.51667) (-3.40000,2.53034) (-3.36667,2.54040) (-3.33333,2.55053) (-3.28333,2.56585) (-3.25000,2.57615) (-3.21667,2.58651) (-3.17362,2.60000) (-3.13333,2.61271) (-3.10000,2.62330) (-3.06667,2.63397) (-3.01690,2.65000) (-2.98333,2.66089) (-2.95000,2.67177) (-2.91477,2.68333) (-2.86667,2.69923) (-2.83333,2.71031) (-2.80000,2.72145) (-2.76466,2.73333) (-2.71667,2.74955) (-2.68333,2.76088) (-2.65000,2.77225) (-2.61667,2.78368) (-2.56929,2.80000) (-2.53333,2.81243) (-2.50000,2.82401) (-2.46667,2.83562) (-2.42553,2.85000) (-2.38333,2.86479) (-2.35000,2.87650) (-2.31667,2.88825) (-2.28333,2.90001) (-2.23622,2.91667) (-2.20000,2.92948) (-2.16667,2.94129) (-2.13333,2.95310) (-2.09506,2.96667) (-2.05000,2.98262) (-2.01667,2.99441) (-1.98333,3.00618) (-1.95000,3.01793) (-1.90621,3.03333) (-1.86667,3.04719) (-1.83333,3.05882) (-1.80000,3.07041) (-1.76268,3.08333) (-1.71667,3.09916) (-1.68333,3.11054) (-1.65000,3.12185) (-1.61595,3.13333) (-1.56667,3.14978) (-1.53333,3.16079) (-1.50000,3.17169) (-1.46406,3.18333) (-1.41667,3.19847) (-1.38333,3.20896) (-1.35000,3.21933) (-1.30420,3.23333) (-1.26667,3.24460) (-1.23333,3.25444) (-1.19111,3.26667) (-1.15000,3.27831) (-1.11667,3.28754) (-1.07053,3.30000) (-1.03333,3.30978) (-1.00000,3.31832) (-0.95000,3.33073) (-0.91667,3.33873) (-0.86798,3.35000) (-0.83333,3.35772) (-0.79154,3.36667) (-0.75000,3.37519) (-0.70825,3.38333) (-0.66667,3.39106) (-0.61667,3.39977) (-0.58333,3.40525) (-0.53333,3.41292) (-0.50000,3.41768) (-0.45000,3.42429) (-0.40000,3.43022) (-0.36667,3.43380) (-0.31667,3.43861) (-0.26667,3.44273) (-0.21667,3.44615) (-0.16667,3.44887) (-0.13333,3.45029) (-0.08333,3.45184) (-0.03333,3.45267) (0.01667,3.45279) (0.06667,3.45220) (0.11667,3.45089) (0.15000,3.44962) (0.20000,3.44714) (0.25000,3.44395) (0.30000,3.44006) (0.35000,3.43548) (0.38333,3.43205) (0.43333,3.42634) (0.48333,3.41996) (0.51667,3.41534) (0.56667,3.40788) (0.61530,3.40000) (0.65000,3.39403) (0.70000,3.38490) (0.73333,3.37849) (0.78333,3.36839) (0.81667,3.36134) (0.86667,3.35030) (0.90000,3.34265) (0.93924,3.33333) (0.98333,3.32251) (1.01667,3.31407) (1.06667,3.30103) (1.10000,3.29208) (1.13333,3.28294) (1.18333,3.26889) (1.21667,3.25930) (1.25000,3.24954) (1.30000,3.23460) (1.33333,3.22445) (1.36667,3.21416) (1.41183,3.20000) (1.45000,3.18785) (1.48333,3.17711) (1.51667,3.16626) (1.56601,3.15000) (1.60000,3.13867) (1.63333,3.12748) (1.66667,3.11621) (1.71421,3.10000) (1.75000,3.08770) (1.78333,3.07619) (1.81667,3.06462) (1.85863,3.05000) (1.90000,3.03551) (1.93333,3.02380) (1.96667,3.01206) (2.00084,3.00000) (2.04798,2.98333) (2.08333,2.97081) (2.11667,2.95901) (2.15000,2.94720) (2.18915,2.93333) (2.23333,2.91769) (2.26667,2.90590) (2.30000,2.89412) (2.33333,2.88237) (2.37800,2.86667) (2.41667,2.85310) (2.45000,2.84143) (2.48333,2.82981) (2.52114,2.81667) (2.56667,2.80090) (2.60000,2.78941) (2.63333,2.77796) (2.66667,2.76656) (2.71535,2.75000) (2.75000,2.73827) (2.78333,2.72705) (2.81667,2.71588) (2.86435,2.70000) (2.90000,2.68820) (2.93333,2.67723) (2.96667,2.66632) (3.01667,2.65007) (3.05000,2.63932) (3.08333,2.62863) (3.12088,2.61667) (3.16667,2.60218) (3.20000,2.59172) (3.23333,2.58132) (3.28069,2.56667) (3.31667,2.55562) (3.35000,2.54545) (3.39006,2.53333) (3.43333,2.52034) (3.46667,2.51042) (3.50190,2.50000) (3.55000,2.48590) (3.58333,2.47621) (3.61667,2.46659) (3.66667,2.45228) (3.70000,2.44283) (3.73374,2.43333) (3.78333,2.41949) (3.81667,2.41028) (3.85412,2.40000) (3.90000,2.38752) (3.93333,2.37853) (3.97772,2.36667) (4.01667,2.35634) (4.05000,2.34758) (4.10000,2.33456) (4.13333,2.32595) (4.16957,2.31667) (4.21667,2.30471) (4.25000,2.29632) (4.30000,2.28385) (4.33333,2.27561) (4.36979,2.26667) (4.41667,2.25527) (4.45000,2.24724) (4.50000,2.23529) (4.53333,2.22740) (4.57907,2.21667) (4.61667,2.20792) (4.65096,2.20000) (4.70000,2.18877) (4.73333,2.18121) (4.78333,2.16996) (4.81667,2.16253) (4.86667,2.15148) (4.90000,2.14418) (4.94991,2.13333) (4.98333,2.12613) (5.00000,2.12256)};
\draw[main,line width=.35pt] plot coordinates {(-5.00000,2.61966) (-4.95000,2.63004) (-4.91667,2.63702) (-4.86667,2.64755) (-4.83333,2.65463) (-4.78333,2.66533) (-4.75000,2.67251) (-4.70017,2.68333) (-4.66667,2.69066) (-4.62430,2.70000) (-4.58333,2.70909) (-4.54944,2.71667) (-4.50000,2.72780) (-4.46667,2.73536) (-4.41667,2.74679) (-4.38333,2.75446) (-4.33333,2.76606) (-4.30000,2.77385) (-4.25971,2.78333) (-4.21667,2.79353) (-4.18333,2.80149) (-4.13333,2.81351) (-4.10000,2.82158) (-4.05180,2.83333) (-4.01667,2.84196) (-3.98333,2.85020) (-3.93333,2.86264) (-3.90000,2.87099) (-3.85112,2.88333) (-3.81667,2.89209) (-3.78333,2.90061) (-3.73333,2.91347) (-3.70000,2.92210) (-3.65693,2.93333) (-3.61667,2.94389) (-3.58333,2.95269) (-3.53333,2.96597) (-3.50000,2.97487) (-3.46667,2.98383) (-3.41667,2.99734) (-3.38333,3.00639) (-3.34574,3.01667) (-3.30000,3.02923) (-3.26667,3.03843) (-3.22500,3.05000) (-3.18333,3.06162) (-3.15000,3.07096) (-3.10605,3.08333) (-3.06667,3.09447) (-3.03333,3.10393) (-2.98867,3.11667) (-2.95000,3.12773) (-2.91667,3.13731) (-2.87263,3.15000) (-2.83333,3.16136) (-2.80000,3.17103) (-2.75768,3.18333) (-2.71667,3.19529) (-2.68333,3.20502) (-2.64357,3.21667) (-2.60000,3.22944) (-2.56667,3.23922) (-2.53000,3.25000) (-2.48333,3.26372) (-2.45000,3.27353) (-2.41667,3.28334) (-2.36667,3.29804) (-2.33333,3.30784) (-2.30000,3.31763) (-2.25000,3.33229) (-2.21667,3.34204) (-2.18333,3.35178) (-2.13333,3.36633) (-2.10000,3.37600) (-2.06667,3.38564) (-2.01679,3.40000) (-1.98333,3.40958) (-1.95000,3.41908) (-1.90000,3.43324) (-1.86667,3.44261) (-1.83333,3.45193) (-1.78333,3.46579) (-1.75000,3.47495) (-1.71667,3.48403) (-1.66667,3.49751) (-1.63333,3.50639) (-1.59438,3.51667) (-1.55000,3.52822) (-1.51667,3.53678) (-1.46667,3.54942) (-1.43333,3.55772) (-1.39686,3.56667) (-1.35000,3.57795) (-1.31667,3.58583) (-1.26667,3.59740) (-1.23333,3.60495) (-1.18333,3.61601) (-1.15000,3.62319) (-1.10171,3.63333) (-1.06667,3.64049) (-1.01864,3.65000) (-0.98333,3.65678) (-0.93333,3.66603) (-0.90000,3.67199) (-0.85000,3.68057) (-0.81667,3.68606) (-0.76667,3.69394) (-0.72610,3.70000) (-0.68333,3.70609) (-0.63333,3.71277) (-0.60000,3.71696) (-0.55000,3.72286) (-0.50000,3.72827) (-0.45000,3.73319) (-0.41667,3.73619) (-0.36667,3.74028) (-0.31667,3.74386) (-0.26667,3.74693) (-0.21667,3.74947) (-0.18333,3.75088) (-0.13333,3.75255) (-0.08333,3.75370) (-0.03333,3.75432) (0.01667,3.75441) (0.06667,3.75397) (0.11667,3.75300) (0.16667,3.75150) (0.20460,3.75000) (0.25000,3.74783) (0.30000,3.74494) (0.35000,3.74153) (0.40000,3.73761) (0.44839,3.73333) (0.48333,3.72996) (0.53333,3.72472) (0.58333,3.71898) (0.61667,3.71489) (0.66667,3.70837) (0.71667,3.70138) (0.75000,3.69647) (0.80000,3.68874) (0.83336,3.68333) (0.88333,3.67489) (0.92981,3.66667) (0.96667,3.65991) (1.01667,3.65038) (1.05000,3.64383) (1.10000,3.63369) (1.13333,3.62673) (1.18029,3.61667) (1.21667,3.60867) (1.25526,3.60000) (1.30000,3.58972) (1.33333,3.58191) (1.38333,3.56995) (1.41667,3.56183) (1.46436,3.55000) (1.50000,3.54102) (1.53333,3.53251) (1.58333,3.51956) (1.61667,3.51080) (1.65736,3.50000) (1.70000,3.48854) (1.73333,3.47950) (1.78016,3.46667) (1.81667,3.45657) (1.85000,3.44728) (1.89967,3.43333) (1.93333,3.42381) (1.96667,3.41434) (2.01667,3.40003) (2.05000,3.39045) (2.08333,3.38083) (2.13219,3.36667) (2.16667,3.35663) (2.20000,3.34691) (2.24644,3.33333) (2.28333,3.32252) (2.31667,3.31274) (2.36003,3.30000) (2.40000,3.28824) (2.43333,3.27843) (2.47334,3.26667) (2.51667,3.25392) (2.55000,3.24412) (2.58673,3.23333) (2.63333,3.21966) (2.66667,3.20990) (2.70053,3.20000) (2.75000,3.18557) (2.78333,3.17587) (2.81667,3.16619) (2.86667,3.15172) (2.90000,3.14210) (2.93333,3.13252) (2.98333,3.11819) (3.01667,3.10867) (3.05000,3.09920) (3.10000,3.08504) (3.13333,3.07564) (3.16667,3.06628) (3.21667,3.05232) (3.25000,3.04305) (3.28513,3.03333) (3.33333,3.02006) (3.36667,3.01094) (3.40686,3.00000) (3.45000,2.98832) (3.48333,2.97935) (3.53072,2.96667) (3.56667,2.95710) (3.60000,2.94829) (3.65000,2.93514) (3.68333,2.92644) (3.72098,2.91667) (3.76667,2.90488) (3.80000,2.89634) (3.85000,2.88362) (3.88333,2.87519) (3.91725,2.86667) (3.96667,2.85433) (4.00000,2.84607) (4.05000,2.83377) (4.08333,2.82563) (4.12027,2.81667) (4.16667,2.80548) (4.20000,2.79750) (4.25000,2.78563) (4.28333,2.77776) (4.33074,2.76667) (4.36667,2.75832) (4.40269,2.75000) (4.45000,2.73916) (4.48333,2.73157) (4.53333,2.72028) (4.56667,2.71281) (4.61667,2.70169) (4.65000,2.69433) (4.70000,2.68337) (4.73333,2.67612) (4.77711,2.66667) (4.81667,2.65818) (4.85513,2.65000) (4.90000,2.64052) (4.93425,2.63333) (4.98333,2.62311) (5.00000,2.61966)};
\draw[main,line width=.35pt] plot coordinates {(-5.00000,3.10644) (-4.95000,3.11629) (-4.91667,3.12290) (-4.86667,3.13288) (-4.83333,3.13957) (-4.78333,3.14967) (-4.75000,3.15645) (-4.70007,3.16667) (-4.66667,3.17354) (-4.61940,3.18333) (-4.58333,3.19085) (-4.53967,3.20000) (-4.50000,3.20836) (-4.46085,3.21667) (-4.41667,3.22609) (-4.38291,3.23333) (-4.33333,3.24403) (-4.30000,3.25127) (-4.25000,3.26218) (-4.21667,3.26950) (-4.16667,3.28054) (-4.13333,3.28794) (-4.08333,3.29911) (-4.05000,3.30659) (-4.00539,3.31667) (-3.96667,3.32545) (-3.93208,3.33333) (-3.88333,3.34450) (-3.85000,3.35218) (-3.80000,3.36375) (-3.76667,3.37150) (-3.71667,3.38318) (-3.68333,3.39101) (-3.64519,3.40000) (-3.60000,3.41069) (-3.56667,3.41862) (-3.51667,3.43055) (-3.48333,3.43854) (-3.43568,3.45000) (-3.40000,3.45861) (-3.36667,3.46668) (-3.31667,3.47883) (-3.28333,3.48696) (-3.23333,3.49918) (-3.20000,3.50735) (-3.16212,3.51667) (-3.11667,3.52786) (-3.08333,3.53609) (-3.03333,3.54846) (-3.00000,3.55672) (-2.95996,3.56667) (-2.91667,3.57742) (-2.88333,3.58572) (-2.83333,3.59817) (-2.80000,3.60647) (-2.75909,3.61667) (-2.71667,3.62724) (-2.68333,3.63554) (-2.63333,3.64799) (-2.60000,3.65628) (-2.55819,3.66667) (-2.51667,3.67696) (-2.48333,3.68521) (-2.43333,3.69756) (-2.40000,3.70577) (-2.35562,3.71667) (-2.31667,3.72620) (-2.28333,3.73432) (-2.23333,3.74646) (-2.20000,3.75450) (-2.15000,3.76651) (-2.11667,3.77447) (-2.07932,3.78333) (-2.03333,3.79417) (-2.00000,3.80197) (-1.95000,3.81358) (-1.91667,3.82125) (-1.86667,3.83264) (-1.83333,3.84016) (-1.78926,3.85000) (-1.75000,3.85867) (-1.71334,3.86667) (-1.66667,3.87672) (-1.63333,3.88381) (-1.58333,3.89428) (-1.55000,3.90116) (-1.50000,3.91130) (-1.46667,3.91794) (-1.41667,3.92772) (-1.38333,3.93412) (-1.33333,3.94351) (-1.29800,3.95000) (-1.25000,3.95862) (-1.20385,3.96667) (-1.16667,3.97299) (-1.11667,3.98125) (-1.08333,3.98659) (-1.03333,3.99437) (-0.99578,4.00000) (-0.95000,4.00664) (-0.90000,4.01359) (-0.86667,4.01804) (-0.81667,4.02444) (-0.76667,4.03049) (-0.73333,4.03434) (-0.68333,4.03982) (-0.63333,4.04495) (-0.58333,4.04970) (-0.55000,4.05267) (-0.50000,4.05681) (-0.45000,4.06057) (-0.40000,4.06395) (-0.35477,4.06667) (-0.31667,4.06872) (-0.26667,4.07106) (-0.21667,4.07300) (-0.16667,4.07454) (-0.11667,4.07568) (-0.06667,4.07642) (-0.01667,4.07676) (0.03333,4.07669) (0.08333,4.07622) (0.13333,4.07535) (0.18333,4.07407) (0.23333,4.07240) (0.28333,4.07032) (0.33333,4.06785) (0.36667,4.06599) (0.41667,4.06287) (0.46667,4.05936) (0.51667,4.05547) (0.56667,4.05121) (0.60000,4.04816) (0.65000,4.04328) (0.70000,4.03804) (0.74215,4.03333) (0.78333,4.02851) (0.83333,4.02234) (0.87701,4.01667) (0.91667,4.01131) (0.96667,4.00426) (1.00000,3.99938) (1.05000,3.99181) (1.10000,3.98394) (1.13333,3.97853) (1.18333,3.97018) (1.21667,3.96446) (1.26667,3.95565) (1.30000,3.94964) (1.35000,3.94041) (1.38744,3.93333) (1.43333,3.92449) (1.47309,3.91667) (1.51667,3.90794) (1.55563,3.90000) (1.60000,3.89081) (1.63558,3.88333) (1.68333,3.87315) (1.71667,3.86595) (1.76667,3.85500) (1.80000,3.84761) (1.85000,3.83641) (1.88333,3.82886) (1.93333,3.81742) (1.96667,3.80972) (2.00846,3.80000) (2.05000,3.79025) (2.08333,3.78238) (2.13333,3.77049) (2.16667,3.76252) (2.21667,3.75048) (2.25000,3.74242) (2.28740,3.73333) (2.33333,3.72212) (2.36667,3.71396) (2.41667,3.70167) (2.45000,3.69345) (2.49094,3.68333) (2.53333,3.67283) (2.56667,3.66456) (2.61667,3.65213) (2.65000,3.64384) (2.69220,3.63333) (2.73333,3.62308) (2.76667,3.61478) (2.81667,3.60232) (2.85000,3.59401) (2.89292,3.58333) (2.93333,3.57328) (2.96667,3.56500) (3.01667,3.55259) (3.05000,3.54434) (3.09451,3.53333) (3.13333,3.52375) (3.16667,3.51555) (3.21667,3.50327) (3.25000,3.49510) (3.29820,3.48333) (3.33333,3.47478) (3.36674,3.46667) (3.41667,3.45458) (3.45000,3.44655) (3.50000,3.43454) (3.53333,3.42656) (3.57487,3.41667) (3.61667,3.40674) (3.65000,3.39886) (3.70000,3.38709) (3.73333,3.37928) (3.78333,3.36762) (3.81667,3.35988) (3.85945,3.35000) (3.90000,3.34067) (3.93333,3.33305) (3.98333,3.32166) (4.01667,3.31411) (4.06667,3.30285) (4.10000,3.29538) (4.15000,3.28424) (4.18333,3.27685) (4.22956,3.26667) (4.26667,3.25853) (4.30583,3.25000) (4.35000,3.24042) (4.38333,3.23324) (4.43333,3.22253) (4.46667,3.21543) (4.51667,3.20484) (4.55000,3.19783) (4.60000,3.18737) (4.63333,3.18044) (4.68333,3.17011) (4.71667,3.16326) (4.76667,3.15306) (4.80000,3.14630) (4.85000,3.13622) (4.88333,3.12954) (4.93333,3.11959) (4.96667,3.11300) (5.00000,3.10644)};
\draw[main,line width=.35pt] plot coordinates {(-5.00000,3.58446) (-4.95000,3.59368) (-4.91596,3.60000) (-4.86667,3.60919) (-4.82679,3.61667) (-4.78333,3.62485) (-4.73856,3.63333) (-4.70000,3.64067) (-4.65124,3.65000) (-4.61667,3.65664) (-4.56667,3.66630) (-4.53333,3.67277) (-4.48333,3.68252) (-4.45000,3.68905) (-4.40000,3.69889) (-4.36667,3.70547) (-4.31667,3.71540) (-4.28333,3.72204) (-4.23333,3.73206) (-4.20000,3.73876) (-4.15000,3.74885) (-4.11667,3.75560) (-4.06667,3.76578) (-4.03333,3.77258) (-3.98333,3.78283) (-3.95000,3.78969) (-3.90006,3.80000) (-3.86667,3.80692) (-3.81974,3.81667) (-3.78333,3.82425) (-3.73990,3.83333) (-3.70000,3.84170) (-3.66050,3.85000) (-3.61667,3.85923) (-3.58148,3.86667) (-3.53333,3.87686) (-3.50000,3.88393) (-3.45000,3.89456) (-3.41667,3.90166) (-3.36667,3.91232) (-3.33333,3.91944) (-3.28333,3.93013) (-3.25000,3.93727) (-3.20000,3.94798) (-3.16667,3.95513) (-3.11667,3.96585) (-3.08333,3.97300) (-3.03517,3.98333) (-3.00000,3.99087) (-2.95743,4.00000) (-2.91667,4.00873) (-2.87956,4.01667) (-2.83333,4.02654) (-2.80000,4.03365) (-2.75000,4.04430) (-2.71667,4.05138) (-2.66667,4.06197) (-2.63333,4.06901) (-2.58333,4.07954) (-2.55000,4.08654) (-2.50000,4.09699) (-2.46667,4.10392) (-2.41667,4.11428) (-2.38333,4.12114) (-2.33333,4.13139) (-2.30000,4.13818) (-2.25000,4.14829) (-2.21667,4.15499) (-2.16667,4.16496) (-2.13333,4.17156) (-2.08333,4.18137) (-2.05000,4.18785) (-2.00000,4.19748) (-1.96667,4.20383) (-1.91667,4.21327) (-1.88333,4.21948) (-1.83333,4.22869) (-1.80000,4.23476) (-1.75000,4.24373) (-1.71458,4.25000) (-1.66667,4.25835) (-1.61809,4.26667) (-1.58333,4.27252) (-1.53333,4.28078) (-1.50000,4.28620) (-1.45000,4.29416) (-1.41248,4.30000) (-1.36667,4.30699) (-1.31667,4.31442) (-1.28333,4.31925) (-1.23333,4.32632) (-1.18333,4.33317) (-1.15000,4.33761) (-1.10000,4.34407) (-1.05240,4.35000) (-1.01667,4.35431) (-0.96667,4.36012) (-0.91667,4.36567) (-0.88333,4.36923) (-0.83333,4.37435) (-0.78333,4.37920) (-0.73820,4.38333) (-0.70000,4.38667) (-0.65000,4.39077) (-0.60000,4.39459) (-0.55000,4.39812) (-0.51667,4.40031) (-0.46667,4.40334) (-0.41667,4.40608) (-0.36667,4.40852) (-0.31667,4.41064) (-0.26667,4.41247) (-0.21667,4.41398) (-0.16667,4.41518) (-0.11667,4.41606) (-0.06667,4.41664) (-0.03333,4.41685) (0.01667,4.41690) (0.06297,4.41667) (0.10000,4.41629) (0.15000,4.41551) (0.20000,4.41441) (0.25000,4.41300) (0.30000,4.41129) (0.35000,4.40926) (0.40000,4.40693) (0.45000,4.40429) (0.50000,4.40135) (0.53333,4.39923) (0.58333,4.39580) (0.63333,4.39208) (0.68333,4.38807) (0.73333,4.38377) (0.76667,4.38075) (0.81667,4.37600) (0.86667,4.37097) (0.90744,4.36667) (0.95000,4.36200) (1.00000,4.35627) (1.05000,4.35029) (1.08333,4.34617) (1.13333,4.33979) (1.18211,4.33333) (1.21667,4.32863) (1.26667,4.32163) (1.30122,4.31667) (1.35000,4.30949) (1.40000,4.30192) (1.43333,4.29677) (1.48333,4.28887) (1.51769,4.28333) (1.56667,4.27529) (1.61667,4.26691) (1.65000,4.26122) (1.70000,4.25256) (1.73333,4.24669) (1.78333,4.23777) (1.81667,4.23173) (1.86667,4.22257) (1.90000,4.21638) (1.95000,4.20699) (1.98681,4.20000) (2.03333,4.19107) (2.07325,4.18333) (2.11667,4.17484) (2.15807,4.16667) (2.20000,4.15832) (2.24152,4.15000) (2.28333,4.14156) (2.32380,4.13333) (2.36667,4.12457) (2.40508,4.11667) (2.45000,4.10738) (2.48553,4.10000) (2.53333,4.09002) (2.56667,4.08304) (2.61667,4.07253) (2.65000,4.06549) (2.70000,4.05491) (2.73333,4.04784) (2.78333,4.03720) (2.81667,4.03010) (2.86667,4.01942) (2.90000,4.01229) (2.95000,4.00159) (2.98333,3.99445) (3.03333,3.98373) (3.06667,3.97658) (3.11287,3.96667) (3.15000,3.95870) (3.19059,3.95000) (3.23333,3.94084) (3.26838,3.93333) (3.31667,3.92300) (3.35000,3.91588) (3.40000,3.90521) (3.43333,3.89811) (3.48333,3.88747) (3.51667,3.88039) (3.56667,3.86980) (3.60000,3.86275) (3.65000,3.85221) (3.68333,3.84520) (3.73333,3.83471) (3.76667,3.82773) (3.81667,3.81731) (3.85000,3.81037) (3.90000,3.80001) (3.93333,3.79313) (3.98090,3.78333) (4.01667,3.77599) (4.06231,3.76667) (4.10000,3.75899) (4.14432,3.75000) (4.18333,3.74211) (4.22698,3.73333) (4.26667,3.72538) (4.31031,3.71667) (4.35000,3.70878) (4.39437,3.70000) (4.43333,3.69232) (4.47918,3.68333) (4.51667,3.67601) (4.56480,3.66667) (4.60000,3.65986) (4.65000,3.65024) (4.68333,3.64385) (4.73333,3.63433) (4.76667,3.62800) (4.81667,3.61857) (4.85000,3.61231) (4.90000,3.60297) (4.93333,3.59677) (4.98333,3.58753) (5.00000,3.58446)};
\draw[main,line width=.35pt] plot coordinates {(-5.00000,4.05534) (-4.95000,4.06390) (-4.91667,4.06962) (-4.86667,4.07825) (-4.83333,4.08402) (-4.78333,4.09272) (-4.74162,4.10000) (-4.70000,4.10729) (-4.65000,4.11609) (-4.61667,4.12197) (-4.56667,4.13083) (-4.53333,4.13675) (-4.48333,4.14567) (-4.45000,4.15163) (-4.40000,4.16061) (-4.36635,4.16667) (-4.31667,4.17564) (-4.27419,4.18333) (-4.23333,4.19075) (-4.18333,4.19987) (-4.15000,4.20595) (-4.10000,4.21511) (-4.06667,4.22123) (-4.01667,4.23043) (-3.98333,4.23657) (-3.93333,4.24581) (-3.90000,4.25198) (-3.85000,4.26125) (-3.81667,4.26744) (-3.76667,4.27674) (-3.73130,4.28333) (-3.68333,4.29228) (-3.64198,4.30000) (-3.60000,4.30784) (-3.55282,4.31667) (-3.51667,4.32343) (-3.46667,4.33279) (-3.43333,4.33903) (-3.38333,4.34839) (-3.35000,4.35463) (-3.30000,4.36398) (-3.26667,4.37022) (-3.21667,4.37956) (-3.18333,4.38578) (-3.13333,4.39510) (-3.10000,4.40130) (-3.05000,4.41059) (-3.01667,4.41677) (-2.96667,4.42602) (-2.92704,4.43333) (-2.88333,4.44137) (-2.83628,4.45000) (-2.80000,4.45663) (-2.75000,4.46573) (-2.71667,4.47177) (-2.66667,4.48080) (-2.63333,4.48679) (-2.58333,4.49573) (-2.55000,4.50165) (-2.50000,4.51049) (-2.46484,4.51667) (-2.41667,4.52507) (-2.36896,4.53333) (-2.33333,4.53946) (-2.28333,4.54798) (-2.25000,4.55362) (-2.20000,4.56200) (-2.16667,4.56753) (-2.11667,4.57575) (-2.07002,4.58333) (-2.03333,4.58923) (-1.98333,4.59717) (-1.95000,4.60241) (-1.90000,4.61016) (-1.85737,4.61667) (-1.81667,4.62280) (-1.76667,4.63021) (-1.73333,4.63508) (-1.68333,4.64226) (-1.63333,4.64930) (-1.60000,4.65391) (-1.55000,4.66070) (-1.50504,4.66667) (-1.46667,4.67166) (-1.41667,4.67802) (-1.37378,4.68333) (-1.33333,4.68823) (-1.28333,4.69413) (-1.23333,4.69984) (-1.20000,4.70354) (-1.15000,4.70894) (-1.10000,4.71414) (-1.06667,4.71749) (-1.01667,4.72236) (-0.96667,4.72702) (-0.91667,4.73147) (-0.88333,4.73432) (-0.83333,4.73842) (-0.78333,4.74229) (-0.73333,4.74594) (-0.68333,4.74937) (-0.65000,4.75153) (-0.60000,4.75457) (-0.55000,4.75738) (-0.50000,4.75996) (-0.45000,4.76229) (-0.40000,4.76439) (-0.35000,4.76624) (-0.31667,4.76734) (-0.26667,4.76879) (-0.21667,4.76999) (-0.16667,4.77095) (-0.11667,4.77165) (-0.06667,4.77211) (-0.01667,4.77232) (0.03333,4.77227) (0.08333,4.77198) (0.13333,4.77144) (0.18333,4.77065) (0.23333,4.76962) (0.28333,4.76833) (0.33333,4.76680) (0.36667,4.76565) (0.41667,4.76372) (0.46667,4.76154) (0.51667,4.75913) (0.56667,4.75647) (0.61667,4.75358) (0.66667,4.75046) (0.70000,4.74825) (0.75000,4.74475) (0.80000,4.74103) (0.85000,4.73708) (0.89499,4.73333) (0.93333,4.73001) (0.98333,4.72549) (1.03333,4.72076) (1.07491,4.71667) (1.11667,4.71242) (1.16667,4.70716) (1.21667,4.70170) (1.25000,4.69795) (1.30000,4.69218) (1.35000,4.68623) (1.38333,4.68216) (1.43333,4.67592) (1.48333,4.66950) (1.51667,4.66514) (1.56667,4.65845) (1.61667,4.65161) (1.65000,4.64697) (1.70000,4.63988) (1.74532,4.63333) (1.78333,4.62776) (1.83333,4.62030) (1.86667,4.61526) (1.91667,4.60759) (1.96537,4.60000) (2.00000,4.59454) (2.05000,4.58656) (2.08333,4.58118) (2.13333,4.57302) (2.17189,4.56667) (2.21667,4.55921) (2.26667,4.55080) (2.30000,4.54515) (2.35000,4.53660) (2.38333,4.53085) (2.43333,4.52217) (2.46667,4.51635) (2.51667,4.50755) (2.55930,4.50000) (2.60000,4.49275) (2.65000,4.48380) (2.68333,4.47780) (2.73333,4.46875) (2.76667,4.46270) (2.81667,4.45359) (2.85000,4.44749) (2.90000,4.43831) (2.93333,4.43217) (2.98333,4.42294) (3.01724,4.41667) (3.06667,4.40750) (3.10700,4.40000) (3.15000,4.39199) (3.19645,4.38333) (3.23333,4.37645) (3.28333,4.36710) (3.31667,4.36087) (3.36667,4.35151) (3.40000,4.34527) (3.45000,4.33591) (3.48333,4.32967) (3.53333,4.32031) (3.56667,4.31408) (3.61667,4.30473) (3.65000,4.29850) (3.70000,4.28917) (3.73333,4.28295) (3.78333,4.27364) (3.82086,4.26667) (3.86667,4.25816) (3.91071,4.25000) (3.95000,4.24273) (4.00000,4.23350) (4.03333,4.22736) (4.08333,4.21817) (4.11667,4.21205) (4.16667,4.20291) (4.20000,4.19682) (4.25000,4.18772) (4.28333,4.18167) (4.33333,4.17262) (4.36667,4.16661) (4.41667,4.15761) (4.45913,4.15000) (4.50000,4.14269) (4.55000,4.13379) (4.58333,4.12787) (4.63333,4.11903) (4.66667,4.11315) (4.71667,4.10437) (4.75000,4.09853) (4.80000,4.08981) (4.83732,4.08333) (4.88333,4.07537) (4.93333,4.06676) (4.96667,4.06104) (5.00000,4.05534)};
\draw[black,->](-3,-2)--(-2.52244,-1.95167);
\draw[black,->](3,-2)--(3.47756,-2.04833);
\draw[black,->](-3,0)--(-2.54073,0.13955);
\draw[black,->](3,0)--(3.45927,-0.13955);
\draw[black,->](-3,3)--(-2.53964,3.13589);
\draw[black,->](3,3)--(3.46036,2.86411);
\draw[black,->](-2,4)--(-1.53035,4.09913);
\draw[black,->](2,4)--(2.46965,3.90087);
\draw[black,->](0,-1.7)--(0.48000,-1.70000);
\end{scope}
\draw[structurecolor,thick](0,0) circle (1);
\fill (0,1) circle (2pt);\fill (-.968246,-.25) circle (2pt);\fill (.968246,-.25) circle (2pt);
\draw[structurecolor](0,2) circle (2pt) node[above right,fill=white,inner sep=1pt]{$2\ii$};
\node at (0,-3.8){$(b)$ 圆柱外的流};
\end{tikzpicture}
\end{center}

(c) 当 $|z|\to\infty$, 有 $\Phi_a'=1+O(1/|z|)$, $\Phi_b'=1+O(1/|z|)$. 两种流场的速度都趋于 $(1,0)$. 流函数的对数项不必趋零, 远场均匀性应由速度导数判断, 这也说明为何两图的流线在远处逐渐接近水平.
\end{solution}

\begin{exercise}
\textbf{原题 PS9.2。}\label{ass:ps9-2}
考虑 $f(z)=z^5-2z$.
(a) $f$ 把圆 $|z|=3$ 绕原点几次? 即若 $\gamma(\theta)$ 参数化该圆, 曲线 $f\circ\gamma$ 绕原点几次?
(b) $f$ 把圆 $|z|=1$ 绕原点几次?
(c) $f$ 把圆 $|z|=3$ 绕点 $-2$ 几次?
\sourceinfo{题面第1页; 官方解答第3--4页}
\end{exercise}
\begin{solution}
据官方解答翻译并补全. 圆均以逆时针方向参数化. 多项式没有极点, 辐角原理把像曲线的绕数化为相应原像圆内的零点数, 按重数计.
(a) 在 $|z|=3$ 上, $|-2z|=6<243=|z^5|$. Rouché 定理说明 $f$ 与 $z^5$ 的圆内零点数相同, 为 $5$, 故绕数为 $5$.
(b) 在 $|z|=1$ 上, $|z^5|=1<2=|-2z|$, 所以 $f$ 与 $-2z$ 都有一个圆内零点, 绕数为 $1$.
(c) $\Ind(f\circ\gamma,-2)=\Ind((f+2)\circ\gamma,0)$. 在 $|z|=3$ 上, $|-2z+2|\le8<243=|z^5|$, 因而 $f+2$ 有 $5$ 个圆内零点, 所求绕数仍为 $5$.
编者校订: 官方第4页将绕 $-2$ 错转为 $f-2$ 的绕零问题, 正确平移是 $f+2$. 两者在本例碰巧都有五个零点, 但正确的论证必须保留目标点的符号.
\end{solution}

\begin{exercise}
\textbf{原题 PS9.3。}\label{ass:ps9-3}
(a) 证明 $f(z)=z^3+9z+30$ 在 $|z|<2$ 内没有零点.
(b) 证明 $f(z)=z^6+4z^2-1$ 在 $|z|<1$ 内恰有两个零点.
\sourceinfo{题面第1页; 官方解答第4页}
\end{exercise}
\begin{solution}
据官方解答翻译并补全.
(a) 在 $|z|=2$ 上, $|9z+30|\ge30-18=12>8=|z^3|$. 因此 Rouché 定理说明 $f$ 与 $9z+30$ 的圆内零点数相等. 后者唯一零点为 $-10/3$, 在圆外, 所以圆内没有零点.
(b) 在 $|z|=1$ 上, $|4z^2-1|\ge4-1=3>1=|z^6|$. Rouché 定理给出 $f$ 与 $4z^2-1$ 的圆内零点数相同; 后者有两个简单零点 $\pm1/2$, 所以 $f$ 的圆内零点数为 $2$, 按重数计. 实际二者不同: 实函数在 $0$ 值为 $-1$, 在 $1$ 值为 $4$, 因而有一个正实根, 偶性又给出其负根; 这两根已用尽上述计数.
\end{solution}

\begin{exercise}
\textbf{原题 PS9.4。}\label{ass:ps9-4}
设 $f$ 在包含 $|z|\le1$ 的区域解析, 且 $|z|=1$ 上 $|f(z)|<1$. 证明 $f(z)-z$ 在 $|z|<1$ 内恰有一个零点.
\sourceinfo{题面第1页; 官方解答第4页}
\end{exercise}
\begin{solution}
据官方解答翻译并补全. 在单位圆上 $|f(z)|<|-z|=1$. 由 Rouché 定理, $f-z$ 与 $-z$ 在圆内的零点总重数相同, 都为 $1$. 所以 $f-z$ 恰有一个零点, 并且这个零点必为简单零点.
\end{solution}

\begin{exercise}
\textbf{原题 PS9.5。}\label{ass:ps9-5}
考虑负反馈线性系统, 反馈增益为 $k$. 若开环系统函数为 $G(s)$, 闭环系统函数为 $G_{\mathrm{CL}}(s)=G(s)/(1+kG(s))$.
(a) 设 $G(s)=(s+1)/[(s-1)(s-2)]$, $k=4$. 解析判断闭环系统是否稳定.
(b) 画 (a) 的 Nyquist 图, 即 $kG\circ\gamma$, 其中 $\gamma$ 是虚轴. 可以使用任何工具; 原题建议使用\href{https://web.mit.edu/jorloff/www/jmoapplets/nyquist/nyquistCrit.html}{课程 Nyquist 小程序}, 并要求用这一加入反馈增益的版本. 曲线绕 $-1$ 几次? 解释为何与 (a) 一致.
\sourceinfo{题面第1--2页; 官方解答第4--5页}
\end{exercise}
\begin{solution}
据官方解答翻译并补全.
(a) 化简得到
$G_{\mathrm{CL}}(s)=(s+1)/[(s-1)(s-2)+4(s+1)]=(s+1)/(s^2+s+6)$.
极点为 $(-1\pm\ii\sqrt{23})/2$, 实部均为 $-1/2<0$, 与分子无消去, 所以连续时间闭环系统稳定.

(b) 编者补充（AI）: 原题未规定虚轴方向. 为与官方绕数 $2$ 的约定一致, 取 $s=\ii\omega$, $\omega$ 从 $-\infty$ 增至 $+\infty$; 反向参数化则绕数为 $-2$.
直接有
\[
 4G(\ii\omega)=
 \frac{8(1-2\omega^2)}{(\omega^2+1)(\omega^2+4)}
 +\ii\frac{4\omega(5-\omega^2)}{(\omega^2+1)(\omega^2+4)}.
\]
以下用 $\omega=\tan t$, $-\pi/2<t<\pi/2$ 绘制精确参数曲线, 两端都趋于原点; 箭头按 $\omega$ 增加的方向.
\begin{center}
\begin{tikzpicture}[x=1.45cm,y=1.45cm,>=Stealth]
\draw[->](-1.8,0)--(2.4,0) node[right]{$\operatorname{Re}(4G)$};\draw[->](0,-2.15)--(0,2.15) node[above]{$\operatorname{Im}(4G)$};
\draw[main,samples=400,domain=-90:90,smooth] plot
({8*cos(\x)^2*(cos(\x)^2-2*sin(\x)^2)/(sin(\x)^2+4*cos(\x)^2)},
{4*sin(\x)*cos(\x)*(5*cos(\x)^2-sin(\x)^2)/(sin(\x)^2+4*cos(\x)^2)});
\foreach \lo/\hi in {-80/-72,12/20}{\draw[main,->,samples=20,domain=\lo:\hi] plot
({8*cos(\x)^2*(cos(\x)^2-2*sin(\x)^2)/(sin(\x)^2+4*cos(\x)^2)},
{4*sin(\x)*cos(\x)*(5*cos(\x)^2-sin(\x)^2)/(sin(\x)^2+4*cos(\x)^2)});}
\fill (-1,0) circle (1.4pt) node[below]{$-1$};\node[below] at (2,0){$2$};\node[above left] at (-1.33333,0){$-4/3$};
\end{tikzpicture}
\end{center}
绕数可由辐角原理严格确认. 用虚轴向上再沿右半平面大半圆返回的闭轮廓, 方向为顺时针. $H(s)=1+4G(s)$ 在右半平面有两个极点 $1,2$, 没有零点, 因为其零点正是 (a) 的两闭环极点. 故像曲线绕零的绕数为 $P-Z=2$. 大半圆上 $4G(s)\to0$ 一致成立, 因而 $H$ 的该段像缩到 $1$; 虚轴像的绕数已是 $2$. 等价地, $4G\circ\gamma$ 绕 $-1$ 为 $2$ 次, 与 Nyquist 稳定判据及 (a) 的结论相合.
\end{solution}

\begin{exercise}
\textbf{原题 PS9.6。}\label{ass:ps9-6}
令 $f(z)=(z+1/z)/2$. 用 Matlab 或 Mathematica 一类软件绘制下列曲线 $\gamma$ 的像 $f\circ\gamma$. 必须使两个坐标轴比例相同, 让圆显示为圆; Matlab 中可在绘图后使用 \texttt{axis equal}.
(a) 圆 $|z|=3/2$;
(b) 圆 $|z+1/2|=3/2$;
(c) 圆 $|\ee^{\ii\pi/4}z+1/2|=3/2$;
(d) 单位圆 $|z|=1$.
\sourceinfo{题面第2页; 官方解答第5--8页}
\end{exercise}
\begin{solution}
据官方解答翻译并补全. 各图横纵坐标使用同一长度单位. 写 $z=x+\ii y\ne0$, 则
$\operatorname{Re}f(z)=\tfrac12x(1+1/(x^2+y^2))$,
$\operatorname{Im}f(z)=\tfrac12y(1-1/(x^2+y^2))$.
因此每问只需给出原圆参数, 再代入这两式; 下列图均以原生 TikZ 按这些参数生成, 左为原曲线, 右为像.

(a) $z=(3/2)\ee^{\ii\theta}$. 像为 $u=(13/12)\cos\theta$, $v=(5/12)\sin\theta$, 即椭圆 $u^2/(13/12)^2+v^2/(5/12)^2=1$.
\begin{center}
\begin{tikzpicture}[x=1cm,y=1cm,>=Stealth]
\begin{scope}
\draw[->](-1.9,0)--(2,0) node[right]{$x$};\draw[->](0,-1.8)--(0,1.9) node[above]{$y$};\draw[main](0,0) circle (1.5);\node at (0,-2.1){$(a)$ 原圆};
\end{scope}
\draw[->](2.4,.4)--(3.5,.4) node[midway,above]{$f$};
\begin{scope}[xshift=5.6cm]
\draw[->](-1.6,0)--(1.7,0) node[right]{$u$};\draw[->](0,-1.8)--(0,1.9) node[above]{$v$};\draw[main](0,0) ellipse (1.083333 and .416667);\node at (0,-2.1){像椭圆};
\end{scope}
\end{tikzpicture}
\end{center}
(b) $z=(3/2)\ee^{\ii\theta}-1/2$. 在 $\theta=0$ 时 $z=1$, $f(1)=1$, 且 $f'(1)=0$, 所以像在右端形成尖点. 左端 $z=-2$ 的像为 $-5/4$.
\begin{center}
\begin{tikzpicture}[x=1cm,y=1cm,>=Stealth]
\begin{scope}
\draw[->](-2.4,0)--(1.5,0) node[right]{$x$};\draw[->](0,-1.8)--(0,1.9) node[above]{$y$};\draw[main](-.5,0) circle (1.5);\node at (-.5,-2.1){$(b)$ 原圆};
\end{scope}
\draw[->](1.9,.4)--(3,.4) node[midway,above]{$f$};
\begin{scope}[xshift=5.2cm]
\draw[->](-1.6,0)--(1.5,0) node[right]{$u$};\draw[->](0,-1.8)--(0,1.9) node[above]{$v$};
\def\xx{(1.5*cos(\x)-.5)}\def\yy{(1.5*sin(\x))}
\draw[main,samples=300,domain=0:360] plot ({.5*\xx*(1+1/(\xx^2+\yy^2))},{.5*\yy*(1-1/(\xx^2+\yy^2))});
\node at (0,-2.1){像曲线};
\end{scope}
\end{tikzpicture}
\end{center}
(c) 令 $v=\ee^{\ii\pi/4}z+1/2$, 则 $v=(3/2)\ee^{\ii\theta}$, 故 $z=\ee^{-\ii\pi/4}((3/2)\ee^{\ii\theta}-1/2)$. 对应
$x=[(3/2)\cos\theta-1/2+(3/2)\sin\theta]/\sqrt2$,
$y=[(3/2)\sin\theta-(3/2)\cos\theta+1/2]/\sqrt2$.
代入通式得到下图翼型状像曲线. 它可以用于把圆柱周围的流动映到新的边界形状; 本图显示形状, 不单独证明相应外部映射的全局单射性.
\begin{center}
\begin{tikzpicture}[x=1cm,y=1cm,>=Stealth]
\begin{scope}
\draw[->](-2.3,0)--(1.6,0) node[right]{$x$};\draw[->](0,-1.5)--(0,2.3) node[above]{$y$};\draw[main](-.353553,.353553) circle (1.5);\node at (-.35,-1.95){$(c)$ 原圆};
\end{scope}
\draw[->](2,.4)--(3.1,.4) node[midway,above]{$f$};
\begin{scope}[xshift=5.3cm]
\draw[->](-1.6,0)--(1.5,0) node[right]{$u$};\draw[->](0,-1.5)--(0,2.3) node[above]{$v$};
\def\xx{((1.5*cos(\x)-.5+1.5*sin(\x))/sqrt(2))}\def\yy{((1.5*sin(\x)-1.5*cos(\x)+.5)/sqrt(2))}
\draw[main,samples=300,domain=0:360] plot ({.5*\xx*(1+1/(\xx^2+\yy^2))},{.5*\yy*(1-1/(\xx^2+\yy^2))});
\node at (0,-1.95){像曲线};
\end{scope}
\end{tikzpicture}
\end{center}
(d) $z=\ee^{\ii\theta}$ 时 $f(z)=(\ee^{\ii\theta}+\ee^{-\ii\theta})/2=\cos\theta$. 所以像为实轴闭线段 $[-1,1]$, 随原圆一周往返各一次.
\begin{center}
\begin{tikzpicture}[x=1cm,y=1cm,>=Stealth]
\begin{scope}
\draw[->](-1.6,0)--(1.7,0) node[right]{$x$};\draw[->](0,-1.4)--(0,1.5) node[above]{$y$};\draw[main](0,0) circle (1);\node at (0,-1.75){$(d)$ 单位圆};
\end{scope}
\draw[->](2.1,.4)--(3.2,.4) node[midway,above]{$f$};
\begin{scope}[xshift=5.2cm]
\draw[->](-1.6,0)--(1.7,0) node[right]{$u$};\draw[->](0,-1.4)--(0,1.5) node[above]{$v$};\draw[main,thick](-1,0)--(1,0);\fill[main](-1,0) circle (1.5pt) node[below]{$-1$};\fill[main](1,0) circle (1.5pt) node[below]{$1$};\node at (0,-1.75){像线段};
\end{scope}
\end{tikzpicture}
\end{center}
为完整保留原题的软件复现要求, 以下 Matlab/Octave 代码可依次绘制四问; 每个子图均设置 \texttt{axis equal}. 参数 $a,b$ 对应 $z=a\ee^{\ii\theta}+b$.
\begin{verbatim}
th = linspace(0, 2*pi, 1601);
aa = [1.5, 1.5, 1.5*exp(-1i*pi/4), 1];
bb = [0, -0.5, -0.5*exp(-1i*pi/4), 0];
for j = 1:4
    z = aa(j)*exp(1i*th) + bb(j);
    w = 0.5*(z + 1./z);
    figure;
    subplot(1,2,1); plot(real(z), imag(z));
    axis equal; grid on;
    subplot(1,2,2); plot(real(w), imag(w));
    axis equal; grid on;
end
\end{verbatim}
\end{solution}
